\documentclass[11pt]{article}
\usepackage[a4paper,margin=19mm]{geometry}
\usepackage{fontspec}
\setmainfont{lmroman10-regular.otf}
\usepackage{amsmath,amssymb,amsthm,mathtools,amscd,xcolor,graphicx,hyperref}
\hypersetup{hidelinks}
\numberwithin{equation}{section}
\newtheorem{theorem}{Theorem}[section]
\newtheorem{lemma}[theorem]{Lemma}
\newtheorem{proposition}[theorem]{Proposition}
\newtheorem{corollary}[theorem]{Corollary}
\newtheorem{claim}[theorem]{Claim}
\newtheorem{fact}[theorem]{Fact}
\newtheorem{conjecture}[theorem]{Conjecture}
\theoremstyle{definition}
\newtheorem{definition}[theorem]{Definition}
\newtheorem{example}[theorem]{Example}
\newtheorem{remark}[theorem]{Remark}
\newtheorem{question}[theorem]{Question}
\theoremstyle{plain}
\newcommand{\C}{\mathbb C}
\newcommand{\R}{\mathbb R}
\newcommand{\Q}{\mathbb Q}
\newcommand{\Z}{\mathbb Z}
\newcommand{\N}{\mathbb N}
\newcommand{\eee}{\mathrm e}
\newcommand{\cclass}[1]{\textsf{#1}}
\newcommand{\cproblem}[1]{\textsc{#1}}
\newcommand{\ie}{i.\,e.}
\newcommand{\Ie}{I.\,e.}
\newcommand{\eg}{e.\,g.}
\newcommand{\Eg}{E.\,g.}
\newcommand{\phd}{Ph.\,D.}
\newcommand{\msc}{M.\,S.}
\newcommand{\bsc}{B.\,S.}
\newcommand{\expref}[2]{\hyperref[#2]{#1~\ref{#2}}}
\newcommand{\expeqref}[2]{\hyperref[#2]{#1~\eqref{#2}}}
\newcommand{\exprefsilent}[2]{\hyperref[#2]{#1}}
\providecommand{\etal}{et~al.}
\usepackage{eprint}
\usepackage{aumacros}
\title{AC0 hitting-set generators characterized by metacomputational and MKtP circuit lower bounds}
\author{Shuichi Hirahara}
\date{doi:10.4086/toc.2023.v019a004; CC BY 3.0}
\begin{document}
\maketitle
\noindent\textit{Editorial note.} Complete original author statement, local proofs, macros and bibliography retained. Standalone wrapper only; no proof certification or claim repair.\par\medskip
\noindent\textit{Allowed standard prerequisites (editorial).}\begin{itemize}
\item Circuit classes, elementary pseudorandomness and resource-bounded Kolmogorov complexity. All source-local black-box reconstruction conventions, universal generators and quantitative reductions retained.. No collector proof or repaired logical assumption.
\end{itemize}\medskip
\section{Introduction}
\label{section_introduction}


A \emph{promise problem}, introduced by Even, Selman, and Yacobi~\cite{EvenSY84_iandc_journals},
is a pair of disjoint sets $(\YES, \NO)$ that are regarded as the sets of $\Yes$ and $\No$ instances, respectively.
The problem is regarded as a problem whose instances are ``promised'' to come from $\YES \cup \NO$.
Specifically,
an algorithm $A$ is said to \emph{solve} a promise problem $(\YES, \NO)$
if
$A$ accepts any instance $x \in \YES$ and rejects any instance $x \in \NO$;
the behavior of $A$ on any ``unpromised'' instance $x \not\in \YES \cup \NO$ can be arbitrary.
The notion of promise problem is crucial for formalizing several important concepts and theorems in complexity theory.
A canonical example is the unique satisfiability problem $(\mathsf{1SAT}, \UNSAT)$, where $\mathsf{1SAT}$ is the set of formulas that has a unique satisfying assignment, and $\UNSAT$ is the set of unsatisfiable formulas.
%
%
The promise problem is a standard $\pr\UP$-complete problem, and
the celebrated theorem of Valiant and Vazirani~\cite{ValiantV86_tcs_journals}
states that it is in fact $\NP$-hard under randomized reductions.
The reader is referred to the survey of Goldreich~\cite{Goldreich06a_birthday_conf} for more background on promise problems.

Usually, it is required that $\YES$ and $\NO$ are disjoint, \ie, $\YES \cap \NO = \emptyset$.
The reason is that if there exists an instance $x \in \YES \cap \NO$,
then
no algorithm can solve the promise problem $(\YES, \NO)$.
Indeed, if there were an algorithm $A$ that solves $(\YES, \NO)$,
then
$A$ must accept $x$ and simultaneously reject $x$, which is impossible.
For this reason,
every definition of promise problems considered before
is, to the best of our knowledge,
always disjoint.

In this paper, we introduce a set of new promise problems which are \emph{conjectured to be non-disjoint}.
We will demonstrate that these ``non-disjoint'' promise problems are worth investigating,
by showing
that hardness results for our promise problems have important consequences in complexity theory.
The fact that the promise problems are conjectured to be non-disjoint
means
that
solving promise problems are conjectured to be \emph{impossible},
no matter how long an algorithm is allowed to run.
Nevertheless,
under several (implausible) assumptions,
we show that ``non-disjoint'' promise problems can be \emph{efficiently solved};
in other words, the promise problems are not only disjoint but also can be shown to be disjoint in a ``constructive'' way.
%
%
Taking the contrapositive,
we prove that
mild hardness results for computing ``non-disjoint'' promise problems, which is conjectured to be \emph{impossible},
are sufficient to resolve important open questions of complexity theory.


To be more specific,
we consider open questions of whether there exists an explicit hitting set generator.
A \emph{hitting set generator} (HSG) $G = \ESet{G_n \colon \binset^{s(n)} \to \binset^n}_{n \in \Nat}$
secure against a class $\mathfrak{C}$
is a family of functions $G_n$
such that
any algorithm $A \in \mathfrak{C}$
that accepts at least a half of the $n$-bit strings
must accept a string $G_n(z)$ for some seed $z \in \binset^{s(n)}$,
for all large $n \in \Nat$.
The existence of a secure hitting set generator
makes it possible to
derandomize any one-sided-error $\mathfrak{C}$-randomized algorithm,
by simply trying all possible $s(n)$-bit seeds $z$
and using $G_n(z)$ as a source of randomness.
A stronger notion called a \emph{pseudorandom generator} (PRG) enables us to derandomize two-sided-error randomized algorithms.


\subsection{Meta-computational view of PRG constructions}
A standard approach for constructing pseudorandom generators
is to use the \emph{hardness versus randomness} framework
developed in, \eg, \cite{Yao82a_focs_conf, BlumM84_siamcomp_journals, NisanW94_jcss_journals, BabaiFNW93_cc_journals, Impagliazzo95_focs_conf, ImpagliazzoW97_stoc_conf, SudanTV01_jcss_journals, KlivansM02_siamcomp_journals}.
One of the landmark results of Impagliazzo and Wigderson~\cite{ImpagliazzoW97_stoc_conf}
states that if there exists a function in $\E = \DTIME(2^{O(n)})$
that is not computable by a circuit of size $2^{\alpha n}$ for some constant $\alpha > 0$,
then there exists a logarithmic-seed-length pseudorandom generator secure against linear-size circuits (and, in particular, $\P = \BPP$ follows).
%
In general,
such a result is proved by using a \emph{black-box pseudorandom generator construction} $G^{(\blank)}$
that converts any hard function $f \not\in \SIZE(2^{o(n)})$ to a pseudorandom generator $G^f \colon \binset^{O(n)} \to \binset^{2^{\alpha n}}$
secure against circuits of size $2^{\alpha n}$,
where $\alpha > 0$ is some constant.
%
%
%
%
%
%
%
%
%
%
%
%
%
%
%
%
%
%
%
%
%
%
%
%
%
%
%
%
%

The underlying theme of this paper is to view black-box PRG constructions from a \emph{meta-computational} perspective.
Usually, $f$ is regarded as a \emph{fixed} hard function such
as $f \not\in \SIZE(2^{o(n)})$.
Instead, here we regard $f$ as an \emph{input} to some ``non-disjoint'' promise problem, and regard a black-box PRG construction $G^{(\blank)}$ as a
reduction
that proves the security of some (universal) hitting set generator
based on the hardness of the ``non-disjoint'' promise problem.
This new perspective can be applied to arbitrary black-box
PRG constructions,
and it gives rise to a ``non-disjoint'' promise problem associated with the black-box PRG construction.
For example, the pseudorandom generator construction of \cite{ImpagliazzoW97_stoc_conf} induces the $\vsProblem{\E}{\SIZE(2^{o(n)})}$ Problem,
which is the problem of distinguishing whether $f \in \E/O(n)$
or $f \not\in \SIZE(2^{o(n)})$,
given the truth table of a function $f \colon \binset^n \to \binset$.
Here, $\E/O(n)$ denotes the class of functions that can be computed in time $2^{O(n)}$ with $O(n)$ bits of advice (see \cref{def_EvsSIZE,section_MCLPdef} for a precise definition).

%
%
%
%
%
%
%
%
%


There are two types of a pseudorandom generator.
One is a \emph{cryptographic} PRG, which is computable in polynomial time in its seed length.
This notion is useful for building secure cryptographic primitives.
We present
in \cref{section_CryptographicHSG}
``non-disjoint'' promise problems
whose hardness gives rise to a cryptographic hitting set generator.
This provides
a new approach for establishing the equivalence between the worst-case and average-case complexity of $\PH$.
The other is a \emph{complexity-theoretic} PRG, which is allowed to be computed in time exponential in its seed length.
This notion is sufficient for the purpose of derandomizing randomized algorithms.
In \cref{section_ComplexityHSG},
we generalize the hardness versus randomness framework by
using the meta-computational view of black-box PRG constructions,
and establish the equivalence between
circuit lower bounds for ``non-disjoint'' promise problems
and
the existence of hitting set generators.
\cref{section_CryptographicHSG,section_ComplexityHSG} can be read independently.

%
%
%
%
%
%
%
%
%
%

%
%
%
%
%
%
%
%
%
%

%
%

%
%
%
%
%
%
%
%




\subsection{Worst-case versus average-case complexity of PH}
\label{section_CryptographicHSG}

Understanding average-case complexity is a fundamental question
in complexity theory.
Average-case hardness of $\NP$ is a prerequisite for building secure cryptographic primitives,
such as one-way functions and cryptographic pseudorandom generators.
Indeed, it is not hard to see that if there exists a polynomial-time-computable hitting set generator $G$,
then checking whether a given string is in the image of $G$ is
a problem in $\NP$ that is hard on average (in the errorless sense).
In this section, we present a new approach for
proving the average-case hardness of $\PH$,
by implicitly
constructing a cryptographic hitting set generator.


A fundamental open question in the theory of average-case complexity, pioneered by \cite{Levin86_siamcomp_journals}, is
to establish the equivalence between the worst-case and average-case complexity of $\NP$.
\begin{open}
  \label{open_Heuristica}
  Does $\P \neq \NP$ imply $\Dist\NP \not\subset \Avg\P$?
\end{open}
Here, $\Dist\NP \not\subset \Avg\P$ is an average-case analogue of $\NP \neq \P$,
and it means that there exist a language $L \in \NP$ and a polynomial-time samplable distribution $\mathcal{D}$
such that no average-case polynomial-time algorithm (equivalently, errorless heuristic scheme) can solve $L$ on instances randomly drawn from $\mathcal{D}$.
%
%
%
The reader is referred to the survey of Bogdanov and Trevisan~\cite{BogdanovT06_siamcomp_journals} for the formal definitions of $\Dist\NP$ and $\Avg\P$
as well as background on average-case complexity.

For large enough complexity classes such as $\PSPACE$ and $\EXP$,
there is a general technique for converting any worst-case hard function $f$ to some two-sided-error average-case hard function $\Enc(f)$ based on error-correcting codes \cite{SudanTV01_jcss_journals, TrevisanV07_cc_journals}.
Here, the encoded function $\Enc(f)$ is computable in $\PSPACE$ given oracle access to $f$; thus, the worst-case and average-case
%
complexities  %
of such large complexity classes are known to be equivalent.
%
%
Viola \cite{Viola05_cc_journals} showed limitations of such an approach:
%
%
$\Enc(f)$ cannot be computed in
the polynomial-time hierarchy $\PH^f$; thus, the proof technique of using error-correcting codes is not sufficient to resolve \cref{open_Heuristica};
%
%
%
%
%
%
it could not even confirm a positive answer to  %
the following weaker open question:

\begin{open}
  \label{open_PHWorstAndAverage}
  Does $\PH \neq \P$ (or, equivalently, $\P \neq \NP$)
  imply
  $\Dist\PH \not\subset \Avg\P$?
\end{open}
Observe that \cref{open_PHWorstAndAverage} is an easier question than \cref{open_Heuristica}, since $\PH = \P$ is known to be equivalent to $\NP = \P$.

There are significant obstacles %
to  %
resolving  %
%
\cref{open_Heuristica}.
One is the relativization barrier due to Impagliazzo~\cite{Impagliazzo11_coco_conf}.  Another is the limits of black-box reductions
due to Feigenbaum and Fortnow~\cite{FeigenbaumF93_siamcomp_journals} and Bogdanov and Trevisan~\cite{BogdanovT06_siamcomp_journals}.
The limits of black-box reductions are partially extended to a ``barrier'' result
for \cref{open_PHWorstAndAverage}
in \cite{HiraharaW20_approx_conf}.

Recently, a non-black-box worst-case to average-case reduction
that is not subject to the latter barrier
was presented in \cite{Hirahara18_focs_conf}.
The reduction shows that solving
the problem $\Gap\MINKT$ of
approximating polynomial-time-bounded Kolmogorov complexity
in the worst case
can be reduced to solving $\MINKT$ on average.
We briefly review the definition of $\Gap\MINKT$ below.
For an integer $t \in \Nat$ and an oracle $A$,
a \emph{$t$-time-bounded $A$-oracle Kolmogorov complexity} $\KtTIME[A]{t}(x)$ of a finite string $x$
is defined as
the shortest length of a program that prints $x$ in $t$ steps
with oracle access to $A$ (see \cref{section_preliminaries} for a precise definition).
The promise problem $\Gap\MINKT = (\YES, \NO)$
%
asks for approximating $\KtTIME{t}(x)$ within an additive error of
$\widetilde{O}(\sqrt{\KtTIME{t}(x)})$,
and is formally defined as follows:
$\YES$ consists of $(x, \unary{s}, \unary{t})$ such that $\KtTIME{t}(x) \le s$;
and
$\NO$ consists of $(x, \unary{s}, \unary{t})$ such that $\KtTIME{\poly(|x|, t)}(x) > s + \widetilde{O}(\sqrt{s})$.

The result of \cite{Hirahara18_focs_conf}
can be seen as providing an approach for establishing the equivalence between worst-case and average-case complexity of $\NP$;
indeed,
proving $\NP$-hardness of $\Gap\MINKT$ is sufficient for resolving \cref{open_Heuristica} affirmatively.  %
However, the approximation error $\widetilde{O}(\sqrt{s})$ caused by the reduction of \cite{Hirahara18_focs_conf} is not optimal,
which makes the question of proving $\NP$-hardness of $\Gap\MINKT$ potentially harder.
%

\subsubsection{\texorpdfstring{$\KvsK{A}$}{Gap(K(A) vs K)}}
We herein introduce the following ``non-disjoint'' promise problem.
%
%
%
%

\begin{definition}
  For
  an oracle $A$ and an approximation quality $\tau \colon \Nat \times \Nat \to \Nat$,
  the problem $\KvsK[\tau]{A}$ %
  is defined as the following promise problem $(\YES, \NO)$.
  \begin{align*}
    \YES & := \ISet{ (x, \unary{s}, \unary{t}) } { \KtTIME[A]{t}(x) \le s },
    \\
    \NO  & := \ISet{ (x, \unary{s}, \unary{t}) } { \KtTIME{\tau(|x|, t)}(x) > s + \log \tau(|x|, t) }.
  \end{align*}
  By default, we assume that $\tau$ is some polynomial
  and write $\KvsK{A} \in \P$ if
  there exists some polynomial $\tau$ such that
  $\KvsK[\tau]{A} \in \P$.
  %
  %
  %
  %
  %
  %
  %
\end{definition}

In this paper, we prove
%
\begin{theorem}
  \label{thm_HeuristicaIntro}
  Let $A$ be any oracle.
  If $\Dist\NP^A \subset \Avg\P$,
  then
  $\KvsK{A} \in \P$.
  %
  %
  %
\end{theorem}

An immediate corollary of \cref{thm_HeuristicaIntro} is an improvement of the reduction of \cite{Hirahara18_focs_conf}, by setting $A := \emptyset$.
In particular, in order to resolve \cref{open_Heuristica},
it suffices to prove $\NP$-hardness of approximating $\KtTIME{t}(x)$ within an additive error of $\log \tau(|x|, t)$ given $(x, \unary{t})$ as input, for any polynomial $\tau$.
A key insight for reducing the approximation error is
that there are two main sources of the approximation error in the reduction of \cite{Hirahara18_focs_conf}:
One comes from fixing a random coin flip sequence,
which we remove by using the (complexity-theoretic) pseudorandom generator
constructed by
Buhrman, Fortnow, and Pavan~\cite{BuhrmanFP05_mst_journals}
under the assumption that $\Dist\NP \subset \Avg\P$.
The other comes from
the advice complexity of a black-box pseudorandom generator construction,
which we reduce by using
a ``$k$-wise direct product generator''
whose advice complexity is small \cite{Hirahara20_stoc_conf}.
%
%
%

More surprisingly,
the promise problem is conjectured to be \emph{non-disjoint} for $A := \SAT$.
That is, it is conjectured to be \emph{impossible} for any algorithm to solve $\KvsK{\SAT}$ --- no matter how long the algorithm is allowed to run.
Nevertheless,
\cref{thm_HeuristicaIntro}
shows that
under the assumption that $\PH$ is easy on average,
there exists a \emph{polynomial-time algorithm} for solving $\KvsK{\SAT}$.
Taking its contrapositive,
this means that,
in order to resolve $\Dist\PH \not\subset \Avg\P$,
it suffices to prove a super-polynomial time lower bound
for solving $\KvsK{\SAT}$,
whose time complexity is conjectured to be \emph{infinity} (in the sense that there exists \emph{no algorithm} that can compute the promise problem).

We now clarify why $\KvsK{\SAT}$ is conjectured to be non-disjoint.
Under the plausible assumption that $\E^\NP \neq \E$,
it is not hard to see that
there are infinitely many strings $x$ such that $x$ is simultaneously a \Yes and \No instance;
here, the string $x$ is defined as the truth table of the characteristic function of $L \in \E^\NP \setminus \E/O(n)$
(see \cref{prop_NondisjointKSATvsK}).


Another corollary of \cref{thm_HeuristicaIntro} is
that
under the assumption that $\Dist\PH \subset \Avg\P$,
any string $x$ that can be compressed with $\SAT$ oracle in polynomial time
can be also compressed without any oracle.
Formally:
\begin{corollary}
  [see also \cref{cor:oracle elimination}]
  \label{cor_KandKSAT}
  If $\Dist\PH \subset \Avg\P$,
  then
  there exists a polynomial $\tau$ such that
  $$\KtTIME{\tau(|x|, t)}(x) \le \KtTIME[\SAT]{t}(x) + \log \tau(|x|, t)$$
  for any $x \in \binset^*$ and $t \in \Nat$.
\end{corollary}
\begin{proof}
  [Proof Sketch]
  Under the assumption, $\KvsK{\SAT}$
  can be solved by \emph{some algorithm}.
  Thus $\KvsK{\SAT}$ problem must be disjoint, from which the result follows immediately.
\end{proof}

%
%
%
\cref{cor_KandKSAT} provides a new approach for resolving \cref{open_PHWorstAndAverage}.
In order to prove $\Dist\PH \not\subset \Avg\P$ under the assumption that $\P \neq \NP$,
it suffices to find a string $x$ that can be compressed with $\SAT$ oracle
but cannot be compressed without $\SAT$ oracle.
In fact, it suffices to find such a string $x$ under the stronger assumption that $\NP \not\subset \Ppoly$.
%
%
%
%
This is because Pavan, Santhanam, and Vinodchandran~\cite{PavanSV06_fsttcs_conf}
proved
$\NP \not\subset \P/\poly$ if
%
%
the answer to \cref{open_PHWorstAndAverage} is negative.  %

More importantly, \cref{thm_HeuristicaIntro} suggests a more general approach to \cref{open_PHWorstAndAverage}.
%
Note that finding a string $x$ compressible with $\SAT$ oracle but incompressible
without any oracle
corresponds to proving the non-disjointness of $\KvsK{\SAT}$;
this amounts to proving
the time complexity of solving $\KvsK{\SAT}$
is \emph{infinity}.
\cref{thm_HeuristicaIntro}
suggests that
it suffices to prove that
a \emph{polynomial-time} algorithm cannot find a difference between
compressible strings under $\SAT$ oracle
and
incompressible strings without any oracle,
under the worst-case hardness assumption for $\NP$.
In other words, %
it suffices to prove $\NP$-hardness of $\KvsK{\SAT}$.

\begin{corollary}
  [A new approach for \cref{open_PHWorstAndAverage}]
  \label{cor_TowardsPH}
  Suppose that the $\KvsK{\SAT}$ problem is ``$\NP$-hard
  under randomized reductions''%
  \footnote{
    Here we use the weak notion of ``$\NP$-hardness'' in order to strengthen the result.
    \cref{cor_TowardsPH} remains true even if one interprets $\NP$-hardness as a randomized reduction from $\NP$.
  }
  in the sense that
  $$
    \NP \not\subset \BPP
    \quad
    \Longrightarrow
    \quad
    \KvsK{\SAT} \not\in \P.
  $$
  Then,
  \cref{open_PHWorstAndAverage} is true; that is,
  $$
    \Dist\PH \not\subset \Avg\P
    \quad
    \Longleftrightarrow
    \quad
    \PH \neq \P.
  $$
\end{corollary}

In a typical proof of $\NP$-hardness of a disjoint promise problem $\Pi = (\YES, \NO)$,
one needs to carefully design
a reduction $R$ from $\SAT$ to $\Pi$
that ``preserves'' a structure of $\SAT$;
\ie,
any formula $\phi \in \SAT$ is mapped to $R(\phi) \in \YES$
and
any formula $\phi \in \UNSAT$ is mapped to $R(\phi) \in \NO$.
The task of proving $\NP$-hardness of $\KvsK{\SAT}$
is potentially much easier:
in principle, $R(\phi)$ can be a fixed input $x$ that is in the intersection of
\Yes and \No instances of $\KvsK{\SAT}$;
more generally, if a promise problem $\Pi$ is non-disjoint,
then any problem is reducible to $\Pi$ via a reduction that maps every instance
to the intersection of \Yes and \No instances of $\Pi$.
%
%
%
%
%
%
%
%
%
%
%
%
It is worth mentioning that
proving $\NP$-hardness of $\KvsK{\SAT}$
is at least as easy as proving $\NP$-hardness of $\Gap\MINKT$
since
$\Gap\MINKT$ is reducible to $\KvsK{\SAT}$ via an identity map.
%
%
%
%

Subsequent to this work, converse directions of \cref{thm_HeuristicaIntro,cor_TowardsPH} are established in the following sense \cite{Hirahara20_focs_conf}:
$\Dist\PH \subset \Avg\P$ if and only if $\KvsK{A} \in \P$ for every oracle $A \in \PH$.
In particular,%
\footnote{using the fact that $\P = \BPP$ if $\Dist\PH \subset \Avg\P$ \cite{BuhrmanFP05_mst_journals}}
$\Dist\PH \subset \Avg\P \iff \PH = \P$ (\cref{open_PHWorstAndAverage}) if and only if $\KvsK{A} \in \P$ for every $A \in \PH$ implies $\NP \subset \BPP$
(``$\NP$-hardness'' of $\KvsK{A}$).
Therefore, the approach to \cref{open_PHWorstAndAverage} suggested in this work
is one of the most general approaches.

\subsubsection{Non-NP-hardness results do not apply}

A line of work presented evidence that $\NP$-hardness of $\MINKT$ is not likely to be established under
deterministic reductions (\eg, \cite{Ko91_siamcomp_journals, MurrayW17_toc_journals, HitchcockP15_fsttcs_conf, HiraharaW16_coco_conf}).
For example, it is not hard to see that the proof technique of Murray and Williams \cite{MurrayW17_toc_journals} (who proved a similar result for $\MCSP$)
can be extended to the case of $\Gap\MINKT$.
\begin{theorem}
  [Essentially in \cite{MurrayW17_toc_journals};
    %
    %
    see \cite{Hirahara20_stoc_conf}]  %
  %
  If $\Gap\MINKT$ is $\NP$-hard under many-one deterministic reductions,
  then
  $\EXP \neq \ZPP$.
\end{theorem}
This result suggests that establishing $\NP$-hardness of $\Gap\MINKT$ under deterministic reductions is a challenging task.
In contrast,
we observe that a similar ``non-$\NP$-hardness'' result cannot be applied to
the ``non-disjoint'' promise problem.
\begin{proposition}
  \label{prop_NoNonNPhardness}
  Assume that
  $\NP$-hardness of $\KvsK{\SAT}$
  under many-one reductions
  implies $\EXP \neq \ZPP$.
  Then, $\EXP \neq \ZPP$ holds unconditionally.
\end{proposition}
The reason is that $\KvsK{\SAT}$ is well defined only if $\EXP \neq \ZPP$.  More formally:
\begin{proof}
  There are two cases.
  Either $\KvsK{\SAT}$ is disjoint or not disjoint.
  In the former case, by \cref{prop_NondisjointKSATvsK}, we have $\E^\NP = \E$,
  which implies $\EXP^\NP = \EXP$ by a standard padding argument;
  thus, we obtain $\EXP = \EXP^\NP = \ZPEXP \neq \ZPP$,
  where the last inequality follows from a standard translation argument (see, \eg, \cite{BuhrmanFS09_icalp_conf}).
  In the latter case, there exists a string $x$ that is simultaneously a \Yes and \No instance.
  A reduction that always maps to $x$ defines a many-one reduction from any problem to
  $\KvsK{\SAT}$; thus, $\EXP \neq \ZPP$ follows from the assumption.
  %
  %
  %
  %
  %
  %
  %
  %
  %
\end{proof}

In light of \cref{prop_NoNonNPhardness}, we leave as an interesting open question whether
there is any barrier explaining
the difficulty of proving $\NP$-hardness of the non-disjoint promise problem.
We mention that
$\Gap\MINKT^\SAT$, which is equivalent to
$\KvsKA{\SAT}{\SAT}$,
is known to be $\Dist\NP$-hard \cite{Hirahara20_stoc_conf}.
In particular, since $\KvsKA{\SAT}{\SAT}$ is reducible to $\KvsK{\SAT}$ via an identity map,
the latter is also $\Dist\NP$-hard.
Therefore, in order to present a barrier for proving $\NP$-hardness of $\KvsK\SAT$,
one must exploit a property that holds for $\NP$ but does not hold for $\Dist\NP$
(unless the notion of reducibility is too strong).


\subsubsection{\texorpdfstring{$\GapVS{F}{F^{-1}}$}{Gap(F vs F(-1))}: Meta-computational view of HILL's PRG}

We also propose another approach towards
\cref{open_Heuristica},
by introducing a promise problem
which asks for distinguishing whether a given function is
computable by small circuits,
or cannot be inverted by small circuits.
Specifically,
for an approximation quality $\tau$,
we define the promise problem $\GapVS[\tau]{F}{F^{-1}}$
as follows.
Given a size parameter $s \in \Nat$ and an integer $n \in \Nat$ and random access to a function $F \colon \binset^n \to \binset^n$,
the task is to distinguish the following two cases:
\begin{description}
  \item
        [\Yes:]
        $F$ is computable by a circuit of size $s$.
  \item
        [\No:]
        $F$ cannot be inverted on average by any $F$-oracle circuit of size $\tau(n, s)$.
\end{description}

We show that ``$\NP$-hardness'' of $\GapVS[\tau]{F}{F^{-1}}$
for every polynomial $\tau$
is enough for resolving \cref{open_Heuristica}.
More specifically,
we prove
\begin{theorem}
  If $\Dist\NP \subset \Avg\P$,
  then there exist a polynomial $\tau$
  and
  a $\co\RP$-type randomized algorithm that solves
  $\GapVS[\tau]{F}{F^{-1}}$
  on input $(n, s)$
  in time $\poly(n, s)$.
  In particular,
  \cref{open_Heuristica} is true
  if $\GapVS[\tau]{F}{F^{-1}}$ is \textup{``}$\NP$-hard\textup{''} for every polynomial $\tau$
  in the following sense:
  $\NP \subset \BPP$ follows from the assumption that
  $\GapVS[\tau]{F}{F^{-1}}$ admits a $\co\RP$-type algorithm.
\end{theorem}
This is proved by viewing the black-box PRG construction
based on any one-way function,
which is given by
H{\aa}stad, Impagliazzo, Levin, and Luby~\cite{HastadILL99_siamcomp_journals},
from the meta-computational perspective.
Our general proof strategy is given in \cref{subsection:proofsketch}.

It is easy to observe that
$\GapVS{F}{F^{-1}}$
is non-disjoint
under
the existence of a one-way function,
which is one of the most standard cryptographic primitives.
Thus,
it is widely believed that
$\GapVS{F}{F^{-1}}$ is
\emph{impossible} to solve.
Nevertheless, $\NP$-hardness of $\GapVS{F}{F^{-1}}$
is sufficient for resolving \cref{open_Heuristica}.

%
\subsection{Meta-computational circuit lower-bound problems; MCLPs}
\label{section_ComplexityHSG}

We now turn our attention to complexity-theoretic hitting set generators.
A standard approach for constructing complexity-theoretic pseudorandom generators
is to use the hardness versus randomness framework,
which reduces the task of constructing
a pseudorandom generator
to the task of finding an explicit hard function, such as $f \in \E \setminus \SIZE(2^{o(n)})$.

It is, however, a widely accepted fact that
proving a circuit size lower bound for an explicit function is extremely hard.
Here by an \emph{explicit} function, we mean that a function is computable in $\E = \DTIME(2^{O(n)})$.
It is an open question
whether there exists an exponential-time-computable function $f \in \E$
that cannot be computed by any circuit of size $4 n$
%
(see \cite{FindGHK16_focs_conf}).  %
On the other hand, a simple counting argument
shows that
most functions $f \colon \binset^n \to \binset$
cannot be computed by circuits of size $2^{\alpha n}$ for any constant $\alpha < 1$.

Why is it so difficult to prove a circuit lower bound for an explicit function?
We propose to view this question from a \emph{meta-computational} perspective.
The fact that it is difficult for \emph{human beings} to show that an explicit function
cannot be computed by small circuits
suggests
that
it should be also difficult for \emph{algorithms} to analyze a circuit lower bound.
Our results indicate that
if we can make this intuition formal, then we get breakthrough results in complexity theory.

Specifically,
we herein introduce a family of new computational problems,
which we call \emph{Meta-computational Circuit Lower-bound Problems} (MCLPs).
These problems ask for distinguishing the truth table of explicit functions
from hard functions.
For example, we propose the following promise problem:
\begin{description}
  \item
        [The $\vsProblem{\E}{\SIZE(2^{o(n)})}$ Problem] (informal)%
        \footnote{This problem may be called the $\vsProblem{\E/O(n)}{\SIZE(2^{o(n)})}$ Problem; however, for the sake of notational simplicity (and for the reason described in the following paragraph), we omit ``$/O(n)$'' from the name of the problem.}
        \hspace{0em}

        Given the truth table of a function $f \colon \binset^n \to \binset$,
        distinguish whether
        $f \in \E / O(n)$
        or
        $f \not\in \SIZE(2^{o(n)})$.
\end{description}

Before defining the problem formally,
let us first observe that the $\vsProblem{\E}{\SIZE(2^{o(n)})}$ Problem is
closely related to
the open question of whether $\E \not\subset \SIZE(2^{o(n)})$.
Indeed,
it is not hard to show that
$\E/O(n) \not\subset {\SIZE(2^{o(n)})}$
if and only if
$\E \not\subset {\SIZE(2^{o(n)})}$
by regarding an advice string as a part of input.
Therefore, the $\vsProblem{\E}{\SIZE(2^{o(n)})}$ Problem
is non-disjoint
if and only if
$\E \not\subset \SIZE(2^{o(n)})$.

We now define the problem formally.
According to the standard notion of advice \cite{KarpL82_advice},
the complexity class $\E/O(n)$ is defined as a subset of functions $f \colon \binset^* \to \binset$ that are defined on all the strings of any length.
Thus, ``$f \in \E/O(n)$'' does not make sense for a function $f \colon \binset^n \to \binset$.
Instead, we interpret advice
by using
the notion of Levin's resource-bounded Kolmogorov complexity \cite{Levin84_iandc_journals}
so that the notion of advice is meaningful for a finite function $f \colon \binset^n \to \binset$.
For a string $x \in \binset^*$,
let $\Kt(x)$ denote
the Levin's Kolmogorov complexity of a string $x$,
which is defined as the minimum of $|M| + \log t$ over all the programs $M$
that output $x$ in time $t$; here, $|M|$ denotes the description length of $M$.
The $\vsProblem{\E}{\SIZE(2^{o(n)})}$ Problem is formally defined as follows.
\begin{definition}
  \label{def_EvsSIZE}
  For any functions $t, s \colon \Nat \to \Nat$,
  let
  $(\YES(t(n)), \NO(s(n)))$
  denote the promise problem defined as
  \begin{align*}
    \YES(t(n)) & := \ISet{ f \in \binset^{2^n} }{ \Kt(f) \le \log t(n), n \in \Nat },
    \\
    \NO(s(n))  & := \ISet{ f \in \binset^{2^n} }{ \size(f) > s(n), n \in \Nat }.
  \end{align*}
  Here, we identity
  a function $f \colon \binset^n \to \binset$
  with its truth table representation $f \in \binset^{2^n}$,
  and $\size(f)$ denotes the minimum size of a circuit that computes $f$.

  \emph{The $\vsProblem{\E}{\SIZE(2^{o(n)})}$ Problem} is defined as the family
  $\ESet{ (\YES(2^{c n}), \NO(2^{\alpha n})) }_{c, \alpha > 0}$
  of the promise problems.
  A family $\ESet{\Pi}$ of problems is said to be \emph{solved} by a class $\mathfrak{C}$
  and denoted by $\ESet{\Pi} \in \mathfrak{C}$
  if every problem in the family is solved by some algorithm in $\mathfrak{C}$.
  %
  %
\end{definition}

The idea behind \cref{def_EvsSIZE} is that
the complexity class $\E / O(n)$ can be characterized as
the class of the functions $f = \ESet{f_n : \binset^n \to \binset}_{n\in\Nat}$
such that,
for some constant $c$, for all large $n \in \Nat$,
$\Kt(f_n) \le c n$ holds.
Indeed, $f \in \E / O(n)$ means that the truth table of $f_n$ can be described by a Turing machine of description length $O(n)$ in time $2^{O(n)}$ for all large $n$.
%
The relationship between complexity classes with advice and resource-bounded Kolmogorov complexity
will be explained in detail in \cref{section_MCLPdef},
where we interpret ``$\DTIME(t(n))/a(n)$'' as a subset of functions $f \colon \binset^n \to \binset$.
%


\subsubsection{Meta-computational view of the hardness vs randomness framework}
%
%
We show that a nearly linear-size $\AC^0 \circ \XOR$ circuit size lower bound for
solving the $\vsProblem{\E}{\SIZE(2^{o(n)})}$ Problem
exactly characterizes
the existence of a hitting set generator secure against $\AC^0 \circ \XOR$.
%
%
Here, $\AC^0 \circ \XOR$ denotes the class of constant-depth circuits that consist of unbounded-fan-in $\AND, \OR$ and $\XOR$ gates, and $\XOR$ gates are allowed to be present only in the bottom later.
\begin{theorem}
  \label{thm_AC0XOREquiv}
  The following (\crefrange{item_HSGagainstAC0XOR}{item_EvsKtLarge}) are equivalent.
  \begin{enumerate}
    \item
          \label{item_HSGagainstAC0XOR}
          There exists a hitting set generator $G = \ESet{ G_n : \binset^{O(\log n)} \to \binset^n }_{n \in \Nat} $ computable in time $n^{O(1)}$ and secure against linear-size $\AC^0 \circ \XOR$ circuits.
    \item
          \label{item_EvsSIZE}
          For all large $N \in \Nat$,
          there exists no $\AC^0 \circ \XOR$ circuit of size
          $N^{1 + o(1)}$
          that computes
          the $\vsProblem{\E}{\SIZE (2^{o(n)})}$ Problem, where $N = 2^n$ denotes the input length.
          \setcounter{TEMPenumi}{\value{enumi}}
  \end{enumerate}
  \begin{upshape}
    The condition can be equivalently stated without referring to the ``non-disjoint'' promise problem.
    Let $\MKtP[O(\log N), N^{o(1)}]$ denote the family of the promise problems
    $\MKtP[c \log N, N^\alpha]$ for constants $c, \alpha > 0$;
    here,
    for functions $s, t \colon \Nat \to \Nat$,
    $\MKtP[s(N), t(N)]$ denotes the promise problem of
    deciding whether $\Kt(x) \le s(|x|)$
    or $\Kt(x) > t(|x|)$ on input $x$.
    Then, the following are equivalent as well.
  \end{upshape}
  \begin{enumerate}
    \setcounter{enumi}{\value{TEMPenumi}}
    \item
          \label{item_EvsKt}
          For all large $N \in \Nat$,
          there exists no $\AC^0 \circ \XOR$ circuit of size
          $N^{1 + o(1)}$
          that computes\\  %
          $\MKtP[O(\log N), N^{o(1)}]$.
    \item
          \label{item_EvsKtLarge}
          For any constant $k \in \Nat$,
          for all large $N \in \Nat$,
          there exists no $\AC^0 \circ \XOR$ circuit of size
          $N^{k}$
          that computes
          $\MKtP[O(\log N), N^{o(1)}]$.
  \end{enumerate}
\end{theorem}

%
%

Observe that \cref{item_HSGagainstAC0XOR} of \cref{thm_AC0XOREquiv} implies
a strongly exponential $\AC^0$ circuit lower bound for $\E$
(\ie, $\E$ cannot be computed by $\AC^0$ circuits of size $2^{o(n)}$),
which also implies that $\EXP \not \subset \NC^1$ (see, \eg, \cite{AllenderHMPS08_siamcomp_journals,OliveiraS18_focs_conf,GoldreichW13_eccc_journals}).
The current best $\AC^0$ circuit lower bound for $\E$ is due to H{\aa}stad~\cite{Hastad87_PhD_Thesis},
who proved that the Parity function on $n$ inputs cannot be computed by depth-$d$ $\AC^0$ circuits
of size $2^{o(n^{1 / (d-1)})}$.
%
\cref{thm_AC0XOREquiv} shows that, in order to improve this state-of-the-art lower bound,
it is sufficient to prove a \emph{nearly linear} $\AC^0 \circ \XOR$ lower bound for
the $\vsProblem{\E}{\SIZE(2^{o(n)})}$ Problem.
In contrast,
the minimum circuit for computing
the $\vsProblem{\E}{\SIZE(2^{o(n)})}$ Problem
is \emph{infinity}
under the standard circuit lower bound assumption that $\E \not\subset \SIZE(2^{o(n)})$.

It is instructive to compare our results with the hardness versus randomness framework.
In order to obtain a hitting set generator
in the latter framework, we need to find an explicit function that
is hard for small circuits to compute.
In our framework,
finding an explicit hard function
corresponds to proving that the minimum circuit size for computing MCLPs is \emph{infinity}
(or, in other words, proving that there exists no circuit of \emph{any size} that computes MCLPs%
\footnote{
  This should be compared with the fact that
  any \emph{disjoint} promise problem can be computed by a circuit of size $O(2^n / n)$
  on inputs of length $n$.
}).
Our results significantly weaken the assumption needed to obtain a hitting set generator:
It suffices to show that a \emph{nearly linear} circuit
cannot find the difference between an explicit function and a hard function.

Our results can be also stated based on the case analysis.
There are two cases.
(1) When the circuit lower bound that $\E \not\subset \SIZE(2^{o(n)})$ holds,
the work of \cite{ImpagliazzoW97_stoc_conf} already implies the existence of a pseudorandom generator.
(2) Even if the circuit lower bound does fail,
\cref{thm_AC0XOREquiv} shows that
a very modest $\AC^0 \circ \XOR$ lower bound for the $\vsProblem{\E}{\SIZE(2^{o(n)})}$ Problem
(which is a disjoint promise problem under the assumption) implies
the existence of a hitting set generator.
In either case, we obtain a hitting set generator.
Our results generalize the hardness versus randomness framework in this sense.


Previously,
based on the hardness versus randomness framework,
it is known that $\E \not\subset \mathfrak{C}$ is equivalent to the existence of a pseudorandom generator secure against $\mathfrak{C}$
for a sufficiently large class $\mathfrak{C}$ (see, \eg, \cite{Fortnow01_coco_conf}).
However, in the previous approach, one needs to transform a worst-case $\mathfrak{C}$-circuit lower bound to an average-case $\mathfrak{C}$-circuit lower bound;
thus $\mathfrak{C}$ needs to be a sufficiently large so that it
can perform local decoding, which requires the majority gate \cite{ShaltielV10_siamcomp_journals}.
For any circuit class $\mathfrak{C}$ smaller than $\TC^0$,
it was not clear whether the existence of a hitting set generator secure against $\mathfrak{C}$ is equivalent to some worst-case $\mathfrak{C}$-circuit lower bound.
\cref{thm_AC0XOREquiv} establishes the first equivalence
for the circuit class $\mathfrak{C} = \AC^0 \circ \XOR$, which is smaller than $\TC^0$ \cite{Razborov87_lowerbound}.


%
%
%
%
%
Our results
can be stated without the non-standard notion of promise problem, as in \cref{item_EvsKt,item_EvsKtLarge} of \cref{thm_AC0XOREquiv}.
Indeed,
any promise problem in the family $\MKtP[O(\log N), N^{o(1)}]$
asks for approximating the $\Kt$-complexity of a given string,
and it is always a disjoint promise problem.
In our terminology,
$\MKtP[O(\log N), N^{o(1)}]$
is equivalent to the $\vsProblem{\E}{\DTIME(2^{2^{o(n)}})/2^{o(n)}}$ Problem.
Since $\SIZE(2^{o(n)}) \subset \DTIME(2^{2^{o(n)}})/2^{o(n)}$,
one can observe that
the $\vsProblem{\E}{\DTIME(2^{2^{o(n)}})/2^{o(n)}}$ Problem
is reducible to
the $\vsProblem{\E}{\SIZE(2^{o(n)})}$ Problem
via an identity map, which explains the implication from \cref{item_EvsKt} to \cref{item_EvsSIZE} in \cref{thm_AC0XOREquiv}.


%
We mention that it is not hard to prove an $\AC^0$ lower bound for
%
%
%
$\MKtP[O(\log N), N^{o(1)}]$     %
and in particular for the $\vsProblem{\E}{\SIZE(2^{o(n)})}$ Problem.
\begin{theoremref}
  {prop_AC0versusMKtP}
  For any constants $\alpha < 1, k, d \in \Nat$,
  there exists a constant $c$ such that
  $$\MKtP[c \log N, N^{\alpha}] \not\in \io \AC_d^0(N^k).$$
\end{theoremref}
This can be proved by using
the pseudorandom restriction method as in \cite{HiraharaS17_coco_conf, CheraghchiKLM19_icalp_conf, ChenHOPRS20_innovations_conf}.
(See \cref{section_AC0MKtP} for a proof.)
Extending \cref{prop_AC0versusMKtP} to the case of $\AC^0 \circ \XOR$
is sufficient for a breakthrough result in light of \cref{thm_AC0XOREquiv}.

For any classes $\mathfrak{C}, \mathfrak{D}$ of functions,
one can define the $\vsProblem{\mathfrak{C}}{\mathfrak{D}}$ Problem.
A particularly interesting problem is
the $\vsProblem{\E}{\widetilde{\AC^0}(2^{o(n)}; \frac{1}{2} - 2^{- o(n)})}$ Problem,
where $\widetilde{\mathfrak{D}}(s; \delta)$
denotes the class of functions that can be computed by a $\mathfrak{D}$-circuit of size $s$
on at least a $(1 - \delta)$ fraction of inputs.
We prove
that, if nearly linear-size $\AC^0$ circuits cannot distinguish
an explicit function from a function that cannot be approximated by small
$\AC^0$ circuits,
then
a logarithmic-seed-length hitting set generator can be obtained.
(Moreover, the converse direction is easy to prove.)


\begin{theorem}
  \label{thm_ACEquiv}
  The following are equivalent.
  \begin{enumerate}
    \item
          \label{item_EvsACnotinAC}
          %
          %
          %
          %
          %
          %
          %
          For all large $N = 2^n$,
          the $\vsProblem{\E}{\widetilde{\AC^0} (2^{o( n)} ; \frac{1}{2} - 2^{- o(n)})}$ Problem
          cannot be decided by $\AC^0$ circuits of size $N^{1 + o(1)}$.
          %
          %
          %
    \item
          \label{item_AC0HSG}
          There exists a hitting set generator $G = \ESet{ G_n : \binset^{O(\log n)} \to \binset^n }_{n \in \Nat} $ computable in time $n^{O(1)}$ and secure against linear-size $\AC^0$ circuits.
    \item
          \label{item_MKtPnotinlinearAC}
          $\MKtP[O(\log N), N - 1] \not\in \io \AC^0(N^{1 + \beta})$ for some constant $\beta > 0$.
    \item
          \label{item_MKtPnotinAC}
          $\MKtP[O(\log N), N - 1] \not\in \io \AC^0(N^{k})$ for any constant $k$.
  \end{enumerate}
\end{theorem}


%
%
%
%
%
%
%
%
%
%
%
%
%

An interesting aspect of \cref{thm_ACEquiv} is its self-referential nature; intuitively, \cref{item_EvsACnotinAC} means that $\AC^0$ circuits cannot
analyze $\AC^0$ circuits itself.  Note that self-reference is crucial for proving, \eg, time hierarchy theorems for uniform computational models.  \cref{thm_ACEquiv} provides an analogue in a non-uniform circuit model.


Why do we consider ``non-disjoint'' promise problems,
despite the fact that \cref{thm_AC0XOREquiv,thm_ACEquiv} can be stated
without referring to the non-standard notion?
A couple of remarks are in order.
First,
\cref{thm_AC0XOREquiv} is obtained by viewing (a variant of) the black-box PRG construction of Impagliazzo and Wigderson~\cite{ImpagliazzoW97_stoc_conf}
from a meta-computational perspective; thus, it is natural to state \cref{thm_AC0XOREquiv} as a connection between the existence of a hitting set generator and a lower bound for the $\vsProblem{\E}{\SIZE(2^{o(n)})}$ Problem.
Second,
an identity map reduces
$\MKtP[O(\log N), N^{o(1)}]$
to
the $\vsProblem{\E}{\SIZE(2^{o(n)})}$ Problem,
and thus it is easier to prove a lower bound for the latter problem.
Third,
the known worst-case-and-average-case equivalence between $\E \subset \SIZE(2^{o(n)})$ and $\E \subset \widetilde{\SIZE}(2^{o(n)}; \frac{1}{2} - 2^{-o(n)})$ \cite{SudanTV01_jcss_journals}
can be naturally regarded as a reduction from
the $\vsProblem{\E}{\SIZE(2^{o(n)})}$ Problem
to
the $\vsProblem{\E}{\widetilde{\SIZE}(2^{o(n)}; \frac{1}{2} - 2^{-o(n)})}$ Problem.
Indeed, \cref{thm_ACEquiv} is proved by viewing the Nisan--Wigderson pseudorandom generator from a meta-computational perspective,
and then \cref{thm_ACEquiv} is translated into \cref{thm_AC0XOREquiv}
by using the worst-case-and-average-case equivalence.
Lastly and most importantly,
in the other settings,
the non-disjointness itself provides new consequences, such as \cref{cor_KandKSAT,cor_KtvsKt}.

%



We also present a potential approach for resolving the $\RL = \L$ question.
Here, $\RL$ is the complexity class of languages that can be solved by a one-sided-error randomized $O(\log n)$-space Turing machine
that reads its random tape \emph{only once}.
A canonical approach for proving $\RL = \L$ is to construct a log-space-computable hitting set generator
of seed length $O(\log n)$ secure against $O(n)$-size read-once branching programs.
%
%
%
%
%
%
%
%
%
%
%
%
%
%
%
%
%
%
%
%
State-of-the-art results are the pseudorandom generator of seed length $O(\log^2 n)$ by Nisan~\cite{Nisan92_combinatorica_journals}
for read-once (known-order) oblivious branching programs,
and the pseudorandom generator of seed length $O(\log^3 n)$
by Forbes and Kelley~\cite{ForbesK18_focs_conf}
for read-once unknown-order oblivious branching programs.\footnote{ In the area of unconditional derandomization of space-bounded randomized algorithms, it is common to assume that a branching program is oblivious and reads the input in the fixed order.  Here, we do not assume these properties.  }

We show that a hitting set generator of seed length $O(\log n)$ can be constructed
if
nearly linear-size read-once co-nondeterministic branching programs
cannot distinguish linear-space-computable functions
from
hard functions.

\begin{theorem}
  \label{thm_ROBP}
  Suppose that,
  for some constants $\beta, \delta \in (0, 1/2)$,
  the $\vsProblem{\DSPACE(n)} { \widetilde{\SIZE}(2^{o(n)}; \delta) }$ Problem cannot
  be computed by read-once co-nondeterministic branching programs of size $N^{1 + \beta}$
  for all large input lengths $N = 2^n$.
  Then there exists a hitting set generator $G = \ESet{ G_n : \binset^{O(\log n)} \allowbreak \to \binset^n }_{n \in \Nat} $ computable in $O(\log n)$ space and secure against linear-size read-once branching programs (and, in particular, $\RL = \L$ follows; moreover, $\BPL = \L$ holds%
  \footnote{This is because of a recent result of Cheng and Hoza~\cite{ChengH20_coco_conf}, which shows $\BPL = \L$
    from the existence of a hitting set generator.}).
\end{theorem}

\cref{thm_ROBP} can be compared with a result of
Klivans and van Melkebeek \cite{KlivansM02_siamcomp_journals}.
Based on the hardness versus randomness framework,
they showed
$\BPL = \L$
under the assumption that $\DSPACE(n)$ requires
%
circuits  %
of size $2^{\Omega(n)}$.
Under the same assumption, by using a worst-case-to-average-case reduction for $\DSPACE(n)$,
it can be shown that the
$\vsProblem{\DSPACE(n)} { \widetilde{\SIZE}(2^{o(n)}; \delta) }$ Problem
is non-disjoint for any constant $\delta < 1 / 2$
%
(see \cref{prop_NonproperKS}).   %
In this case, the minimum size of a co-nondeterministic branching program for computing the MCLP is \emph{infinity};
thus, \cref{thm_ROBP} generalizes the result of \cite{KlivansM02_siamcomp_journals}.

It should be noted that
the   %
limits of
the computational power of read-once non-deterministic branching programs
are ``well understood.''
For example,
Borodin, Razborov, and Smolensky~\cite{BorodinRS93_cc_journals}
presented an explicit function that cannot be computed by any read-$k$-times
nondeterministic branching program of size $2^{o(n)}$
for any constant $k$.
\cref{thm_ROBP} shows that,
in order to resolve $\BPL = \L$,
it suffices to similarly analyze the read-once co-nondeterministic branching program size for computing the MCLP.
This approach could be useful; by using the Nechiporuck method, it can be shown that
neither nondeterministic nor co-nondeterministic branching programs of size $o(N^{1.5}/\log N)$
can compute $\MKtP$ \cite{CheraghchiHMY20}, which is a much more general lower bound than read-$k$-times nondeterministic branching programs.


We also mention that
a partial converse of \cref{thm_ROBP} is easy to prove:
If there exists a log-space-computable hitting set generator secure
against linear-size read-once nondeterministic branching programs,
then the $\vsProblem {\DSPACE(n)} { \widetilde{\SIZE}(2^{o(n)}; \delta) }$ Problem
cannot be computed by a read-once co-nondeterministic branching programs
of size $N^k$, where $N = 2^n$ and $\delta < 1 / 2$ is an arbitrary constant.
More generally, any results showing the existence of a hitting set generator
secure against $\mathfrak{C}$
must entail
a $\co \mathfrak{C}$-lower bound for MCLPs
%
(see \cref{prop_Converse}). %

%
%
%
%
%
%
%
%
%
%
%
%
%
%
%

\subsubsection{Non-trivial derandomization and lower bounds for MKtP}

Our proof techniques can be also applied to \emph{uniform} algorithms.
We consider the question of whether one-sided-error uniform algorithms can be non-trivially derandomized in time $2^{n - \omega(\sqrt{n} \log n)}$.
We say that an algorithm $A$ is a \emph{derandomization algorithm for $\DTIME(t(n))$} if,
for any machine $M$ running in time $t(n)$,
$A$ takes $\unary{n}$ and a description of $M$ as input
and
outputs $y \in \binset^n$ such that $M(y) = 1$
for infinitely many $n \in \Nat$ such that
$\Pr_{x \usample \binset^n} \squarebrackets{ M(x) = 1 } \ge \frac{1}{2}$.
Unlike the standard setting of a derandomization algorithm for non-uniform computational models,
the description length of $M$ is at most a constant independent of input length $n$; thus, our notion of derandomization algorithm is essentially equivalent
to the existence of a hitting set generator secure against $\DTIME(t(n))$.
Applying our proof techniques to this setting,
we establish the following equivalence between the existence of a derandomization algorithm for uniform algorithms and
a lower bound for approximating $\Kt$ complexity.


\begin{theorem}
  [informal; see \cref{thm_NontrivialDerandomizationFormal}]
  \label{thm_NontrivialDerandomizationInformal}
  For any constant $0 < \epsilon < 1$,
  the following are equivalent:
  \begin{enumerate}
    \item
          There exists a derandomization algorithm for $\DTIME(2^{O(\sqrt{N} \log N)})$ that runs in time $2^{N - \omega( \sqrt{N} \log N )}$.
    \item
          $\MKtP[N - \omega( \sqrt{N} \log N ), N-1] \not\in \DTIME(2^{O( \sqrt{N} \log N )})$.
    \item
          $\MKtP[N^{\epsilon}, N^{\epsilon} + \omega( \sqrt{N^{\epsilon}} \log N )] \not\in \DTIME(2^{O( \sqrt{N^{\epsilon}} \log N )})$.
  \end{enumerate}
\end{theorem}

Usually, the time complexity is measured with respect to the input size.
Our result, however, suggests that
the time complexity of $\MKtP[s(N), s(N) + {\omega}(\sqrt{N} \log N)]$ is well captured by the size parameter $s(N)$ rather than the input size $N$:
Indeed,
\cref{thm_NontrivialDerandomizationInformal} implies that
$$\MKtP[N^{\epsilon}, N^{\epsilon} + \omega( \sqrt{N^{\epsilon}} \log N )] \in \DTIME(2^{O( \sqrt{N^{\epsilon}} \log N )})$$
is equivalent to
$$\MKtP[N^{\delta}, N^{\delta} + \omega( \sqrt{N^{\delta}} \log N )] \in \DTIME(2^{O( \sqrt{N^{\delta}} \log N )})$$
for any $0 < \epsilon, \delta < 1$.


%
%
%
%

\cref{thm_NontrivialDerandomizationInformal} highlights the importance of a lower bound for $\MKtP$.
In fact,
it is a long-standing open question whether $\MKtP \not\in \P$, despite the fact that $\MKtP$ is an $\EXP$-complete problem under non-uniform reductions \cite{AllenderBKMR06_siamcomp_journals}.
Towards resolving this open question,
we present several results that might be useful.
We show that
some ``non-disjoint'' promise problem can be solved
under the assumption that $\MKtP \in \P$.
\begin{theorem}
  \label{thm_KtvsKt}
  Assume that $\MKtP \in \P$.
  Then, there exists a $2^{\widetilde{O}(\sqrt{n})}$-time algorithm that solves the $\vsProblem{\Kt}{\KtTIME{t}}$ Problem, which is defined as follows:
  Given a string $x \in \binset^*$ of length $n$ and a parameter $s \in \Nat$,
  distinguish whether
  $\Kt(x) \le s$
  or
  $\KtTIME{t}(x) \ge s + O(\sqrt{s} \log n + \log^2 n)$,
  where $t := n^{\log n}$.
\end{theorem}
Using the disjointness of the $\vsProblem{\Kt}{\KtTIME{t}}$ Problem and setting $s := \Kt(x)$, we obtain
\begin{corollary}
  \label{cor_KtvsKt}
  If $\MKtP \in \P$,
  then
  $\KtTIME{t}(x) \le \Kt(x) + O(\sqrt{\Kt(x)} \log n + \log^2 n)$
  for any $x \in \binset^n$,
  where  $t := n^{\log n}$.
\end{corollary}

Since $\Kt(x) \le \KtTIME{t}(x) + O(\log^2 n)$ holds for $t := n^{\log n}$ unconditionally,
\cref{cor_KtvsKt} shows that
$\Kt(x)$ and $\KtTIME{t}(x)$ are
equal up to an additive term of $\widetilde{O}(\sqrt{n})$
under the assumption that $\MKtP \in \P$.
In particular,
in order to resolve $\MKtP \not\in \P$,
it suffices to prove that
$\KtTIME{t}(x)$ cannot be approximated
within an additive error $\widetilde{O}(\sqrt{n})$
in polynomial time.
It was shown in \cite{Hirahara20_stoc_conf} that
there is no polynomial-time algorithm
that computes
$\KtTIME{t}(x)$ \emph{exactly} on input $x$.
The large additive error term $\widetilde{O}(\sqrt{n})$ prevents us from combining the lower bound proof technique of \cite{Hirahara20_stoc_conf} and \cref{cor_KtvsKt}.


\subsubsection{Related work: Minimum Circuit Size Problem}
The definitions of MCLPs  are inspired by
the \emph{Minimum Circuit Size Problem} (MCSP).
%
%
While the history of $\MCSP$ is said to date back to as early as 1950s \cite{Trakhtenbrot84_annals_journals},
its importance was not widely recognized
until
Kabanets and Cai~\cite{KabanetsC00_stoc_conf}
named the problem as $\MCSP$
and investigated it
based on the natural proof framework of Razborov and Rudich~\cite{RazborovR97_jcss_journals}.
The task of $\MCSP$
is
to decide whether there exists a size-$s$ circuit that computes $f$,
given the truth table of a function $f \colon \binset^n \to \binset$
and a size parameter $s \in \Nat$.
It turned out
that
$\MCSP$
is one of the central computational problems
in relation to wide research areas of complexity theory,
including
circuit lower bounds \cite{RazborovR97_jcss_journals,OliveiraS18_focs_conf},
learning theory~\cite{CarmosinoIKK16_coco_conf},
cryptography~\cite{RazborovR97_jcss_journals, AllenderBKMR06_siamcomp_journals, AllenderH19_toct_journals, Santhanam20_innovations_conf},
and
average-case complexity~\cite{Hirahara18_focs_conf}.


It has been recognized that $\MCSP$ lacks one desirable mathematical property --- \emph{monotonicity} with respect to an underlying computational model.
$\MCSP$ can be defined for any circuit classes $\mathfrak{C}$;
for example,
$\mathfrak{C}\hyphen\MCSP$ stands for a version of $\MCSP$ where the task is to find the minimum $\mathfrak{C}$-circuit size;
$\MCSP^A$ stands for the minimum $A$-oracle circuit size problem.
We are tempted to conjecture that,
as a computational model becomes stronger,
the corresponding minimization problem becomes harder;
\eg,
$\MCSP^A$ should be harder than $\MCSP$ for any oracle $A$.
However, this is not the case ---
Hirahara and Watanabe~\cite{HiraharaW16_coco_conf}
showed that there exists an oracle $A$ such that $\MCSP \not\le^p_T \MCSP^A$
unless $\MCSP \in \P$.
Moreover, $\mathrm{DNF}\hyphen\MCSP$ \cite{Masek79,AllenderHMPS08_siamcomp_journals} and $(\mathrm{DNF}\circ\mathrm{XOR})\hyphen\MCSP$ \cite{HiraharaOS18_coco_conf}
are known to be $\NP$-complete, whereas $\NP$-completeness of $\MCSP$ is a long-standing open question.
%

Why does the monotonicity of $\MCSP$ fail?
For two circuit classes $\mathfrak{C}$ and $\mathfrak{D}$,
define the $\vsProblem{\mathfrak{C}}{\mathfrak{D}}$ Problem
to be the promise problem of distinguishing
whether $f$ has a $\mathfrak{C}$-circuit of size $s$
or $f$ does not have a $\mathfrak{D}$-circuit of size $s$,
given the truth table of a function $f \colon \binset^n \to \binset$
and a size parameter $s$.
Then,
$\mathfrak{C}\hyphen\MCSP$ can be regarded as a special case of the $\vsProblem{\mathfrak{C}}{\mathfrak{D}}$ Problem where $\mathfrak{C} = \mathfrak{D}$.
It is easy to observe that
the $\vsProblem{\mathfrak{C}}{\mathfrak{D}}$ Problem is reducible to
the $\vsProblem{\mathfrak{C'}}{\mathfrak{D'}}$ Problem via an identity map
if $\mathfrak{C} \subset \mathfrak{C'}$
and
$\mathfrak{D} \supset \mathfrak{D'}$;
thus, MCLPs have monotonicity properties in this sense.
In contrast,
the monotonicity property of MCLPs fails when
$\mathfrak{C} = \mathfrak{D} \subset \mathfrak{C}' = \mathfrak{D}'$,
which corresponds to the case of $\MCSP$.

In an attempt to remedy the monotonicity issue,
Hirahara and Santhanam~\cite{HiraharaS17_coco_conf}
observed that
average-case complexity of $\MCSP$ is monotone increasing.
%
%
Carmosino, Impagliazzo, Kabanets, and Kolokolova~\cite{CarmosinoIKK16_coco_conf}
implicitly showed
that
the complexity of $\MCSP$ is monotone increasing under non-black-box reductions.

In contrast,
MCLPs incorporate
the monotonicity property
in the definition itself,
which makes a mathematical theory cleaner.
For example,
recall that
it can be shown that
$\E \not\subset \SIZE(2^{o(n)})$
if and only if
$\E \not\subset \widetilde{\SIZE}(2^{o(n)}; \frac{1}{2} - 2^{-o(n)})$
by using error-correcting codes \cite{SudanTV01_jcss_journals}.
Viewing this equivalence from a meta-computational perspective,
it can be interpreted as an efficient reduction
from
the $\vsProblem{\E}{\SIZE(2^{o(n)})}$ Problem
to
the $\vsProblem{\E}{\widetilde{\SIZE}(2^{o(n)}; \frac{1}{2} - 2^{-o(n)})}$ Problem
%
(see \cref{thm_ECCReduction}). %
A similar reduction was presented for $\MCSP$ under the assumption that $\EXP \subset \Ppoly$ \cite{ChenHOPRS20_innovations_conf};
finding such a reduction for $\MCSP$ requires certain creativity,
whereas working with the definition of MCLPs makes it trivial to find the reduction.
%


\subsubsection{Related work: Hardness Magnification}

A recent line of work \cite{OliveiraS18_focs_conf, OliveiraPS19_coco_conf, McKayMW19_stoc_conf, ChenJW19_focs_conf, ChenHOPRS20_innovations_conf} exhibit surprising phenomena, which are termed as ``hardness magnification phenomena.''
Oliveira, Santhanam, and Pich~\cite{OliveiraS18_focs_conf,OliveiraPS19_coco_conf}
showed that
very weak lower bounds for $\MCSP$ and related problems
are sufficient for resolving
long-standing open questions
about circuit lower bounds, such as $\EXP \not\subset \NC^1$.
An interesting aspect of hardness magnification phenomena is that,
as argued in \cite{AllenderK10_jacm_journals, OliveiraS18_focs_conf},
the argument does not seem to be subject to the natural proof barrier of Razborov and Rudich~\cite{RazborovR97_jcss_journals},
which is one of the major obstacles of complexity theory.
Our results can be seen as a new interpretation of hardness magnification phenomena
from the perspective of
a construction of a hitting set generator.

It is worthwhile to point out that our reductions have very similar structures
to hardness magnification phenomena.
For example, it was shown in \cite{OliveiraS18_focs_conf, OliveiraPS19_coco_conf} that, for a parameter $s = s(N)$,
the problem $\MKtP[s, s + O(\log N)]$
can be solved by
an $\AND_{O(N)} \circ D_{\poly(s)} \circ \XOR$ circuit,
where $D_{\poly(s)}$ is an oracle gate computable in $\EXP$ that takes an input of length $\poly(s)$.
This reduction shows that
$\EXP \subset \Ppoly$ implies $\MKtP[s, s + O(\log N)] \in \SIZE(N \cdot \poly(s))$.
Taking its contrapositive,
it can be interpreted as a hardness magnification phenomenon:
for $s(N) \ll N$, a (seemingly) weak lower bound $\MKtP[s, s + O(\log N)] \not\in \SIZE(N \cdot \poly(s))$
can be ``hardness magnified'' to
a strong circuit lower bound $\EXP \not\subset \Ppoly$.

\cref{thm_AC0XOREquiv} has exactly the same structure with the reduction mentioned above.
Given a circuit $D_{N^{o(1)}}$ that avoids a hitting set generator,
we construct a nearly linear-size $\AND_{O(N)} \circ D_{N^{o(1)}} \circ \XOR$ circuit that computes
the $\vsProblem{\E}{\SIZE(2^{o(n)})}$ Problem.
Thus, our reductions are as efficient as the reductions presented
in the line of work of hardness magnification phenomena.

More importantly, our results significantly strengthen the consequences of hardness magnification:
Not only circuit lower bounds, but also hitting set generators can be obtained.
This is especially significant for the case of read-once branching programs.
Since
there is already an exponential size lower bound
for read-once branching programs \cite{BorodinRS93_cc_journals},
it does not make sense to try to hardness-magnify
a lower bound for read-once branching programs.
In contrast,
our results (\cref{thm_ROBP})
indicate that
a nearly linear size lower bound for co-nondeterministic read-once branching programs
is enough for resolving $\BPL = \L$.

An intriguing question about hardness magnification is this:
By using hardness magnification phenomena,
can we prove any new consequences, such as
circuit lower bounds or derandomization?
%
\cref{thm_ROBP} adds a new computational model, \ie, co-nondeterministic read-once branching programs,
for exploring this question.

Chen, Hirahara, Oliveira, Pich, Rajgopal, and Santhanam~\cite{ChenHOPRS20_innovations_conf}
proposed a barrier for the question, termed as a ``locality barrier.''
Briefly speaking, the idea there is to regard hardness magnification phenomena
as a (black-box) reduction to oracles with small fan-in,
and then to show that most circuit lower bound proofs can be extended to rule
out such a reduction;
thus, such a circuit lower bound proof technique cannot be combined
with hardness magnification phenomena.
A salient feature of our reductions is that our reductions are \emph{non-black-box}
in the sense that we exploit the efficiency of oracles;
the non-black-box property appears in the definition of $\No$ instances of the $\vsProblem{\mathfrak{C}}{\mathfrak{D}}$ Problem.
In light of this,
our results provide a potential approach for bypassing the locality barrier:
Try to develop a circuit lower bound proof technique that crucially exploits
the structure of the $\No$ instances of  the $\vsProblem{\mathfrak{C}}{\mathfrak{D}}$ Problem.
The existing circuit lower bound proof techniques for $\MCSP$ and related problems fail to exploit such a structure.

%

%
%



%
%
%
%
%
%
%
%
%
%
%
%
%
%
%
%
%
%
%
%
%
%
%
%
%
%
%
%
%
%
%
%
%
%
%


\subsection{Proof techniques: Meta-computational view of PRG constructions}
\label{subsection:proofsketch}

All of our results are given by a \emph{single principle} --- that views
any black-box pseudorandom generator construction from a meta-computational perspective.
The differences among our theorems simply originate from the fact that we use a different black-box pseudorandom generator construction.
The underlying principle is this:
%
%
%
%
%
%
%
\begin{quote}
  Any black-box construction of
  a pseudorandom generator $G^f$
  based on a hard function $f \not\in \mathcal{R}$
  gives rise to
  a non-black-box security reduction
  for a hitting set generator
  based on the hardness of a non-disjoint promise problem
  (\eg, the $\vsProblem{\E}{\mathcal{R}}$ Problem).
\end{quote}

For the purpose of exposition, we take a specific PRG construction of Impagliazzo and Wigderson~\cite{ImpagliazzoW97_stoc_conf},
and prove a connection between the PRG construction and the $\vsProblem{\E}{\SIZE(2^{o(n)})}$ Problem.
Specifically, in this section, we present a self-contained proof of the following result.
\begin{theorem}
  \label{thm_ProofTechniques}
  There exists a universal constant $\beta$ such that the following are equivalent.
  \begin{enumerate}
    \item The $\vsProblem{\E}{\SIZE(2^{o(n)})}$ Problem cannot be solved by a circuit of size $N^{\beta}$ on inputs of length $N = 2^n$ for all large $N$.
          \label{item_ProofTechniquesLowerBound}
    \item There exists a hitting set generator $H = \ESet{ H_m : \binset^{O(\log m)} \to \binset^m }_{m \in \Nat}$ computable in time $m^{O(1)}$ and secure against linear-size circuits.
          \label{item_ProofTechniquesHSG}
  \end{enumerate}
\end{theorem}
The implication from \cref{item_ProofTechniquesHSG} to \cref{item_ProofTechniquesLowerBound}
is by now a standard fact in the literature of $\MCSP$ and Kolmogorov complexity \cite{RazborovR97_jcss_journals,AllenderBKMR06_siamcomp_journals}.
The implication from \cref{item_ProofTechniquesLowerBound} to \cref{item_ProofTechniquesHSG}
is shown by applying our principle to the following black-box pseudorandom generator construction of \cite{ImpagliazzoW97_stoc_conf}.
\begin{lemma}
  [Impagliazzo and Wigderson~\cite{ImpagliazzoW97_stoc_conf}]
  \label{lemma_IWBBConstruction}
  There is a constant $\beta$ such that,
  for any constant $\alpha > 0$,
  there exist
  a constant $0 < \alpha' < 1/2$ and
  a ``black-box pseudorandom generator construction''
  $G^f \colon \binset^{O(n)} \to \binset^{2^{\alpha' n}}$
  that takes a function $f \colon \binset^n \to \binset$ for any $n \in \Nat$
  and satisfies the following.
  \begin{description}
    \item[Explicitness]
          $G^f(z)$ is computable in time $2^{O(n)}$
          given the truth table of $f$ and a seed $z$.
    \item[Constructivity]
          For every fixed seed $z \in \binset^{O(n)}$,
          there is a circuit $C_z$ of size $2^{(\beta - 2) n}$ that
          takes as input the $2^n$-bit truth table of a function $f$ and
          computes $G^f(z)$.
    \item[Security]
          For every function $D \colon \binset^{2^{\alpha' n}} \to \binset$ that
          $\frac{1}{4}$-\emph{distinguishes}%
          \footnote{That is, $\Pr_z \squarebrackets{ D(G^f(z)) = 1} - \Pr_w \squarebrackets{D(w) = 1} \ge 1/4$; see \cref{section_preliminaries} for the definition of PRG.}
          the output distribution
          of $G^f(\blank)$ from the uniform distribution,
          then
          $f \in \SIZE^D(2^{(\alpha - \alpha') n})$.
  \end{description}
\end{lemma}

%
%
%
%
%
%
%
%
%
%
%
%
%
%
%
%
%
Here, by ``black-box'',%
\footnote{In the taxonomy of \cite{ReingoldTV04_tcc_conf}, $G^{(\blank)}$ is \emph{fully black-box}
  in the sense that the construction from $f$ to $G^f$ also works for every function $f$.}
we mean that the security of the PRG is proved by a (black-box) reduction, \ie,
the security reduction works for \emph{every} function $D$;
otherwise, the reduction is called non-black-box (\eg, in the case when the reduction works only if the oracle is efficient).
This is in the same spirit with the non-black-box reduction of \cite{Hirahara18_focs_conf},
which overcomes the black-box reduction barrier of Bogdanov and Trevisan~\cite{BogdanovT06_siamcomp_journals}.
We explain below how a black-box PRG construction gives rise to a \emph{non-black-box} security reduction of some
hitting set generator.

%
%
%
%
%
%
%
%
%
%
%
%
%
%
%
%
%
%

\begin{proof}
  [Proof of \cref{thm_ProofTechniques}]
  (\cref{item_ProofTechniquesLowerBound} $\Rightarrow$ \ref{item_ProofTechniquesHSG})
  The goal is to construct some secure hitting set generator
  $H = \ESet{ H_m \colon \allowbreak \binset^{c_0 \log m} \to \binset^m }_{m \in \Nat}$
  for some constant $c_0$.
  As a choice of $H$, we simply take
  a ``universal'' hitting set generator:
  Let $U$ be a universal Turing machine, \ie, a machine that simulates every Turing machine efficiently.
  Then we define $H_m(z)$ to be the output of $U$ on input $z$
  if $U$ outputs a string of length $m$ in $2^{|z|}$ steps, where $z \in \binset^{c_0 \log m}$.
  The important property of $H_m$ is that
  every string $x \in \binset^m$
  with $\Kt(x) \le c_0 \log m$ is contained in the image of $H_m$.
  (The choice of $H$ is universal, in the sense that
  the existence of some exponential-time computable HSG
  implies that $H$ is also secure.)

  The strategy for proving the security of a hitting set generator $H$
  is to regard a function $f \colon \binset^n \to \binset$ as an input of the $\vsProblem{\E}{\SIZE(2^{o(n)})}$ Problem,
  and
  to view the black-box PRG construction $G^f$
  as a (non-black-box) reduction that proves the security of $H$.
  %

  We claim that $H$ is a hitting set generator
  secure against linear-size circuits,
  under the assumption
  that
  the $\vsProblem{\E}{\SIZE(2^{o(n)})}$ Problem cannot be solved
  by small circuits.
  To this end, we present a non-black-box reduction $R^{(\blank)}$ from
  the task of solving the $\vsProblem{\E}{\SIZE(2^{o(n)})}$ Problem
  to the task of ``avoiding'' $H$.
  That is, we prove the contrapositive of \cref{item_ProofTechniquesLowerBound} $\Rightarrow$ \ref{item_ProofTechniquesHSG}:
  Assume that $H$ is not secure (for every constant $c_0$, which decides the seed length of $H$).
  Then,
  for infinitely many $m \in \Nat$,
  there exists a linear-size circuit $D$ that \emph{avoids} $H_m$, \ie,
  every string in the image of $H_m$ is rejected by $D$ whereas $D$ accepts at least a half of all the $m$-bit inputs.
  We claim that there exists a small oracle circuit $R^D$ that solves
  the $\vsProblem{\E}{\SIZE(2^{o(n)})}$ Problem.
  The randomized circuit $R^D$ is extremely simple:
  \begin{description}
    \item[
          A randomized algorithm $R^D$ for solving the $\vsProblem{\E}{\SIZE(2^{o(n)})}$ Problem
          with $D$ oracle]
          \hspace{0em}

          \emph{
            Given $f$ as an input,
            pick a random seed $z$ of $G^f$, and accept if and only if $D(G^f(z)) = 0$.
          }
  \end{description}
  More formally,
  recall that the $\vsProblem{\E}{\SIZE(2^{o(n)})}$ Problem is defined as
  the family of promise problems $\Pi_{c, \alpha} = (\YES(2^{c n}), \NO(2^{\alpha n}))$ for every $c, \alpha > 0$
  as in \cref{def_EvsSIZE}.
  We need to prove that, for any constants $c, \alpha > 0$,
  for infinitely many $n \in \Nat$,
  there exists a randomized circuit $R^D$ that solves $\Pi_{c, \alpha}$
  on inputs of length $N = 2^n$.
  We take $\alpha'$ of \cref{lemma_IWBBConstruction} (depending on $c, \alpha$), and
  choose $n \in \Nat$
  so that
  $m = 2^{\alpha' n}$.%
  \footnote{More precisely, we consider infinitely many $m$ such that $D$ avoids $H_m$;
    then, setting $n := (\log m) / \alpha'$, we prove that the reduction $R^D$ solves
    $\Pi_{c, \alpha}$ on inputs of length $n$.}

  The correctness of the reduction $R^D$ can be proved as follows.
  Assume that $f$ is a $\Yes$ instance of the promise problem $\Pi_{c, \alpha}$;
  this means that $\Kt(f) \le \log 2^{c n} = c n$.
  It follows from the explicitness of $G^f$ (\cref{lemma_IWBBConstruction})
  that
  we can take a large constant $c_0$ such that
  $\Kt(G^f(z)) \le c_0 \alpha' n = c_0 \log m$ for every seed $z \in \binset^{O(n)}$.
  By the definition of the universal hitting set generator $H_m$,
  one can observe that $G^f(z) \in \binset^m$ is in the image of $H_m$
  (for infinitely many $m$);
  therefore, $G^f(z)$ is rejected by $D$
  for every $z$,
  and hence the reduction $R^D$ accepts $f$ with probability $1$.

  Conversely,
  we prove that any $\No$ instance $f$ of $\Pi_{c, \alpha}$
  is rejected by the reduction $R^D$ with probability at least $\frac{1}{4}$.
  Intuitively, this is because of the fact that if $f$ is a hard function, then
  $D$ cannot distinguish $G^f(\blank)$
  from the uniform distribution.
  More formally, we prove the contrapositive.
  Assume that $R^D$ rejects $G^f(z)$ with probability at most $\frac{1}{4}$.
  This means that
  $$
    \Pr_z[ D(G^f(z)) = 1 ] \le \frac{1}{4}.
  $$
  On the other hand, since $D$ avoids $H_m$, we have $\Pr_w [ D(w) = 1 ] \ge \frac{1}{2}$.
  Therefore,
  $$
    \Pr_w [D(w) = 1] - \Pr_z[ D(G^f(z)) = 1 ] \ge \frac{1}{4},
  $$
  which means that
  $\neg D$ $\frac{1}{4}$-distinguishes $G^f$ from the uniform distribution.
  By using the black-box security property of $G^f$ (\cref{lemma_IWBBConstruction}),
  we obtain that $f \in \SIZE^D(2^{(\alpha - \alpha') n})$.
  Since $D$ is a linear-size circuit (\ie, of size $m = 2^{\alpha' n}$), we conclude that $f \in \SIZE(2^{\alpha n})$,
  which means that $f$ is not a \No instance of the promise problem $\Pi_{c, \alpha}$.
  Note here that we rely on the efficiency of $D$, which makes the security proof of the HSG $H$ non-black-box.


  To summarize,
  the promise problem $\Pi_{c, \alpha}$
  can be solved by the randomized circuit $R^D$
  of size $N^{\beta - 2} + 2^{\alpha' n} \le N^{\beta - 1.5}$, where the size bound follows from
  the constructivity of $G^f$ (\cref{lemma_IWBBConstruction}).
  By using the standard trick of Adleman~\cite{Adleman78_focs_conf},
  the randomized circuit $R^D$ of size $N^{\beta - 1.5}$ can be converted to a
  deterministic circuit of size $O(N^{\beta-0.5})$
  (by amplifying the success probability and fixing a good random coin flip).
  We conclude that
  the $\vsProblem{\E}{\SIZE(2^{o(n)})}$ Problem can be solved by a circuit of size $N^{\beta}$.

  (\cref{item_ProofTechniquesHSG} $\Rightarrow$ \ref{item_ProofTechniquesLowerBound})
  Conversely,
  if there exists a hitting set generator $H$ secure against linear-size circuits,
  then the $\vsProblem{\E}{\SIZE(2^{o(n)})}$ Problem cannot be solved by circuits of size $N^{\beta'}$
  for any constant $\beta'$.
  Indeed, we can prove the contrapositive (roughly) as follows.
  Given a circuit $D$ that solves the $\vsProblem{\E}{\SIZE(2^{o(n)})}$ Problem,
  we claim that the circuit $\neg D$ avoids $H$.
  On one hand,
  any string $x$ in the image of $H$ satisfies that $\Kt(x) = O(\log |x|)$;
  therefore, $x$ is a \Yes instance of
  the $\vsProblem{\E}{\SIZE(2^{o(n)})}$ Problem,
  and hence rejected by $\neg D$.
  On the other hand,
  a string $x$ chosen uniformly at random
  has circuit complexity at least $2^{\alpha n}$ with high probability
  for any constant $\alpha < 1$;
  therefore, $\neg D$ accepts most inputs.
  (See \cref{prop_Converse} for a formal proof.)
\end{proof}


Note that the proof above shows a generic connection
between
a black-box pseudorandom generator construction
and
a ``non-disjoint'' promise problem.
The efficiency of the security reduction depends on the choice of a
black-box PRG construction.
For example, \cref{thm_AC0XOREquiv} is proved by giving an efficient implementation
of the black-box PRG construction $G^{(\blank)}$ of \cite{ImpagliazzoW97_stoc_conf}
such that $G^f(z)$ is computable by one layer of XOR gates for every fixed seed $z$.
In the rest of this paper, we present some of the instantiations of
the proof ideas above
based on several specific constructions of PRGs;
however, we emphasize that our reductions are not limited to those specific instantiations,
and new black-box PRG constructions can lead to a more efficient reduction
and a ``non-disjoint'' promise problem that is easy to analyze.

\subsection{Perspective: Meta-computational view of complexity theory}

More broadly, we propose to view complexity theory from a meta-computational perspective.
%

In order to explain the view, it is helpful to regard an algorithm that tries to solve MCLPs
as a malicious prover that tries to falsify a circuit lower bound.
To be more specific, consider the $\vsProblem{\E}{\SIZE(2^{o(n)})}$ Problem.
As we explained earlier,
the existence of
any algorithm (of any complexity) that solves the $\vsProblem{\E}{\SIZE(2^{o(n)})}$ Problem
implies that $\E \subset \SIZE(2^{o(n)})$,
and moreover the converse direction is also true (for non-uniform algorithms).
In this sense, any algorithm that solves an MCLP can be regarded as an adversary that falsifies circuit lower bounds such as $\E \not\subset \SIZE(2^{o(n)})$.

Complexity theory can be regarded as a game between us (\ie, complexity theorists, who try to prove circuit lower bounds) and provers (\ie, algorithms that solve MCLPs).
We lose the game if some prover can solve MCLPs (and hence circuit lower bounds fail).
We win the game if we find an explicit function whose circuit complexity is high.
This is equivalent to finding a witness for the non-disjointness of the $\vsProblem{\E}{\SIZE(2^{o(n)})}$ Problem,
and thus it is equivalent to showing that there exists no prover that can solve the MCLP.
In other words,
prior to our work,
we implicitly tried to fight against \emph{every prover without any restriction on efficiency}.

What we showed in this work is that we do not have to fight against every non-efficient prover.
Instead,
in order to obtain a circuit lower bound (which is implied by the existence of a hitting set generator),
it suffices to show that no \emph{efficient} algorithms such as nearly linear-size $\AC^0 \circ \XOR$ circuits
can find
the difference between explicit functions and hard functions.
In principle,
it should be easier to prove a nearly linear circuit size lower bound for some problem
when we believe that the problem does not admit any algorithm (because of the non-disjointness).
While we have not found any existing method for proving such a lower bound that is sufficient for breakthrough results,
we believe that this is simply because of the fact that MCLPs were not investigated explicitly.
We leave as an important open question to develop a proof technique to analyze MCLPs.


\section{Preliminaries}
\label{section_preliminaries}


\subsection{Notation}

For a Boolean function $f \colon \binset^n \to \binset$,
we denote by $\truthtable(f)$ the truth table of $f$, \ie, the concatenation of $f(x)$ for all $x \in \binset^n$ in the lexicographical order.
We often identify $f \colon \binset^n \to \binset$ with $\truthtable(f) \in \binset^{2^n}$.
We denote by $\size(f)$ the minimum circuit size for computing the Boolean function $f \colon \binset^n \to \binset$.
For a parameter $\delta$, we denote by $\ApproxCC{\delta}{f}$
the minimum size of a circuit $C$ such that $f(x) = C(x)$ on at least a $(1 - \delta)$ fraction of the $n$-bit inputs $x$.
%
%
%

%
%
%
%
%
%

For a function $f \colon \binset^n \to \binset$ and an integer $k \in \Nat$, we denote by $f^{k} \colon (\binset^n)^k \to \binset^k$
the direct product of $f$.
We denote by $f^{\oplus k} \colon (\binset^n)^k \to \binset$ the function $\oplus_k \circ f^{k}$, where $\oplus_k$ is the parity function on $k$-bit inputs.


\subsection{Pseudorandomness}
Let $G \colon \binset^d \to \binset^m$
and $D \colon \binset^m \to \binset$
be functions.
For an $\epsilon > 0$,
we say that $D$ $\epsilon$-\emph{distinguishes}
the output distribution of $G(\blank)$ from
the uniform distribution
if
$$
  \Pr_{z \usample \binset^d} [ D(G(z)) = 1 ]
  - \Pr_{w \usample \binset^m} [ D(w) = 1 ]
  \ge
  \epsilon
$$
In this case, we refer $D$ as an $\epsilon$-\emph{distinguisher} for $G$.
Conversely, $G$ is said to $\epsilon$-\emph{fool} $D$ if $D$ is not an $\epsilon$-distinguisher for $G$.
Similarly,
we say that $D$ $\epsilon$-\emph{avoids}
$G$
if
$\Pr_{w \usample \binset^m} [D(w) = 1] \ge \epsilon$
and
$D(G(z)) = 0$ for every $z \in \binset^d$.
By default, we assume that $\epsilon := \frac{1}{2}$.

For a circuit class $\mathfrak{C}$ and functions $s \colon \Nat \to \Nat$ and $\epsilon \colon \Nat \to [0,1]$,
a family of functions $G = \ESet{ G : \binset^{s(n)} \to \binset^n }_{n \in \Nat}$
is said to be a \emph{pseudorandom generator} that $\epsilon$-fools $\mathfrak{C}$
if,
for all large $n \in \Nat$ and any circuit $C \in \mathfrak{C}$ of $n$ inputs,
$G$ $\epsilon(n)$-fools $C$.
We say that $G$ is a \emph{hitting set generator} secure against $\mathfrak C$
if for all large $n \in \Nat$,
there is no circuit $D \in \mathfrak C$ on $n$ inputs
that avoids $G_n$.


\subsection{Circuits}
We measure circuit size by the number of gates (except for the input gates).
For a circuit type $\mathfrak{C}$ and $s \colon \Nat \to \Nat$ and $\delta \in [0, 1]$,
we denote by $\widetilde{\mathfrak{C}}(s; \delta)$
the class of functions $f \colon \binset^n \to \binset$
such that
there exists a circuit of size $s(n)$
such that
$\Pr_{x \sim \binset^n} \squarebrackets{ f(x) = C(x) } \ge 1 - \delta$.
We define $\mathfrak{C}(s) := \widetilde{\mathfrak{C}}(s; 0)$.
For the standard circuit class,
we use the notation $\widetilde{\SIZE}(s; \delta)$ and $\SIZE(s)$.


\subsection{Time-bounded kolmogorov complexity}
We fix an efficient universal Turing machine $U$.
Time-bounded Kolmogorov complexity is defined as follows.
\begin{definition}
  [Time-bounded Kolmogorov Complexity]
  The time-$t$ $A$-oracle Kolmogorov complexity of a string $x \in \binset^*$
  is defined as
  $$
    \KtTIME[A]{t}(x) := \min \ISet{ |d| } { \text{$U^{A}$ outputs $x$ in $t$ steps on input $d \in \binset^*$ } },
  $$
  where $A$ is an oracle and $t \in \Nat$.
\end{definition}


\section{Meta-computational view of cryptographic PRG constructions}

In this section,
we apply the meta-computational principle to cryptographic pseudorandom generator constructions.
The worst-case-to-average-case reduction for $\KvsK{\SAT}$
is given in \cref{section_Heuristica}.
We view the pseudorandom generator construction of
H{\aa}stad, Impagliazzo, Levin, and Luby~\cite{HastadILL99_siamcomp_journals}
from a meta-computational perspective
in \cref{section_HILL}.

\subsection{\texorpdfstring{$\KvsK{\SAT}$}{Gap(K(SAT) vs K)}}
\label{section_Heuristica}

We present a proof of the following result.

\begin{remindertheorem}
  {thm_HeuristicaIntro}
  Let $A$ be any oracle.
  If $\Dist\NP^A \subset \Avg\P$,
  then $\KvsK{A} \in \P$.
\end{remindertheorem}



%
%
%
%
%
%
%
%
%
%
%
%
%
%

At the core of the proof of \cref{thm_HeuristicaIntro} is to use a black-box PRG construction whose advice complexity is small.
%
%
Following \cite{Hirahara20_stoc_conf},
we observe that a $k$-wise direct product generator, which is one of the simplest constructions of pseudorandom generators,
has small advice complexity.

\begin{theorem}
  [$k$-wise direct product generator \cite{Hirahara20_stoc_conf}]
  \label{thm_DPConstruction}
  For any parameters $n, k \in \Nat$ and $\epsilon > 0$ with $k \le 2n$,
  there exists a ``black-box pseudorandom generator construction'' $(\DP_k, A^{(\blank)}, R^{(\blank)})$
  satisfying the following.
  \begin{itemize}
    \item
          $\DP_k \colon \binset^n \times \binset^d \to \binset^{d + k}$
          is called a \emph{pseudorandom generator function}
          and computable in time $\poly(n / \epsilon)$.
          \uprightbrackets{Intuitively, $\DP_k$ takes the $n$-bit truth table of a candidate ``hard'' function as well as a $d$-bit seed and outputs a pseudorandom sequence of length $d + k$.}
    \item
          $A^D \colon \binset^n \times \binset^r \to \binset^a$
          is called an \emph{advice function}
          and computable in time $\poly(n / \epsilon)$ given oracle access to
          $D \colon \binset^{d+k} \to \binset$.
    \item
          $R^D \colon \binset^a \times \binset^r \to \binset^n$
          is called a \emph{reconstruction procedure}
          and computable in time $\poly(n / \epsilon)$ given oracle access to
          $D \colon \binset^{d+k} \to \binset$.
    \item
          The seed length $d$ is at most $\poly(n / \epsilon)$,
          the advice complexity $a$ is at most $k + O(\log (k / \epsilon))$,
          and
          the randomness complexity $r$ is at most $\poly(n / \epsilon)$.
    \item
          For any string $x$ and any function $D$ that $\epsilon$-distinguishes
          the output distribution of $\DP_k(x, \blank)$
          from the uniform distribution,
          it holds that
          $$
            \Pr_{w \usample \binset^r}[ R^D(A^D(x, w), w) = x ] \ge 3 / 4.
          $$
  \end{itemize}
\end{theorem}

For completeness, we present a proof of \cref{thm_DPConstruction}.
We make use of the following list-decodable error-correcting code,
which can be constructed by concatenating a Reed-Solomon code with an Hadamard code.
\begin{lemma}
  [\cite{Sudan97_jc_journals}; see, \eg, \cite{Vadhan12_fttcs_journals}]
  \label{lemma_ListDecodable}
  For any parameters $n \in \Nat$ and $\epsilon > 0$,
  there exists a function $\Enc_{n, \epsilon} \colon \binset^n \to \binset^{\widehat{n}}$ such that
  \begin{itemize}
    \item
          $\widehat{n} = 2^\ell$ for some integer $\ell \in \Nat$ and $\widehat{n} \le \poly(n / \epsilon)$,
    \item
          $\Enc_{n, \epsilon}$ is computable in time $\poly(n / \epsilon)$, and
    \item
          given $y \in \binset^{\widehat{n}}$,
          one can find a list of size $\poly(1 / \epsilon)$
          that contains all the strings $x \in \binset^n$ such that $y$ and $\Enc_{n, \epsilon}(x)$ agree on at least a $(1 / 2 + \epsilon)$-fraction of coordinates.
  \end{itemize}
\end{lemma}



\begin{proof}
  [Proof of \cref{thm_DPConstruction}]
  We describe the pseudorandom generator function $$\DP_k \colon \binset^n \times \binset^d \to \binset^{d+k},$$
  which we call a \emph{$k$-wise direct product generator}.
  Let $\Enc_{n, \delta}$ denote the error-correcting code of \cref{lemma_ListDecodable},
  where $\delta := \epsilon / 2 k$,
  and let
  $\widehat{x} \colon \binset^{\ell} \to \binset$
  be
  the function specified by the truth table $\Enc_{n, \delta}(x) \in \binset^{2^{{\ell}}}$,
  where
  ${\ell} = O(\log (n / \epsilon))$.
  The $k$-wise direct product generator is defined as
  $$
    \DP_k(x, \bar{z}) := (z^1, \cdots, z^k, \widehat{x}(z^1), \cdots, \widehat{x}(z^k))
  $$
  for any $\bar{z} = (z^1, \cdots, z^k) \in \roundbrackets{\binset^{\ell}}^k$.
  The seed length of $\DP_k$ is $d := \ell k = O(k \log (n/\epsilon))$.

  We first present a reconstruction procedure $R_0$ with small success probability;
  then we will amplify the success probability to $3/4.$%
  \footnote{
    We amplify the success probability so that the proof of \cref{thm_HeuristicaIntro} is
    clearer,
    although
    a proof can be given without the amplification.
  }
  \begin{claim}
    \label{claim_weaksecurityofDP}
    There exist an advice function $A_0 \colon \binset^{n} \times \binset^{r_0} \to \binset^{k}$
    and
    a reconstruction procedure $R_0^D \colon \binset^{a_0} \times \binset^{r_0} \to \binset^n$
    such that
    $$
      \Pr_w [R_0^D(A_0(x, w), w) = x] \ge \frac{\epsilon}{2k} \cdot \frac{1}{L},
    $$
    where $r_0 = O(d)$ and $L = \poly(k/\epsilon)$ denotes the list size of \cref{lemma_ListDecodable}.
  \end{claim}

  \begin{claimproof}
    \cref{claim_weaksecurityofDP} can be proved by using a standard hybrid argument (as in \cite{NisanW94_jcss_journals,Vadhan12_fttcs_journals}).
    Assume that $D \colon \binset^{d+k} \to \binset$ satisfies
    $$
      \Pr_{\bar{z}} \squarebrackets{ D(z^1, \cdots, z^k, \widehat{x}(z^1), \cdots, \widehat{x}(z^k)) = 1 }
      -
      \Pr_{\bar{z}, b} \squarebrackets{ D(z^1, \cdots, z^k, b_1, \cdots, b_k) = 1 }
      \ge \epsilon.
    $$
    For any $i \in \ESet{0, \cdots, k}$,
    define the $i$th hybrid distribution $H_i$ as
    the distribution of
    $$(z^1, \cdots, z^k, \widehat{x}(z^1), \cdots, \widehat{x}(z^i), b_{i+1}, \cdots, b_k),$$
    where $\bar{z} = (z^1, \cdots, z^k) \usample \roundbrackets{ \binset^{\ell} }^k$
    and $b_{i+1}, \cdots, b_k \usample \binset$.
    By this definition, $H_0$ is identically distributed with the uniform distribution, and $H_k$ is an identical distribution with $\DP_k(x, \bar{z})$.
    Therefore,
    \[
      \Exp_{\substack{i \usample \NumSet{k} \\ \bar{z}, b}} \squarebrackets {
        D(H_i) - D(H_{i-1})
      }
      \ge \frac{\epsilon}{k}.
    \]

    By an averaging argument,
    we obtain
    \begin{equation}
      \label{eq_HybridAdvantage}
      \Pr_{\substack{i \usample \NumSet{k}, b \\ z^1, \cdots, z^{i-1}, z^{i+1}, \cdots, z^k}} \squarebrackets {
        \Exp_{z^i \usample \binset^{\ell}} \squarebrackets{
          D(H_i) - D(H_{i-1})
        }
        \ge \frac{\epsilon}{2 k}
      }
      \ge \frac{\epsilon}{2 k}.
    \end{equation}
    %
    %
    Define the advice function $A_0$ as
    $A_0(x, w) := (\widehat{x}(z^1), \cdots, \widehat{x}(z^{i-1}), b_i, \cdots, b_k) \in \binset^k$,
    where $w$ is a coin flip sequence that contains
    the random choice of
    $i \usample \NumSet{k}$, $z^1, \cdots, z^{i-1} \usample \binset^{\ell}$, and $b \usample \binset^k$.
    By a standard calculation (see, \eg, \cite[Proposition 7.16]{Vadhan12_fttcs_journals}),
    it follows from \cref{eq_HybridAdvantage} that
    \begin{equation}
      \label{eq_HybridApproximateNextBit}
      \Pr_{z^i \usample \binset^\ell} \squarebrackets{ D(\bar{z}, A_0(x, w)) \oplus 1 \oplus b_i = \widehat{x}(z^i) } \ge \frac{1}{2} + \frac{\epsilon}{2 k }
    \end{equation}
    with probability at least $\epsilon / 2 k$ over the random choice of
    $(i, z^{\NumSet{k} \setminus \ESet{i}}, b)$.



    Now we describe the reconstruction procedure $R_0^D(\alpha, w)$.
    Given an advice string $\alpha \in \binset^k$ and a coin flip $w$,
    we regard the random bits $w$
    as $i \usample \NumSet{k}$, $z^1, \cdots, z^{i-1}, z^{i+1}, \cdots, z^k \usample \binset^{\ell}$, $b \usample \binset^k$, and $j \usample \NumSet{L}$.
    Define a string $y \in \binset^{2^\ell}$ as
    $$
      y_{z^i} :=
      D(\bar{z}, \alpha) \oplus 1 \oplus b_i
    $$
    for every $z^i \in \binset^\ell$.
    The reconstruction procedure uses the list-decoding algorithm of \cref{lemma_ListDecodable}
    to obtain a list of all the $2^\ell$-bit strings that agree with $y$ on a $(1/2 + \epsilon / 2k)$-fraction of indices,
    and outputs the $j$th string in the list.

    We claim that
    $\Pr_w [R_0^D(A_0(x, w), w) = x] \ge \epsilon / 2 k L$.
    By \cref{eq_HybridApproximateNextBit}, with probability at least $\epsilon / 2 k$ over the random choice of $(i, z^{\NumSet{k} \setminus \ESet{i}}, b)$,
    the list-decoding algorithm outputs a list of size $L$ that contains $x$.
    With probability at least $1 / L$ over the choice of $j \usample \NumSet{L}$,
    it holds that
    $x$ is the $j$th string in the list,
    in which case $R_0^D(A_0(x, w), w) = x$.
    Therefore, we obtain
    $$\Pr_w [R_0^D(A_0(x, w), w) = x] \ge \frac{\epsilon}{2 k} \cdot \frac{1}{L}.$$
  \end{claimproof}

  It remains to amplify the success probability.
  Intuitively,
  we repeat picking coin flip sequences $t = O(k L / \epsilon)$ times,
  and provide the successful index $i \in \NumSet{t}$ as an advice string.
  Details follow.

  The advice function $A^D \colon \binset^n \times (\binset^{r_0})^t \to \binset^{k + O(\log t)}$
  takes an $n$-bit string $x$ and coin flip sequences $(w_1, \cdots, w_t) \in (\binset^{r_0})^t$
  and outputs
  $(A_i(x, w_i), i) \in \binset^k \times \NumSet{t}$,
  where $i$ is the first index $i \in \NumSet{t}$ such that $R_0^D(A_0(x, w_i), w_i) = x$;
  if there is no such $i$, the output of $A^D$ is arbitrary.
  Given an advice string $(\alpha, i) \in \binset^k \times \NumSet{t}$
  and coin flip sequences $\bar{w} = (w_1, \cdots, w_t) \in (\binset^{r_0})^t$,
  we define $R^D((\alpha, i), \bar{w}) = R_0^D(\alpha, w_i)$.
  The failure probability is
  $$
    \Pr_{\bar{w}} \squarebrackets{
      R^D(A^D(x, \bar{w}), \bar{w}) \neq x
    }
    =
    \Pr_{\bar{w}} \squarebrackets{
      \forall i \in \NumSet{t},
      R_0^D(A_0(x, w_i), w_i) \neq x
    }
    \le (1 - \epsilon / 2 k L)^t \le 1 / 4,
  $$
  where the last inequality holds for a sufficiently large $t =O(k L / \epsilon)$.
\end{proof}

%
%
%
%
%
%
%
%
%
%

One of important ingredients of the proof of \cref{thm_HeuristicaIntro} is
a pseudorandom generator constructed by
Buhrman, Fortnow, and Pavan~\cite{BuhrmanFP05_mst_journals}.
\begin{lemma}
  [Buhrman, Fortnow, and Pavan~\cite{BuhrmanFP05_mst_journals}]
  \label{lemma_BFP05}
  If $\Dist\NP \subset \Avg\P$,
  then there exist a constant $c$ and a pseudorandom generator
  $G = \ESet{G_n : \binset^{c \log n} \to \binset^n }_{n \in \Nat}$
  that $(1/n)$-fools size-$n$ circuits.
\end{lemma}

The direct product generator construction (\cref{thm_DPConstruction})
achieves small advice complexity $a$; however,
a disadvantage is that
the randomness complexity $r$ is large.
The pseudorandom generator of \cref{lemma_BFP05} enables us to reduce the randomness complexity effectively.

Another ingredient is the fact that
$\Dist\NP^A \subset \Avg\P$
implies that
a dense subset of $A$-oracle time-bounded Kolmogorov-random strings
can be rejected
in polynomial time.
%

\begin{lemma}
  [\cite{Hirahara18_focs_conf}]
  \label{fact_DistNPtoHSG}
  Assume that $\Dist\NP^A \subset \Avg\P$.
  Then,
  there exists a polynomial-time algorithm $M$ such that
  \begin{enumerate}
    \item
          $M(x, \unary{t}) = 1$ for every $x$ such that $\KtTIME[A]{t}(x) < |x| - 1$, and
    \item
          $\Pr_{x \usample \binset^n} \squarebrackets{ M(x, \unary{t}) = 0} \ge \frac{1}{4}$
          for every $n \in \Nat$ and every $t \in \Nat$.
  \end{enumerate}
  %
  %
  %
  %
  %
  %
\end{lemma}
\begin{proof}
  [Proof Sketch]
  We consider a distributional problem $(L, \mathcal{D})$ defined as follows.
  The language $L \in \NP^A$ is the problem of deciding whether $\KtTIME[A]{t}(x) < |x| - 1$
  on input $(x, \unary{t})$.
  A distribution $\mathcal{D}_m$ picks $t \usample \NumSet{m}$ and $x \usample \binset^{m - t}$
  and outputs $(x, \unary{t})$.
  Since
  the fraction of $\Yes$ instances in $L$ is at most $\frac{1}{2}$ for any fixed $t$,
  any errorless heuristic algorithm that solves $(L, \mathcal{D})$
  must reject a $\frac{1}{4}$-fraction of inputs.
  Details can be found in \cite{Hirahara18_focs_conf}.
\end{proof}

Now we are ready to prove \cref{thm_HeuristicaIntro}.
\begin{proof}
  [Proof of \cref{thm_HeuristicaIntro}]
  Under the assumption that $\Dist\NP^A \subset \Avg\P$, we have the secure pseudorandom generator
  $G = \ESet{G_m : \binset^{c_0 \log m} \to \binset^m }_{m \in \Nat}$
  of \cref{lemma_BFP05}.
  In particular,
  by replacing random bits used by randomized algorithms
  with the output of $G$,
  we obtain $\pr\BPP = \pr\P$;
  thus
  it suffices to present a randomized polynomial-time algorithm $M_1$
  for computing $\KvsK[\tau]{A}$ for some polynomial $\tau$.


  Fix any input $(x, \unary{s}, \unary{t})$, where $x \in \binset^*$ is a string of length $n \in \Nat$
  and $s, t \in \Nat$.
  Take the $k$-wise direct product generator $\DP_k(x, \blank) \colon \binset^d \to \binset^{d + k}$ of \cref{thm_DPConstruction},
  where $k$ is some parameter chosen later and $\epsilon := \frac{1}{8}$.
  Let $M_0$ be the polynomial-time algorithm $M$ of \cref{fact_DistNPtoHSG}.
  Let $\tau_0$ be some polynomial chosen later.


  The randomized algorithm $M_1$ operates as follows.
  On input $(x, \unary{s}, \unary{t})$,
  $M_1$ samples a string $z \usample \binset^d$ uniformly at random.
  Then,
  $M_1$ simulates $M_0$ on input $(\DP_k(x, z), \unary{t'})$, where $t' := \tau_0(n, t)$,
  and accepts if and only if $M_0$ accepts.
  (That is, $M_1(x, \unary{s}, \unary{t})$ is defined to be $M_0(\DP_k(x, z), \unary{t'})$ for a random $z \usample \binset^d$.)

  We claim the correctness of the algorithm $M_1$ below.
  \begin{claim}
    \label{claim_CorrectnessOfHeuristica}
    \hspace{0em}
    \begin{enumerate}
      \item
            If $\KtTIME[A]{t}(x) \le s$,
            then $M_1(x, \unary{s}, \unary{t})$ accepts with probability $1$.
      \item
            If $\KtTIME{\tau(n, t)}(x) > s + \log \tau(n, t)$,
            then $M_1(x, \unary{s}, \unary{t})$ rejects with probability at least $\frac{1}{8}$.
    \end{enumerate}
  \end{claim}
  %
  %

  We claim the first item.
  Fix any ${z} \in \binset^d$.
  Since the output $\DP_k(x, {z})$ of
  the direct product generator
  can be described by
  $n, k \in \Nat$,
  the seed ${z} \in \binset^d$
  and
  the program for describing $x$ of size $\KtTIME[A]{t}(x)$
  in time $t' = \tau_0(n, t)$,
  where $\tau_0$ is some polynomial,
  it holds that
  \begin{align}
    \label{eq_YESBoundInHeuristica}
    \KtTIME[A]{t'}(\DP_k(x, {z})) \le d + \KtTIME[A]{t}(x) + c_1 \log n,
  \end{align}
  for some constant $c_1$.
  We set $k := s + c_1 \log n + 2$.
  Note that under the assumption that $\KtTIME[A]{t}(x) \le s$, \cref{eq_YESBoundInHeuristica} is less than $d + k - 1$.
  Therefore, by \cref{fact_DistNPtoHSG},
  $M_0(\DP_k(x, {z})) = 1$ for every ${z} \in \binset^d$;
  thus $M_1$ accepts.

  Next, we claim the second item, by proving its contrapositive.
  Assume that $M_1(x, \unary{s}, \unary{t})$ rejects with probability less than $\frac{1}{8}$.
  This means that
  $$
    \Pr_{{z} \usample \binset^d} \squarebrackets{ M_0(\DP_k(x, {z}), \unary{t}) = 0 } < \frac{1}{8}.
  $$
  On the other hand, by \cref{fact_DistNPtoHSG}, we also have
  $$
    \Pr_{w \usample \binset^{d+k}} \squarebrackets{ M_0(w, \unary{t}) = 0 } \ge \frac{1}{4}.
  $$
  Therefore,
  $$
    \Pr_{w \usample \binset^{d+k}} \squarebrackets{ M_0(w, \unary{t}) = 0 }
    -
    \Pr_{{z} \usample \binset^d} \squarebrackets{ M_0(\DP_k(x, {z}), \unary{t}) = 0 }
    \ge \frac{1}{8},
  $$
  which means that a function $D$ defined as $D(w) := \neg M_0(w, \unary{t})$ distinguishes $\DP_k(x, \blank)$ from the uniform distribution.
  By the reconstruction property of \cref{thm_DPConstruction},
  \begin{align}
    \label{eq_ReconstructionInHeuristica}
    \Pr_{\rho \usample \binset^r} \squarebrackets{ R^D(A^D(x, \rho), \rho) = x } \ge \frac{3}{4},
  \end{align}
  where the advice complexity $a$ is at most $k + O(\log (k / \epsilon)) = s + O(\log (n / \epsilon))$.%
  \footnote{We may assume without loss of generality that $k := s + O(\log n) \le 2 n$,
    as otherwise $\KtTIME[A]{t}(x) \le s$ always holds. }
  Now we derandomize the random choice of $\rho$ of \cref{eq_ReconstructionInHeuristica}
  by using the secure pseudorandom generator $G = \ESet{G_m}_{m \in \Nat}$.
  That is,
  we argue that $\rho$ can be replaced with $G_m(\rho_0)$ for some short seed $\rho_0$,
  which enables us to obtain a short description for $x$.
  To this end, consider a statistical test $T \colon \binset^r \to \binset$ that checks the condition of \cref{eq_ReconstructionInHeuristica}; that is,
  $T(\rho) = 1$ if and only if $R^D(A^D(x, \rho), \rho) = x$.
  One can easily observe that $T$ can be implemented by a circuit of size $m := \poly(n, t)$.

  %
  %
  %
  %
  %
  %
  %
  %
  %
  %
  %
  %

  Now we replace the random bits $\rho \in \binset^r$ with the first $r$ bits of the pseudorandom sequence $G_m(\rho_0)$.
  By \cref{eq_ReconstructionInHeuristica},
  we have
  $$
    \Pr_{\rho_0 \usample \binset^{c_0 \log m}} \squarebrackets{ T(G(\rho_0)) = 1 } \ge
    \Pr_{\rho \usample \binset^r} \squarebrackets{ T(\rho) = 1 } - o(1) > 0.
  $$
  In particular, there exists a seed $\rho_0 \in \binset^{c_0 \log m}$
  such that $T(G_m(\rho_0)) = 1$.

  We are ready to present the algorithm for describing $x$.
  In order to describe $x$, it takes as a description
  $n, t, m \in \Nat$,
  the seed $\rho_0 \in \binset^{c_0 \log m}$,
  and
  the advice string $\alpha := A^D(G_m(\rho_0)) \in \binset^a$.
  Since $T(G_m(\rho_0)) = 1$,
  the string $x$ is equal to $R^D(\alpha, G_m(\rho_0))$,
  which can be computed in time $\tau_1(n, t)$ for some polynomial $\tau_1$.
  Therefore,
  $$
    \KtTIME{\tau_1(n, t)}(x) \le a + c_0 \log m + O(\log n t) \le s + O(\log n t).
  $$
  In particular, by choosing a polynomial $\tau$ large enough, we have
  $$\KtTIME{\tau(n, t)}(x) \le \KtTIME{\tau_1(n, t)}(x) \le s + \log \tau(n, t).$$
  This completes the proof of \cref{claim_CorrectnessOfHeuristica}.

  Since $M_1$ is a one-sided-error algorithm, the success probability can be amplified by repeating the computation of $M_1$
  for independent random coin flips.
  We thus conclude that $\KvsK[\tau]{A}$ is in $\pr\BPP = \pr\P$.
\end{proof}

An important corollary is that $\NP$-hardness of $\KvsK{A}$ for some $A \in \PH$
is sufficient for proving an equivalence between worst-case and average-case complexity of $\PH$.
\begin{restatetheorem}
  {cor_TowardsPH}
  Assume that $\KvsK{A}$ is ``$\NP$-hard'' for some $A \in \PH$
  in the sense that
  $$
    \NP \not\subset \BPP
    \quad
    \Longrightarrow
    \quad
    \KvsK{A} \not\in \P.
  $$
  Then,
  $$
    \Dist\PH \not\subset \Avg\P
    \quad
    \Longleftrightarrow
    \quad
    \PH \neq \P.
  $$
\end{restatetheorem}
\begin{proof}
  It is obvious that $\PH = \P$ implies that $\Dist\PH \subset \Avg\P$.
  We prove the converse direction.
  Assume that $\Dist\PH \subset \Avg\P$.
  By \cref{thm_HeuristicaIntro},
  we obtain
  $\KvsK{A} \in \P$
  for any $A \in \PH$.
  It follows from the contrapositive of the assumption that
  $\NP \subset \BPP$;
  moreover, by \cref{lemma_BFP05}, we also have $\BPP = \P$.
  Therefore, we obtain $\NP = \P$, which is equivalent to $\PH = \P$.
\end{proof}



Another corollary is
a relationship between time-bounded Kolmogorov complexity
and its $\PH$-oracle version,
which can be proved by using the disjointness of $\KvsK{A}$.
\begin{corollary}
  \label{cor:oracle elimination}
  If $\Dist\PH \subset \Avg\P$,
  then,
  for any $A \in \PH$,
  there exists a polynomial $\tau$ such that
  $$\KtTIME{\tau(|x|, t)}(x) \le \KtTIME[A]{t}(x) + \log \tau(|x|, t)$$
  for any $x \in \binset^*$ and $t \in \Nat$.
\end{corollary}
\begin{proof}
  By \cref{thm_HeuristicaIntro},
  under the assumption that $\Dist\NP^A \subset \Dist\PH \subset \Avg\P$,
  there exists an algorithm $M$ such that
  $M(x, \unary{s}, \unary{t}) = 1$ for every $x$ such that $\KtTIME[A]{t}(x) \le s$
  and
  $M(x, \unary{s}, \unary{t}) = 0$ for every $x$ such that $\KtTIME{\tau(|x|, t)}(x) > s + \log \tau(|x|, t)$,
  for any $s \in \Nat$ and $t \in \Nat$.
  For any $x \in \binset^*$ and $t \in \Nat$, define $s := \KtTIME[A]{t}(x)$;
  since $0 \neq M(x, \unary{s}, \unary{t}) = 1$, we obtain
  $\KtTIME{\tau(|x|, t)}(x) \le s + \log \tau(|x|, t)$.
\end{proof}


%
Under the plausible assumption that $\E^\NP \neq \E$,
we observe that $\KvsK{\SAT}$ is non-disjoint.
\begin{proposition}
  \label{prop_NondisjointKSATvsK}
  If $\E^\NP \neq \E$,
  then,
  for some polynomial $\tau_0$ and any polynomial $\tau$,
  there are infinitely many strings $x$
  such that
  $\KtTIME[\SAT]{t}(x) = O(\log |x|)$
  and
  $\KtTIME{\tau(|x|)}(x) > \log \tau(|x|)$,
  where $t := \tau_0(|x|)$.
\end{proposition}
\begin{proof}
  By \cite{BuhrmanH92_fsttcs_conf}, the assumption is equivalent to $\E^\NP \not\subset \E/O(n)$.
  Take a language $L \in \E^\NP \setminus \E/O(n)$.
  For each $n \in \Nat$,
  define $x_n$ to be the truth table of length $2^n$ that encodes the characteristic function of $L \cap \binset^n$.
  Since $x_n$ can be described in time $\tau_0(|x_n|) = \poly(|x_n|)$ given $n \in \Nat$ and oracle access to $\SAT$,
  we have $\KtTIME[\SAT]{t}(x_n) = O(\log |x_n|)$.
  On the other hand, if $\KtTIME{\tau(|x_n|)}(x_n) \le \log \tau(|x_n|)$ for all large $n \in \Nat$,
  there exists an advice string of length $\log \tau(|x_n|) = O(n)$ that
  makes it possible
  to compute $L \cap \binset^n$ in time $\poly(\tau(|x_n|)) = 2^{O(n)}$,
  which contradicts $L \not\in \E/O(n)$.
\end{proof}


Finally, we observe that the complexity of $\KvsK{\SAT}$
is closely related to the complexity of $\MINKT$.
\begin{proposition}
  For any polynomial $\tau$, the following hold.
  \begin{itemize}
    \item
          $\GapVS[\tau]{\K}{\K}$ is reducible to $\GapVS[\tau]{\K^\SAT}{\K}$ via an identity map.
    \item
          If $\GapVS[\tau]{\K^\SAT}{\K}$ is disjoint,
          then
          $\GapVS[\tau]{\K^\SAT}{\K}$ is reducible to
          $\MINKT$; in particular,
          $\GapVS[\tau]{\K^\SAT}{\K} \in \NP$.
  \end{itemize}
\end{proposition}

\begin{proof}
  The first item is obvious because $\KtTIME[\SAT]{t}(x) \le \KtTIME{t}(x)$ for any $t \in \Nat$ and $x \in \binset^*$.
  For the second item,
  let $(\YES, \NO)$ denote
  $\GapVS[\tau]{\K^\SAT}{\K}$.
  Since $(\overline{\NO}, \NO)$ is a problem of checking whether $\KtTIME{\tau(|x|, t)}(x) \le s + \log \tau(|x|, t)$
  given $(x, \unary{s}, \unary{t})$ as input,
  it is reducible to $\MINKT$.
  In particular, $(\YES, \NO) \in \NP$ holds under the assumption that it is disjoint.
\end{proof}


\subsection{\texorpdfstring{$\GapVS{F}{F^{-1}}$}{Gap(F vs F(-1))}: PRG construction from one-way functions}
\label{section_HILL}
H{\aa}stad, Impagliazzo, Levin, and Luby~\cite{HastadILL99_siamcomp_journals}
showed that
a cryptographic pseudorandom generator can be constructed from any one-way function.
In this section, we view the black-box PRG construction from a meta-computational perspective,
which leads us to
the promise problem
$\GapVS{F}{F^{-1}}$,
which is a problem of asking whether a given function $f$ is
computable by a small circuit
or
$f$ is hard to invert by any small circuit.
Here we assume that the function $f \colon \binset^n \to \binset^n$ is given as oracle,
and we focus on an algorithm that runs in time $\poly(n)$.
In other words, we consider a sublinear-time algorithm that is given random access to the truth table of $f$.
%

\begin{definition}
  For a function $\tau \colon \Nat \times \Nat \to \Nat$,
  $\GapVS[\tau]{F}{F^{-1}}$
  is a promise problem $(\YES, \NO)$ defined as follows.
  The input consists of a size parameter $s \in \Nat$, an integer $n \in \Nat$,
  and black-box access to a function $f \colon \binset^n \to \binset^n$.
  \begin{itemize}
    \item
          The set $\YES$ consists of inputs $(n, s, f)$
          such that $\size(f) \le s$.
    \item
          The set $\NO$ consists of
          inputs $(n, s, f)$
          such that, for any oracle circuit $C$ of size $\tau(n, s)$,
          $$
            \Pr_{x \usample \binset^n} \squarebrackets{ C^f(f(x)) \in f^{-1}(f(x)) } < \frac{1}{2}.
          $$
  \end{itemize}
\end{definition}

\begin{theorem}
  \label{thm_FvsFinv}
  If $\Dist\NP \subset \Avg\P$,
  then there exist
  a polynomial $\tau$ and
  a $\co\RP$-type randomized algorithm $M$ that solves $\GapVS[\tau]{F}{F^{-1}} = (\YES, \NO)$
  on input $(n, s)$
  in time $\poly(n, s)$.
  That is, $M$ is a randomized oracle algorithm such that
  \begin{enumerate}
    \item
          $\Pr_M [ M^f(n, s) = 1 ] = 1$ for every $(n, s, f) \in \YES$,
    \item
          $\Pr_M [ M^f(n, s) = 0 ] \ge \frac{1}{2}$ for every $(n, s, f) \in \NO$, and
    \item
          $M^f(n, s)$ runs in time $\poly(n, s)$.
  \end{enumerate}
\end{theorem}

For the proof, we make use of the following black-box construction
of a pseudorandom generator based on any one-way function.
\begin{lemma}
  [A black-box PRG Construction from Any OWF \cite{HastadILL99_siamcomp_journals}]
  \label{lemma_BBPRGfromOWF}
  There exists a polynomial $d = d(n)$
  such that,
  for a parameter $m \in \Nat$,
  there exists a polynomial-time oracle algorithm $G_{m}^{(\blank)} \colon \binset^{d(n)} \to \binset^{m}$
  that takes a function $f \colon \binset^n \to \binset^n$
  and there exists an oracle algorithm $R$
  such that, for any function $D \colon \binset^m \to \binset$, if $m > d$ and
  $$
    \Pr_{z \usample \binset^n} \squarebrackets{ D(G^f(z)) = 1 }
    -
    \Pr_{w \usample \binset^m} \squarebrackets{ D(w) = 1 }
    \ge
    \frac{1}{8},
  $$
  then
  $$
    \Pr_{x, R} [R^{f, D}(f(x)) \in f^{-1}(f(x))] \ge \frac{1}{2}.
  $$
  The running time of $G_{m}^{(\blank)}$ and $R$ is at most $\poly(n, m)$.
\end{lemma}

\begin{proof}
  [Proof Sketch]
  Since a weakly one-way function exists if and only if a strongly one-way function exists (\cite{Yao82a_focs_conf}),
  it suffices to present a reduction that inverts $f$ with probability at least $1  / \poly(n, m)$.
  (To be more specific,
  we first amplify the hardness of $f$ by taking a direct product
  $f^t(x_1, \cdots, x_t) := (f(x_1), \cdots, f(x_t))$,
  where $t$ is some appropriately chosen parameter, and
  then apply the HILL construction to $f^t$ described below.)

  We invoke the pseudorandom generator construction $G_\mathrm{HILL}^{f}$ of
  H{\aa}stad, Impagliazzo, Levin, and Luby~\cite{HastadILL99_siamcomp_journals}
  based on $f$.
  They presented a security reduction $R$
  such
  that
  if there exists a function $D$ that distinguishes $G_\mathrm{HILL}^{f}$ from
  the uniform distribution,
  then an oracle algorithm $R^{f, D}(f({x}))$ can compute an element in $f^{-1}(f({x}))$
  with probability at least $1 / \poly(n, m)$ over the choice of $x \usample \binset^{n}$ and the
  internal randomness of $R$.
\end{proof}

\begin{proof}
  [Proof of \cref{thm_FvsFinv}]  %
%
 \quad Let\ $G^{(\blank)}$\ be\ the\ black-box\ pseudorandom\ generator\
 construction\ of\ \cref{lemma_BBPRGfromOWF},
  and let $R$ be the security reduction of \cref{lemma_BBPRGfromOWF}.

  Under the assumption that $\Dist\NP \subset \Avg\P$,
  by \cref{fact_DistNPtoHSG},
  there exists a polynomial-time algorithm $M$ such that
  $M(x, \unary{t}) = 1$ for every $x$ such that $\KtTIME{t}(x) < |x| - 1$, and
  %
  \[ \Pr_{x \usample \binset^n} \squarebrackets{ M(x, \unary{t}) = 0}
    \ge \frac{1}{4} \]
  for every $n \in \Nat$ and every $t \in \Nat$.

  The algorithm $M'$ for computing $\GapVS{F}{F^{-1}}$ is defined as follows.
  Let $f \colon \binset^n \to \binset^n$ and $s \in \Nat$ be inputs.
  Define $m := p(n, s)$ for some polynomial $p$ chosen later.
  Pick a random $z \in \binset^n$.
  Then
  $M'$ accepts if and only if
  $M(G_{m}^{f}(z), \unary{t})$ accepts for a sufficiently large $t = \poly(n, s)$.

  We claim the correctness of $M'$ below.
  \begin{claim}
    \hspace{0em}
    \begin{enumerate}
      \item
            $M'$ accepts
            any $f$ such that $\size(f) \le s$
            with probability $1$.
      \item
            $M'$ rejects
            any $f$ such that $\Pr_x \squarebrackets{ C^f(f(x)) \in f^{-1}(f(x)) } < \frac{1}{2}$
            with probability at least $\frac{1}{8}$.
    \end{enumerate}
  \end{claim}

  We claim that any $f$ such that $\size(f) \le s$ is accepted by the algorithm $M'$.
  Indeed,
  since $f$ is computable by some circuit of size $s$ and $G_{m}^f$ is polynomial-time-computable,
  the output of the generator $G_{m}^f(z)$ can be described
  by using
  the description of the circuit of size $s$, the seed $z$ of length $d(n)$, and $n, m \in \Nat$ in polynomial time;
  thus, for a sufficiently large $t = \poly(n, s)$,
  $$
    \KtTIME{t}(G_{m}^f(z)) \le d(n) + \widetilde{O}(s) + O(\log n).
  $$
  Choosing $m = p(n, s)$ large enough, this is bounded by $m - 2$.
  Thus $M'$ accepts.

  Conversely, suppose that
  the algorithm $M'$ accepts with probability at least $\frac{7}{8}$.
  This means that
  $$
    \Pr_{z \usample \binset^n} \squarebrackets{ M(G^f_{m}(z), \unary{t}) = 1 } \ge \frac{7}{8}.
  $$
  On the other hand, by \cref{fact_DistNPtoHSG}, we have
  $$
    \Pr_{w \usample \binset^m} \squarebrackets{ M(w, \unary{t}) = 1 } \le \frac{3}{4}.
  $$
  Therefore, $M( \blank, \unary{t})$ is a distinguisher for $G_{ m}^f$.
  It follows from the property of the security reduction $R$
  that
  $$
    \Pr_{x, R} \squarebrackets{ R^{f, M(\blank, \unary{t})}(f(x)) \in f^{-1}(f(x)) } \ge \frac{1}{2}.
  $$
  By fixing the internal randomness of $R$ and
  simulating the polynomial-time algorithm $R^{f, M(\blank, \unary{t})}$ by a polynomial-size circuit $C^f$,
  we conclude that $f$ can be inverted by the oracle circuit $C^f$ of size $\poly(n, s) =: \tau(n, s)$.
  Thus $f$ is not a $\No$ instance of $\GapVS[\tau]{F}{F^{-1}}$.
\end{proof}

\begin{remark}
  In fact, the assumption that $\Dist\NP \subset \Avg\P$ of \cref{thm_FvsFinv} can be weakened to
  the assumption that there exists a $\P$-natural property useful against $\SIZE(2^{o(n)})$ (which is essentially equivalent to
  an errorless heuristic algorithm for $\MCSP$ \cite{HiraharaS17_coco_conf,Hirahara18_focs_conf}).
  Indeed,
  as in \cite{AllenderBKMR06_siamcomp_journals,RazborovR97_jcss_journals,AllenderH19_toct_journals},
  the pseudorandom function generator construction of \cite{GoldreichGM86_jacm_journals}
  can be used to construct a black-box pseudorandom generator $G^f$ based on a one-way function $f$
  that satisfies
  $\size(G^f(z)) \le \poly(|z|, \log |G^f(z)|)$ for any seed $z \in \binset^d$;
  such a pseudorandom generator $G^f$ can be distinguished from the uniform distribution by using the natural property.
\end{remark}

We now explain that $\GapVS{F}{F^{-1}}$ is conjectured to be non-disjoint.
An \emph{auxiliary-input one-way function} (AIOWF) $f = \ESet{ f_x : \binset^{p(|x|)} \to \binset^{q(|x|)} }_{x \in \binset^*}$
is
a polynomial-time-computable function such that,
for some infinite set $I$,
for any non-uniform polynomial-time algorithm $A$,
$\Pr_y \squarebrackets{ A(x, f_x(y)) \in f_x^{-1}(f_x(y)) } < 1 / n^{\omega(1)}$ for all large $n \in \Nat$ and any $x \in I$.
This is
a weaker cryptographic primitive than a one-way function (\ie, the existence of a one-way function implies that of an auxiliary-input one-way function).
Ostrovsky~\cite{Ostrovsky91_coco_conf}
showed that non-triviality of $\SZK$ implies the existence of an auxiliary-input one-way function.
(see also \cite{Vadhan06_siamcomp_journals}).
We observe that the existence of an auxiliary-input one-way function implies the non-disjointness of $\GapVS{F}{F^{-1}}$.

\begin{proposition}
  \label{prop_AIOWF}
  If there exists an auxiliary-input one-way function
  $f = \ESet{ f_x : \binset^{p(|x|)} \to \binset^{q(|x|)} }_{x \in \binset^*}$,
  then,
  for any polynomial $\tau$,
  $\GapVS[\tau]{F}{F^{-1}}$ is non-disjoint.
\end{proposition}
\begin{proof}
  Take an infinite set $I \subset \binset^*$ that is hard for polynomial-size circuits to invert $\ESet{f_x}_{x \in I}$.
  Since $f$ is polynomial-time-computable, $\size(f_x) \le n^c$ for some constant $c$,
  where $n = |x|$.
  We set the size parameter $s := n^c$.
  On the other hand,
  by the property of AIOWF,
  for any circuit $A$ of size $\tau(n, s)$,
  it holds that
  $$\Pr_y \squarebrackets{ A(x, f_x(y)) \in f_x^{-1}(f_x(y)) } < \frac{1}{2},$$
  for a sufficiently large $x \in I$.
  This means that $f_x$ is a \Yes and \No instance of $\GapVS[\tau]{F}{F^{-1}}$.
\end{proof}

An immediate corollary of
\cref{thm_FvsFinv}
and
\cref{prop_AIOWF}
is that the existence of AIOWF implies $\Dist\NP \not\subset \Avg\P$
(which is already shown in \cite{Hirahara18_focs_conf}).
An interesting open question is to prove ``$\NP$-hardness'' of $\GapVS{F}{F^{-1}}$,
which has the following important consequence:
\begin{corollary}
  If,
  for any polynomial $\tau$,
  it is \textup{``}$\NP$-hard\textup{''} to solve $\GapVS[\tau]{F}{F^{-1}}$ in time $\poly(n, s)$,
  then
  the worst-case and average-case complexity of $\NP$ is equivalent
  in the sense that $\P \neq \NP$ iff $\Dist\NP \not\subset \Avg\P$.
\end{corollary}

\begin{proof}
  The assumption that
  $\GapVS{F}{F^{-1}}$
  is ``$\NP$-hard''
  means that, for any polynomial $\tau$,
  if there exists a $\co\RP$-type algorithm
  that solves $\GapVS[\tau]{F}{F^{-1}}$
  on input $(n, s)$ in time $\poly(n, s)$,
  then
  $\NP \subset \BPP$.
  If $\Dist\NP \subset \Avg\P$,
  then
  \cref{thm_FvsFinv}
  implies that
  there exists some polynomial $\tau$ such that
  $\GapVS[\tau]{F}{F^{-1}}$ can be solved by a $\co\RP$-type algorithm
  in time $\poly(n, s)$.
  By the assumption, we obtain $\NP \subset \BPP = \P$, where the last equality is from \cref{lemma_BFP05}.
\end{proof}


\section{Meta-computational view of complexity-theoretic PRG constructions}
\label{section_Genral}

%
This section provides a meta-computational view of
\emph{complexity-theoretic} PRG constructions.
The section is organized as follows.
In \cref{subsection_UniversalHSG},
we give definitions of resource-bounded Kolmogorov complexity and universal hitting set generators.
In \cref{section_MCLPdef},
we interpret the notion of advice
using resource-bounded Kolmogorov complexity.
Using the (new) notion of advice,
we provide the definition of MCLPs in \cref{subsection_DefiningMCLPs}.
In \cref{subsection_MCLPsFromPRG}, we present
a generic connection from black-box PRG constructions to MCLPs.
\cref{section_GeneralNW} applies the connection
to the specific PRG construction given by Nisan and Wigderson~\cite{NisanW94_jcss_journals},
which is reviewed in \cref{subsection_ReviewNWGenerator}.
\cref{section_KtandHSG} provides
a meta-computational view of hardness amplification theorems,
which gives rise to reductions among different MCLPs.
In \cref{section_KSandROBP},
we apply our principle to space-bounded algorithms.

%
%
%
%
%
%
%
%
%
%

\subsection{Universal hitting set generators and Kolmogorov complexity}
\label{subsection_UniversalHSG}


We review the notion of $\Kt$-complexity and its space-bounded version, and present
definitions of universal hitting set generators.
Recall the definition of Levin's resource-bounded Kolmogorov complexity.
\begin{definition}
  [\cite{Levin84_iandc_journals}]
  For a string $x \in \binset^*$, Levin's $\Kt$ complexity of $x$ is defined as
  $$
    \Kt(x) := \min \ISet { |d| + t } { \text{$U$ outputs $x$ on input $d$ in time $2^t$} },
  $$
  where $U$ is an efficient universal Turing machine.
\end{definition}

It is possible to enumerate all the strings $x$ such that $\Kt(x) \le s$
in time $2^{O(s)}$.
The enumeration algorithm gives rise to a universal time-bounded hitting set generator.
\begin{definition}
  [Universal time-bounded hitting set generator]
  For a function $s \colon \Nat \to \Nat$,
  define
  a \emph{universal time-bounded hitting set generator}
  $\UnivHt^s = \ESet{ \UnivHt^s_n : \binset^{s(n)+1} \to \binset^n }_{n \in \Nat}$ so that
  $\UnivHt^s_n(d 0 1^t)$ is equal to the output of $U$ on input $d$
  if $U(d) \in \binset^n$ and $U(d)$ halts in $2^t$ time steps,
  where $d \in \binset^*$ and $t \in \Nat$.
\end{definition}

The hitting set generator $\UnivHt$ is universal in the sense that, if there exists a secure hitting set generator $G_n \colon \binset^{s} \to \binset^n$ computable in time $2^s$,
then
$\UnivHt^{O(s)}_n \colon \binset^{O(s)} \to \binset^n$ is also secure.
Indeed, the image of $G_n$ is a subset of that of $\UnivHt^{O(s)}_n$ because
of the following property.
\begin{proposition}
  [Universality of $\UnivHt$]
  \label{prop_UniversalHSGforTime}
  For every function $s \colon \Nat \to \Nat$ and $n \in \Nat$,
  the image of $\UnivHt^s_n$ contains every string $x \in \binset^n$ such that $\Kt(x) \le s(n)$.
  Moreover, the image of $\UnivHt^s_n$ can be enumerated in time $\widetilde O(2^{s(n) + \log n})$.
\end{proposition}
\begin{proof}
  Consider any string $x$ of length $n$ such that $\Kt(x) \le s(n)$.
  By the definition of $\Kt$-complexity,
  there exists a description $d \in \binset^*$ such that
  $U(d)$ outputs $x$ in time $2^t$ for some $t \le s(n)$ and $|d| \le s(n) - t$.
  By the definition of $\UnivHt^s_n$, we have
  $\UnivHt^s_n(d 0 1^t) = x$.
\end{proof}


The notion of $\KS$-complexity was introduced in \cite{AllenderBKMR06_siamcomp_journals}
as a space-bounded analogue of $\Kt$-complexity.
\begin{definition}
  [\cite{AllenderBKMR06_siamcomp_journals}]
  For a string $x \in \binset^*$, the $\KS$-complexity of $x$ is defined as
  $$
    \KS(x) := \min \ISet { |d| + s } { \text{$U$ outputs $x$ on input $d$ in space $s$} }.
  $$
\end{definition}
One can modify the definition of the universal time-bounded hitting set generator
to a space-bounded version.
\begin{definition}
  [Universal space-bounded hitting set generator]
  For a function $s \colon \Nat \to \Nat$,
  define
  a \emph{universal space-bounded hitting set generator}
  $\UnivGS^s = \ESet{ \UnivGS^s_n : \binset^{s(n)} \to \binset^n }_{n \in \Nat}$
  so that
  $\UnivGS^s_n(d 0 1^t)$ is equal to
  the output of $U$ on input $d$ if
  $U(d) \in \binset^n$ and $U(d)$ uses at most $t$ space,
  where $d \in \binset^*$ and $t \in \Nat$.
\end{definition}
\begin{proposition}
  [Universality of $\UnivGS$]
  \label{prop_UniversalHSGforSpace}
  For every function $s \colon \Nat \to \Nat$ and $n \in \Nat$,
  the image of $\UnivGS^s_n$ contains every string $x \in \binset^n$ such that $\KS(x) \le s(n)$.
  Moreover, $\UnivGS^s_n$ can be computed in $O(s(n) + \log n)$ space.
\end{proposition}

\subsection{Length-wise advice and resource-bounded Kolmogorov complexity}
\label{section_MCLPdef}

In order to define meta-computational circuit lower-bound problems,
we modify the standard notion of advice.
Usually, a complexity class with advice such as $\E / O(n)$ is defined as a subset of functions $f \colon \binset^* \to \binset$ that are defined on all the strings of any length.
Here,
for any $n \in \Nat$,
we define
``$\DTIME(2^{O(n)}) \Nadvice O(n)$'' as a subset of functions $f \colon \binset^n \to \binset$,
where the superscript $n$ in ``$\Nadvice$'' is appended in order to emphasize that it depends on $n$.
\begin{definition}
  [Length-wise advice]
  \label{def_AdviceLengthwise}
  For any integers $t, a, n \in \Nat$,
  we denote by $\DTIME(t) \Nadvice a$ the class of functions $f \colon \binset^n \to \binset$
  such that
  there exists a Turing machine $M$ whose description length%
  \footnote{%
    The \emph{description length} $|d_M|$ of a machine $M$ can be formally defined by using a universal Turing machine $U$.  Specifically, a description $d_M \in \binset^*$ of $M$
    is a string that satisfies $U(d_M, x) = M(x)$ for every $x$.}
  is $a$
  and
  that outputs $f(x)$ on input $x \in \binset^n$ in time $t$.
  Similarly, let $\DSPACE(t) \Nadvice a$ denote the class of functions $f \colon \binset^n \to \binset$
  such that
  there exists a Turing machine $M$ whose description length
  is $a$
  and
  that outputs $f(x)$ on input $x \in \binset^n$ in \emph{space} $t$.
\end{definition}

%
%
%
%
%
%
%
%
%
%
%


%
This definition is slightly different from the standard notion of complexity classes with advice,
but these are essentially the same.
In order to clarify the difference,
for functions $t, a \colon \Nat \to \Nat$,
let $\DTIME(t) \KLadvice a$ denote
the complexity class $\DTIME(t)$ with $a$-bit advice strings
in the standard sense of Karp and Lipton~\cite{KarpL82_advice}.%
\footnote{
  It is also common to use the notation $\DTIME(t(n)) \KLadvice a(n)$,
  where $n$ is an indeterminate.
  %
  %
}
That is, a function $f \colon \binset^* \to \binset$ is in $\DTIME(t) \KLadvice a$
if and only if
there exists a Turing machine $M$ such that, for any $n \in \Nat$,
there exists an advice string $\alpha_n \in \binset^{a(n)}$
such that
$M$ outputs $f(x)$ on input $(x, \alpha_n)$
in time $t(n)$
for every $x \in \binset^n$.
Then, the length-wise advice ``$\Nadvice$'' and the Karp--Lipton advice ``$\KLadvice$'' are
equivalent in the following sense.
\begin{fact}
  For any functions $t, a \colon \Nat \to \Nat$
  and any family of functions $f = \ESet{ f_n \colon \binset^n \to \binset }_{n \in \Nat}$
  \textup{(}which is identified with a function $f \colon \binset^* \to \binset$\textup{)},
  the following are equivalent.
  \begin{enumerate}
    \item
          There exists a constant $c$ such that $f \in \DTIME(t') \KLadvice a'$,
          where
          $t'(n) := t(n)^c + c$ and $a'(n) := c \cdot a(n) + c$.
    \item
          There exists a constant $c$ such that, for any $n \in \Nat$,
          $f_n \in \DTIME(t(n)^c + c) \Nadvice c \cdot a(n) + c$.
  \end{enumerate}
\end{fact}
\begin{proof}
  [Proof Sketch]
  If $f \in \DTIME(t') \KLadvice a'$,
  then there exists a machine $M$ that takes an advice string $\alpha_n$ on inputs of length $n$.
  For each $n \in \Nat$, define $M_n$ to be the machine that, on input $x$, simulates $M$ on input $(x, \alpha_n)$; the description length of $M_n$ is at most $O(|\alpha_n|)$, and thus $f_n \in \DTIME(t(n)^{O(1)}) \Nadvice O(a(n))$.
  Conversely, if, for any $n \in \Nat$, there exists a Turing machine $M_n$ whose description length is $O(a(n))$, then
  a universal Turing machine $U$ witnesses $f \in \DTIME(t') \KLadvice a'$.
\end{proof}

The advice ``$\Nadvice$'' is equivalent to $\Kt$-complexity
up to a constant factor
in the following sense.
\begin{fact}
  \label{fact_KtandAdvice}
  For any function $t \colon \Nat \to \Nat$
  such that $t(n) \ge n$
  and for any family of functions $f = \ESet{ f_n \colon \binset^n \to \binset }_{n \in \Nat}$,
  the following are equivalent.
  \begin{enumerate}
    \item
          $f_n \in \DTIME(t(n)^{O(1)}) \Nadvice O(\log t(n))$ for all large $n \in \Nat$.
    \item
          $\Kt(f_n) = O(\log t(n))$ for all large $n \in \Nat$.
  \end{enumerate}
\end{fact}

\subsection{Meta-computational circuit lower-bound problems (MCLPs)}
\label{subsection_DefiningMCLPs}
We define promise problems of
distinguishing
the truth table of explicit functions (\eg, computable in $\DTIME(2^{c n}) \Nadvice c n$)
from
the truth table of hard functions (\eg, that cannot be computed in $\SIZE(2^{\epsilon n})$).
We call these
Meta-computational Circuit Lower-bound Problems (MCLPs).

\begin{definition}
  [Meta-computational Circuit Lower-bound Problems; MCLPs]
  Let $\mathcal{E}, \mathcal{D}$ be families of functions.
  \emph{The $\vsProblem{\mathcal{E}}{\mathcal{D}}$ Problem} is defined as the following promise problem $(\YES, \NO)$.
  \begin{align*}
    \YES & := \ISet{ \truthtable(f) }{ f \in \mathcal{E} },
    \\
    \NO  & := \ISet{ \truthtable(f) } { f \not\in \mathcal{D} }.
  \end{align*}
  We will mainly consider a non-uniform computational model for computing the $\vsProblem{\mathcal{E}}{\mathcal{D}}$ Problem;
  for $N \in \Nat$,
  we denote
  by $(\vsProblem{\mathcal{E}}{\mathcal{D}})_N$
  the problem restricted to the input length of $N$.

  We denote by $\roundbrackets{\vsProblem{\E}{\mathcal{D}}}$ a family of problems $\ESet { \vsProblem{\E_c }{\mathcal{D}} }_{c \in \Nat}$,
  where $\E_c := \bigcup_{n \in \Nat}  \DTIME(2^{c n}) \Nadvice c n$.
  For a circuit class $\mathfrak{C}$,
  we say that $(\vsProblem{\E}{\mathcal{D}})$ is solved by a $\mathfrak{C}$-circuit of size $s(N)$ and denote by
  $\roundbrackets{\vsProblem{\E}{\mathcal{D}}} \in \io\mathfrak{C}(s(N))$
  if, for every constant $c$, there exists a family of $\mathfrak{C}$-circuits $\ESet{ C_N }_{N \in \Nat}$ of size $s(N)$
  such that $C_N$ solves the promise problem $(\vsProblem{\E_c}{\mathcal{D}})_N$ for infinitely many $N$.
  We also denote by $\roundbrackets{\vsProblem{\E}{\SIZE(2^{o(n)})}}$
  a family of problems $\ESet{ \vsProblem{\E_c}{\SIZE(2^{\alpha n})} }_{c \in \Nat, \alpha > 0}$.
  %
  %

  Similarly, $\roundbrackets{\vsProblem{\DSPACE(n)}{\SIZE(2^{o(n)})}}$ denotes $\ESet { \vsProblem{\bigcup_{n \in \Nat}\DSPACE(c n) \Nadvice c n}{\SIZE(2^{\alpha n})} }_{c \in \Nat, \alpha > 0}.$
\end{definition}

The definition of the $\vsProblem{\E}{\SIZE(2^{o(n)})}$ Problem given here is slightly different from \cref{def_EvsSIZE} given earlier.
In \cref{def_EvsSIZE}, we defined the $\vsProblem{\E}{\SIZE(2^{o(n)})}$ Problem by using the notion of $\Kt$-complexity;
however,
these definitions are essentially equivalent in light of \cref{fact_KtandAdvice}.

We justify the notation of $\roundbrackets{\vsProblem{\DSPACE(n)}{\SIZE(2^{o(n)})}}$ below.
One can observe that
the open question %
%
whether $\DSPACE(n) \not\subset \io{\SIZE}(2^{o(n)})$
is closely related to the ${\vsProblem{\DSPACE(n)}{\SIZE(2^{o(n)})}}$ Problem.
%
%

\begin{proposition}
  \label{prop_NonproperKS}
  The following are equivalent.
  \begin{enumerate}
    \item
          \label{item_DSPACElowerbound}
          $\DSPACE(n) \not\subset \io{\SIZE}(2^{\epsilon n})$ for some constant  $\epsilon> 0$.
    \item
          \label{item_NondisjointDSPACESIZE}
          No circuit can solve
          the $\vsProblem{\DSPACE(n)}{{\SIZE}(2^{o(n)})}$ Problem for all large $n \in \Nat$.
    \item
          \label{item_NondisjointDSPACEApproxSIZE}
          No circuit can solve
          the $\vsProblem{\DSPACE(n)}{\widetilde{\SIZE}(2^{o(n)}; \frac{1}{2} - 2^{-o(n)})}$ Problem for all large $n \in \Nat$.
    \item
          \label{item_NondisjointDSPACEApproxSIZERestated}
          There exist some constants $c \in \Nat, \alpha > 0$ such that,
          for all large $n \in \Nat$,
          there exists a function $f \colon \binset^n \to \binset$ such that
          $f \in \dspace{c n}$
          and $f \not\in \widetilde{\SIZE}(2^{\alpha n}; \frac{1}{2} - 2^{-\alpha n})$.
  \end{enumerate}
  %
  %
  %
  %
  %
  %
  %
  %
  %
  %
  %
  %
\end{proposition}
\begin{proof}
  First, observe that \cref{item_NondisjointDSPACEApproxSIZE} is equivalent to \cref{item_NondisjointDSPACEApproxSIZERestated}.
  Indeed, \cref{item_NondisjointDSPACEApproxSIZERestated} implies \cref{item_NondisjointDSPACEApproxSIZE} by the definition.
  Conversely, if \cref{item_NondisjointDSPACEApproxSIZERestated} does not hold,
  the $\vsProblem{\DSPACE(n)}{\widetilde{\SIZE}(2^{o(n)}; \frac{1}{2} - 2^{-o(n)})}$ Problem is disjoint; thus, there exists an exponential size circuit that computes the problem,
  which means that \ref{item_NondisjointDSPACEApproxSIZERestated} does not hold.

  Next, we claim that \cref{item_DSPACElowerbound} implies \cref{item_NondisjointDSPACEApproxSIZERestated}.
  By using a locally-decodable error-correcting code, %
  it can be shown that \cref{item_DSPACElowerbound} is equivalent to $\DSPACE(n) \not\subset \io\widetilde{\SIZE}(2^{\alpha n}; \frac{1}{2} - 2^{- \alpha n})$
  for some constant $\alpha > 0$
  %
  (see  %
  \cite{SudanTV01_jcss_journals, KlivansM02_siamcomp_journals}).
  Take a problem $L \in \DSPACE(n) \setminus \io\widetilde{\SIZE}(2^{\alpha n}; \frac{1}{2} - 2^{- \alpha n})$.
  For each $n \in \Nat$,
  let $f_n \colon \binset^n \to \binset$ be the characteristic function of $L \cap \binset^n$.
  Then,
  there exists a constant $c$ such that, for all large $n \in \Nat$, $f_n \in \dspace{c n}$
  and
  $f_n \not\in \widetilde{\SIZE}(2^{\alpha n}; {\frac{1}{2} - 2^{- \alpha n}})$,
  which completes the proof of \cref{item_NondisjointDSPACEApproxSIZERestated}.
  It is immediate from the definition that \cref{item_NondisjointDSPACEApproxSIZERestated} implies \cref{item_NondisjointDSPACESIZE}.

  We claim that \cref{item_NondisjointDSPACESIZE} implies \cref{item_DSPACElowerbound}.
  Let $f = \ESet{ f_n : \binset^n \to \binset }_{n \in \Nat}$ be a family of functions
  such that
  $f_n \in \dspace{c n}$ and $f_n \not\in \SIZE(2^{\alpha n})$ for all large $n \in \Nat$.
  In particular, for all large $n \in \Nat$, there exists some machine $M_n$ of description length $c n$
  that computes $f_n(x)$ on input $x \in \binset^n$ in space $c n$.
  Let $L \subset \binset^*$ be a language such that $(x, M) \in L$ if and only if
  the description length of a Turing machine $M$ is $c |x|$ and
  $M$ accepts $x$ in space $c |x|$.
  It is clear that $L \in \DSPACE(n)$.
  Moreover, for all large $n \in \Nat$, the characteristic function of $\ISet{ x \in \binset^n }{ (x, M_n) \in L }$
  is equal to $f_n$; thus, it requires circuits of size $2^{\alpha n}$.
  Therefore, $L \not\in \io\SIZE(2^{\alpha n / (c + 1)})$.
\end{proof}


It is also possible to define
MCLPs whose non-disjointness characterizes other circuit lower bounds.
For example,
the $\vsProblem{\EXP/\poly}{\SIZE(2^{o(n)})}$ Problem defined as $$\left\{ \vsProblem{\bigcup_{n \in \Nat}\DTIME(2^{n^c}) \Nadvice n^c}{\SIZE(2^{\alpha n})} \right\}_{c \in \Nat, \alpha > 0}$$
is non-disjoint if and only if $\EXP/\poly \not\subset \SIZE(2^{o(n)})$.



\subsection{MCLPs from PRG and HSG constructions}
\label{subsection_MCLPsFromPRG}

We formalize the notion of black-box PRG and HSG construction
and then we present a general connection
between
black-box PRG and HSG constructions
and
meta-computational circuit lower bound problems.

\begin{definition}
  [Black-Box PRG and HSG Construction]
  \label{def_BBConstruction}
  Let $G^{(\blank)} \colon \binset^d \to \binset^m$ be an oracle algorithm
  that expects an oracle of the form $f \colon \binset^\ell \to \binset$.
  Let $\mathfrak{C}$ be a circuit class and $\mathcal{R}^{(\blank)}$ be an oracle circuit class.
  The algorithm $G$ is referred to as
  a \emph{black-box $\mathfrak{C}$-pseudorandom generator
    (resp.\ $\mathfrak{C}$-hitting set generator)
    construction with $\mathcal{R}$-reconstruction} and error parameter $\epsilon$
  if the following hold.
  \begin{description}
    \item[$\mathfrak{C}$-Construction]
          For any seed $z \in \binset^d$,
          there exists a $\mathfrak{C}$-circuit
          that takes $\truthtable(f)$ as input
          and outputs $G^f(z)$.
    \item[$\mathcal{R}$-Reconstruction]
          For any function $f \colon \binset^\ell \to \binset$ and
          any function $D \colon \binset^m \to \binset$,
          if $D$ is an $\epsilon$-distinguisher for $G^f$
          (resp.\ $D$ $\epsilon$-avoids $G^f$),
          then $f \in \mathcal{R}^D$.
  \end{description}
\end{definition}

%
We now present a generic connection between
MCLPs
and HSG constructions.
\begin{theorem}
  %
  \label{thm_Generic}
  Let $G^{(\blank)} \colon \binset^s \to \binset^m$
  be a black-box $\mathfrak{C}$-construction with $\mathcal{R}$-reconstruction
  that takes a function $f \colon \binset^n \to \binset$.
  Define $N := 2^n$.
  Let $H \colon \binset^{d} \to \binset^m$ be a function such that $G^f(z) \in \image(H)$
  for every $f \in \mathcal{E}$ and every $z \in \binset^{d}$.
  Let $D \colon \binset^m \to \binset$ be any function that $\epsilon$-avoids $H$.
  Then, the following hold.
  \begin{enumerate}
    \item
          If
          $G^f$ is a HSG construction with error parameter $\epsilon$,
          then
          $(\vsProblem{\mathcal{E}}{\mathcal{R}^D})_N \in \AND_{2^s} \circ \NOT \circ D \circ \mathfrak{C}$.
    \item
          If
          $G^f$ is a PRG construction with error parameter $\epsilon / 2$,
          then
          $(\vsProblem{\mathcal{E}}{\mathcal{R}^D})_N \in \AND_{O(N/\epsilon)} \circ \NOT \circ D \circ \mathfrak{C}$.
  \end{enumerate}
\end{theorem}

\begin{proof}
  Let $G^{(\blank)}$ be a HSG construction with error parameter $\epsilon$.
  We first present a randomized circuit for solving $(\vsProblem{\mathcal{E}}{\mathcal{R}^D})$
  on inputs of length $N$.
  Let $f \colon \binset^n \to \binset$ denote an input.
  Consider a circuit $D_1$ such that $D_1(f; z) := D\roundbrackets{ G^f(z) }$,
  where $z$ is an auxiliary input (that will be regarded as non-deterministic bits or random bits).
  We claim that $D_1$ can solve $(\vsProblem{\mathcal{E}}{\mathcal{D}})$ co-nondeterministically.
  \begin{claim}
    \label{claim_Generic}
    \hspace{0em}
    \begin{enumerate}
      \item
            If $f \in \mathcal{E}$, then $D_1(f; z) = 0$ for every $z \in \binset^s$.
      \item
            If $f \not\in \mathcal{R}^D$, then $D_1(f; z) = 1$ for some $z \in \binset^s$.
    \end{enumerate}
  \end{claim}
  %
  Suppose that $f$ is a \Yes instance of $(\vsProblem{\mathcal{E}}{\mathcal{D}})$, that is, $f \in \mathcal{E}$.
  By the assumption, we have $G^f(z) \in \image(H)$.
  Therefore, by the property of $D$, we obtain $D_1(f; z) = D\roundbrackets{ G^f(z) } = 0$.

  Conversely, suppose that $D\roundbrackets{ G^f(z) } = 0$ for every $z \in \binset^s$.
  This means that $D$ $\epsilon$-avoids $G^f$.
  By the reconstruction property of $G^f$,
  we obtain $f \in \mathcal{R}^D$.
  Taking its contrapositive, it follows that
  if $f \not\in \mathcal{R}^D$ then $D_1(f; z) = D \roundbrackets{ G^f(z) } = 1$ for some $z \in \binset^s$.
  This completes the proof of \cref{claim_Generic}.
  %

  Now consider a circuit $D_2$ defined as $D_2(f) := \bigwedge_{z \in \binset^s} \neg D_1(f; z)$.
  Then, it follows from \cref{claim_Generic} that the circuit $D_2$ solves $(\vsProblem{\mathcal{E}}{\mathcal{R}^D})$
  on inputs of length $N$.

  We move on to the case when $G^{(\blank)}$ is a PRG construction.
  In this case, we claim that the second item of \cref{claim_Generic} can be strengthened to the following:
  If $f \not\in \mathcal{R}^D$, then $\Pr_{z \usample \binset^s} \squarebrackets{ D_1(f; z) = 1 } \ge \frac{\epsilon}{2}$.

  %
  We prove the contrapositive of this claim.
  Assume that $\Pr_{z \usample \binset^s} \squarebrackets{ D_1(f; z) = 1 } < \frac{\epsilon}{2} .$
  Since $D$ $\epsilon$-avoids $H$, we have $\Pr_{w \usample \binset^m} \squarebrackets{ D(w) = 1 } \ge \epsilon$.
  Therefore, $D$ $\frac{\epsilon}{2}$-distinguishes $G^f(\blank)$ from the uniform distribution;
  by the reconstruction of $G^{(\blank)}$, we obtain $f \in \mathcal{R}^D$, as desired.

  Therefore, $D_1(\blank; z)$ is a one-sided-error randomized circuit that computes $(\vsProblem{\mathcal{E}}{\mathcal{R}^D})$
  with probability at least $\epsilon / 2$.
  Now define a randomized circuit $D_2'$ such that $D_2'(f) := \bigwedge_{i = 1}^k \neg D_1(f; z_i)$,
  where $z_1, \cdots, z_k \sim \binset^s$ are chosen independently and $k$ is a parameter chosen later.
  Then, it is easy to see that $D_2'(f) = 1$ for any $f \in \mathcal{E}$; on the other hand, for any $f \not\in \mathcal{R}^D$,
  $D_2'(f) = 1$ with probability at most $(1 - \epsilon / 2)^k$, which is less than $2^{-N}$ by choosing $k = O(N / \epsilon)$ large enough.
  By using a union bound, one can hardwire random bits in $D_2'$ as in Adleman's trick \cite{Adleman78_focs_conf},
  and obtain a deterministic $\AND \circ \NOT \circ D \circ \mathfrak{C}$ circuit that computes $(\vsProblem{\mathcal{E}}{\mathcal{R}^D})$.
\end{proof}




\subsection{The Nisan--Wigderson generator}
\label{subsection_ReviewNWGenerator}
The pseudorandom generator construction of Nisan and Wigderson~\cite{NisanW94_jcss_journals} is particularly efficient.
Indeed, each output bit of the Nisan--Wigderson generator
depends on only $1$ bit of the truth table of a candidate hard function.

\begin{theorem}
  [Nisan--Wigderson Pseudorandom Generator Construction]
  \label{thm_NWGenerator}
  %
  For every constant $\gamma > 0$ and any $\ell, m \in \Nat$,
  there exists a $\mathfrak{C}$-pseudorandom generator construction $G^{(\blank)} \colon \binset^d \to \binset^m$
  with $\mathcal{R}$-reconstruction and error parameter $\epsilon$
  that takes a function $f \colon \binset^\ell \to \binset$
  such that
  \begin{itemize}
    \item
          $d = O(\ell)$, $m \le 2^\ell$,
    \item
          $\mathfrak{C} = \NC^0_1$, and
    \item
          $\mathcal{R}^D = \widetilde{D \circ \AC^0_2}(m \cdot 2^{\gamma \ell}; \frac{1}{2} - \frac{\epsilon}{m})$.
  \end{itemize}
  Moreover, the output $G^{f}(z)$ of the PRG
  can be computed in space $O(\ell)$
  given a function $f$ and a seed $z$ as input.
\end{theorem}

We first recall the construction of the Nisan--Wigderson generator.
In order to have
a space-efficient algorithm for computing the Nisan--Wigderson generator,
we use the following construction of a combinatorial design.
\begin{lemma}
  [Klivans and van Melkebeek \cite{KlivansM02_siamcomp_journals}, Viola \cite{Viola05_cc_journals}]
  \label{lemma_design}
  For every constant $\gamma > 0$, for any $\ell \in \Nat$,
  there exist $\ell$-sized subsets $S_1, \cdots, S_{2^\ell}$ of $\NumSet{d}$ for some $d = O(\ell)$
  such that
  \begin{enumerate}
    \item
          $|S_i \cap S_j| \le \gamma \cdot \ell$ for every distinct $i, j \in \NumSet{2^\ell}$, and
    \item
          $S_1, \cdots, S_{2^\ell}$ can be constructed in space $O(\ell)$.
  \end{enumerate}
\end{lemma}

A nearly disjoint generator $\ND$ is defined based on the design.
\begin{definition}
  \label{def_NDgenerator}
  For a string $z \in \binset^d$ and a subset $S \subset \NumSet{d}$ of indices,
  $z_S$ denotes the string obtained by concatenating the $i$th bit $z_i$ of $z$ for every $i \in S$.
  For $\ell$-sized subsets $\mathcal{S} = \ESet{S_1, \cdots, S_{m}}$ of $\NumSet{d}$, define
  a \emph{nearly disjoint generator} $\ND \colon \binset^d \to (\binset^\ell)^m$ as
  $\ND(z) := (z_{S_1}, \cdots, z_{S_m})$ for every $z \in \binset^d$.
\end{definition}
%

Using the nearly disjoint generator, the Nisan--Wigderson generator is defined as follows.
\begin{definition}
  [Nisan--Wigderson generator \cite{NisanW94_jcss_journals}]
  \label{def_NWgenerator}
  For a Boolean function $f \colon \binset^\ell \to \binset$ and
  for $\ell$-sized subsets $\mathcal{S} = \ESet{S_1, \cdots, S_{m}}$ of $\NumSet{d}$,
  the Nisan--Wigderson generator $\NW^f \colon \binset^d \to \binset^m$ is defined as
  $$\NW^f(z) := f^{m} \circ \ND (z) = f(z_{S_1}) \cdots f(z_{S_m}).$$
\end{definition}

It was shown in \cite{NisanW94_jcss_journals}
that the pseudorandom generator construction is secure in the following sense.

\begin{lemma}
  [Security of the Nisan--Wigderson Generator \cite{NisanW94_jcss_journals}]
  \label{lemma_SecurityOfNW}
  For any functions $T \colon \binset^m \to \binset$ and $f \colon \binset^\ell \to \binset$, for any $\epsilon > 0$,
  if
  $$
    \Pr_{w \usample \binset^m} \left[ T(w) = 1 \right] - \Pr_{z \usample \binset^d} \left[ T(\NW^f(z)) = 1 \right] \ge \epsilon,
  $$
  then
  there exists a one-query depth-2 oracle circuit $C^T \in T \circ \AC^0_2$ of size $O(m \cdot 2^{\gamma \ell})$
  such that
  $$
    \Pr_{x \usample \binset^\ell} \left[ C^T(x) = f(x) \right] \ge \frac{1}{2} + \frac{\epsilon}{m}.
  $$
\end{lemma}

\begin{proof}
  [Proof of \cref{thm_NWGenerator}]
  We use the Nisan--Wigderson generator $\NW^{(\blank)}$ defined in \cref{def_NWgenerator} (\ie, we define $G^{(\blank)} := \NW^{(\blank)}$).
  $\NW^{(\blank)}$ is an $\NC^0_1$-PRG construction because each bit of $\NW^{f}(z)$ for any fixed seed $z$
  is equal to one bit of the truth table of $f$.
  The $\mathcal{R}$-reconstruction property of $\NW^{(\blank)}$ follows from \cref{lemma_SecurityOfNW}.
\end{proof}


\subsection{Meta-computational view of the Nisan--Wigderson generator}
\label{section_GeneralNW}

Using the Nisan--Wigderson generator construction,
we present a general connection between the existence of hitting set generators
and lower bounds for meta-computational circuit lower-bound problems.
We consider any family of circuits that is closed under taking projections in the following sense.
\begin{definition}
  A class $\mathfrak{C}$ of circuits is said to be
  \emph{closed under taking projections}
  if, for any $s \in \Nat$,
  for every size-$s$ circuit $C \in \mathfrak{C}$ of $n$ inputs,
  a circuit $C'$ defined as $C'(x_1, \cdots, x_n) = C(x_{\sigma(1)}, \cdots, x_{\sigma(n)})$ for some function $\sigma \colon \NumSet{n} \to \NumSet{n}$
  can be simulated by a size-$s$ $\mathfrak{C}$-circuit.
\end{definition}

\begin{theorem}
  \label{thm_GeneralNW}
  Let $\mathfrak{C}$ be any circuit class that is closed under taking projections.
  Suppose that
  $(\vsProblem{\E}{\widetilde{\mathfrak{C} \circ \AC^0_2} (2^{\alpha n} ; \frac{1}{2} - 2^{- \alpha n})} ) \not\in \io \AND \circ \NOT \circ \mathfrak{C} (N^{1 + \beta})$
  for some constants $\alpha, \beta > 0$.
  Then,
  there exists a hitting set generator $G = \ESet{ G_n : \binset^{O(\log n)} \to \binset^n }_{n \in \Nat} $ computable in time $n^{O(1)}$ and secure against linear-size $\mathfrak{C}$ circuits.
\end{theorem}

\begin{proof}
  We prove the contrapositive.
  Assume that,
  for every function $s(m) = O(\log m)$,
  there exists a linear-size $\mathfrak{C}$ circuit $D$ that avoids the universal hitting set generator $\UnivHt_m^s$
  for infinitely many $m \in \Nat$.
  Given arbitrary constants $c, \alpha, \beta > 0$,
  we will choose a small constant $\gamma > 0$,
  and define $s(m) := c' \log m / \gamma$ for some large constant $c'$.
  Then using $D$ that avoids $\UnivHt_m^s$,
  we present a $\AND \circ \NOT \circ \mathfrak{C}$-circuit of size $N^{1 + \beta}$ that solves the
  $\vsProblem{\E_c}{\widetilde{\mathfrak{C} \circ \AC^0_2} (2^{\alpha n} ; \frac{1}{2} - 2^{- \alpha n})}$ Problem.
  %

  Let $f \colon \binset^n \to \binset$ denote the input of the MCLP.
  We use the Nisan--Wigderson generator construction $\NW^{f} \colon \binset^{d} \to \binset^{m}$ of \cref{thm_NWGenerator},
  where $d = O(n)$, $m = 2^{\gamma n}$, and $\epsilon := \frac{1}{4}$.
  By \cref{thm_NWGenerator}, this is a black-box $\NC^0_1$-PRG construction with $\widetilde{\AC^0_2}(2^{2 \gamma n}; \frac{1}{2} - \frac{1}{4 m})$ reconstruction.

  In order to apply \cref{thm_Generic},
  we claim that $\NW^f(z) \in \image(\UnivHt_m^s)$
  for every $f \in \E_c$ and every $z \in \binset^d$.
  By \cref{thm_NWGenerator},
  the output $\NW^f(z)$ can be described by
  using a seed $z$ and a description for $f$ in time $N^{O(1)}$,
  and thus
  \begin{align*}
    \Kt(\NW^{f}(z)) \le |z| + \Kt(f) + O(\log N) = O(n).
  \end{align*}
  In particular, for a large enough constant $c'$, we have
  \begin{align*}
    \Kt(\NW^{f}(z)) \le c' n = s(m).
  \end{align*}
  By the universality of $\UnivHt_m^s$ (\cref{prop_UniversalHSGforTime})
  we obtain that $\NW^f(z) \in \image(\UnivHt_m^s)$.

  By applying \cref{thm_Generic},
  we have $(\vsProblem{\E_c}{\widetilde{D \circ \AC^0_2}}(2^{2 \gamma n}; \frac{1}{2} - \frac{1}{4 m}))_N \in \AND_{O(N)} \circ \NOT \circ D \circ \NC^0_1$.
  Since $D$ is a $\mathfrak{C}$-circuit of size $m = 2^{\gamma n}$,
  it follows that
  $(\vsProblem{\E_c}{\widetilde{\mathfrak{C} \circ \AC^0_2}(2^{O(\gamma n)}; \frac{1}{2} - \frac{1}{4 m})})_N \in \AND \circ \NOT \circ \mathfrak{C} (O(N^{1 + \gamma}))$.
  Note that this holds for infinitely many $N = m^{1 / \gamma}$.
  By choosing $\gamma$ small enough depending on $\alpha, \beta$, the result follows.
\end{proof}



We observe that the converse direction also holds.
In particular, for any circuit class $\mathfrak{C} \subset \Ppoly$ such that $\mathfrak{C}$ is closed under taking projections and $\AND \circ \NOT \circ \mathfrak{C} = \mathfrak{C}$,
our reductions in fact establish the equivalence
between the existence of a hitting set generator secure against $\mathfrak{C}$
and a $\mathfrak{C}$-lower bound for the $\vsProblem{\E}{\widetilde{\SIZE}(2^{o(n)}; \frac{1}{2} - 2^{- o(n)})}$ Problem.
\begin{proposition}
  [Converse of \cref{thm_GeneralNW}]
  \label{prop_Converse}
  Let $\mathfrak{C}$ be any circuit complexity class.
  Suppose that
  there exists a hitting set generator $G = \ESet{ G_n : \binset^{O(\log n)} \to \binset^n }_{n \in \Nat} $ computable in time $n^{O(1)}$ and secure against linear-size $\mathfrak{C}$ circuits.
  Then,
  for any constant $k \in \Nat$,
  for all sufficiently large $N \in \Nat$,
  no $\NOT \circ \mathfrak{C}$ circuit of size $N^{k}$
  can solve $\roundbrackets{ \vsProblem{\E}{\widetilde{\SIZE}(2^{\alpha n}; \frac{1}{2} - 2^{- \alpha n})} }_N$
  for $\alpha := 1/4$
  nor $\MKtP[O(\log N), N - 1]$ on input length $N$.
\end{proposition}

\begin{proof}
  We first observe that
  $\MKtP[O(\log N), N - 1]$
  is reducible to
  $(\vsProblem{\E}{\widetilde{\SIZE}(2^{\alpha n}; \frac{1}{2} - 2^{- \alpha n})})$
  via an identity map:
  Take any function $f \in \widetilde{\SIZE}({2^{\alpha n}; \frac{1}{2} - 2^{- \alpha n}})$.
  Since the truth table of $f$ can be described by a circuit of size $N^\alpha$ and $\log \binom{N}{\le N/2 - N^{1 - \alpha}}$ bits of information
  in time $N^{O(1)}$,
  the $\Kt$-complexity of $\truthtable(f)$ is at most
  $$
    \widetilde{O}(N^\alpha) + N \cdot \mathrm{H}_2(1 / 2 - N^{-\alpha}) + O(\log N) \le N - \Omega(N^{1 - 2 \alpha}),
  $$
  which is much smaller than $N - 1$.
  Therefore, it suffices to prove the result only for $\MKtP[O(\log N), N - 1]$.

  We prove the contrapositive.
  Suppose that, for some constant $k \in \Nat$,
  for any constant $c$,
  for infinitely many $N \in \Nat$,
  there exists a $\NOT \circ \mathfrak{C}$ circuit of size $N^{k}$ that solves
  the promise problem $\MKtP[c \log N, N - 1]$.
  Consider any family of functions $G = \ESet{ G_m : \binset^{d \log m} \to \binset^m }_{m \in \Nat} $ computable in time $m^{d}$ for a constant $d$.
  Let $c := 4 k d$, and take a $\NOT \circ \mathfrak{C}$ circuit $\neg C$
  of size $N^k$ that solves $\MKtP[c \log N, N - 1]$ on inputs of length $N$.

  We regard $C$ as a circuit that takes $m := N^k$ input bits by ignoring $m - N$ input bits,
  and in what follows we claim that the linear-size circuit $C$ avoids $G_m$.
  For a string $w \in \binset^m$, denote by $w \restriction_N$ the first $N$ bits of $w$.

  Let $z \in \binset^{d \log m}$ be any seed of $G_m$.
  Since $G_m(z) \restriction_N$ can be described by $N \in \Nat$ and $z \in \binset^{d \log m}$ in time $m^d$,
  its $\Kt$ complexity is
  $$
    \Kt(G_m(z) \restriction_N) \le \log N + |z| + d \log m + o(\log m) \le 4 k d \log N,
  $$
  which means that $G_m(z) \restriction_N$ is a \Yes instance of $\MKtP[c \log N, N - 1]$ and thus $G_m(z)$ is rejected by $C$.

  Now consider a string $w \usample \binset^m$ chosen uniformly at random.
  By a standard counting argument, $\Kt(w \restriction_N) \ge N - 1$ with probability at least $\frac{1}{2}$;
  thus $C$ accepts at least a half of all inputs.
  Therefore, the function $G_m$ is not secure against $C$.
\end{proof}

Applying \cref{thm_GeneralNW} to depth-$d$ circuits $\mathfrak{C} := \AC^0_d$,
we obtain the following.
\begin{corollary}
  \label{cor_NWAC0}
  Let $d$ be a constant.
  Suppose that
  $(\vsProblem{\E}{\widetilde{\AC^{0}_{d+2}} (2^{\alpha n} ; \frac{1}{2} - 2^{- \alpha n})} ) \not\in \io \AC^0_{d+1} (N^{1 + \beta})$
  for some constants $\alpha, \beta > 0$.
  Then,
  there exists a hitting set generator $G = \ESet{ G_n : \binset^{O(\log n)} \to \binset^n }_{n \in \Nat} $ computable in time $n^{O(1)}$ and secure against linear-size $\AC^0_{d}$ circuits.
\end{corollary}
This means that, in order to obtain a nearly optimal hitting set generator for $\AC^0_d$,
it suffices to prove that nearly linear-size $\AC^0_{d+1}$ circuits cannot distinguish the truth tables of functions in $\E/O(n)$
from the truth tables of functions that cannot be approximated by $\AC^0_{d+2}$ circuits.

We present a proof of \cref{thm_ACEquiv}.
\begin{restatetheorem}
  {thm_ACEquiv}
  The following are equivalent.
  \begin{enumerate}
    \item
          \label{item_EvsACnotinACRestated}
          For any constants $d, d'$,
          there exists a constant $\beta > 0$ such that
          $(\vsProblem{\E}{\widetilde{\AC_{d'}^0} (2^{o( n)} ; \frac{1}{2} - 2^{- o(n)})}) \not\in \io\AC^0_d(N^{1 + \beta})$.
    \item
          \label{item_AC0HSGRestated}
          For any constant $d$,
          there exists a hitting set generator $G = \ESet{ G_n : \binset^{O(\log n)} \to \binset^n }_{n \in \Nat} $ computable in time $n^{O(1)}$ and secure against linear-size $\AC_d^0$ circuits.
    \item
          \label{item_MKtPnotinlinearACRestated}
          For any constant $d$, there exist constants $c, \beta > 0$ such that
          $\MKtP[c \log N, N - 1] \not\in \io \AC_d^0(N^{1 + \beta})$.
    \item
          \label{item_MKtPnotinACRestated}
          For any constants $d, k$, there exists a constant $c > 0$ such that
          $\MKtP[c \log N, N - 1] \not\in \io \AC_d^0(N^{k})$.
  \end{enumerate}
\end{restatetheorem}


\begin{proof}
  [Proof of \cref{thm_ACEquiv}]
  \cref{item_EvsACnotinACRestated} $\Longrightarrow$
  \cref{item_AC0HSGRestated}
  follows from \cref{cor_NWAC0}.
  \cref{item_AC0HSGRestated}
  $\Longrightarrow$
  \cref{item_MKtPnotinACRestated}
  follows from \cref{prop_Converse}.
  \cref{item_MKtPnotinACRestated}
  $\Longrightarrow$
  \cref{item_MKtPnotinlinearACRestated}
  is obvious.
  \cref{item_MKtPnotinlinearACRestated}
  $\Longrightarrow$
  \cref{item_EvsACnotinACRestated}
  holds because the truth table of any function in $\widetilde{\AC^0}(2^{o(n)}; \frac{1}{2} - 2^{- o(n)})$
  has $\Kt$ complexity less than $N - 1$.
\end{proof}

%
%
%



\subsection{Hardness amplification and MCLPs}
\label{section_KtandHSG}

The main advantage of studying MCLPs is that
hardness amplification can be naturally regarded as a reduction
between two different MCLPs.
In this section, we present such reductions.

Impagliazzo and Wigderson~\cite{ImpagliazzoW97_stoc_conf} gave a derandomized hardness amplification theorem.
We use the following generalized version of their result.
\begin{theorem}
  [Derandomized hardness amplification]
  \label{lemma_HardnessAmplification}
  Let $\gamma > 0$ be an arbitrary constant.
  There exists a hardness amplification procedure $\Amp$
  that takes a function $f \colon \binset^n \to \binset$ and parameters $\delta, \epsilon > 0$,
  and returns
  a Boolean function
  $\Amp_{\epsilon, \delta}^f \colon \binset^{O(n + \log (1/\epsilon))} \to \binset$
  satisfying the following:
  \begin{enumerate}
    \item
          \label{item_AmpApprox}
          If $\ApproxCC{1/2 - \epsilon}{\Amp_{\epsilon, \delta}^f} \le s$,
          then
          $\ApproxCC{\delta}{f} \le s \cdot \poly(1/\epsilon, 1/\delta)$.
    \item  \label{item_AmpXOR}
          For any fixed $y \in \binset^{O(n + \log (1 / \epsilon))}$,
          there exist strings $v_1, \cdots, v_k \in \binset^{n}$
          for some $k = O(\log (1 / \epsilon) / \delta)$
          such that
          $\Amp_{\epsilon, \delta}^f(y) =  f(v_1) \oplus \cdots \oplus f(v_k)$.
          Moreover, if $y$ is distributed uniformly at random,
          for each $i \in \NumSet{k}$,
          $v_i$ is distributed uniformly at random.
    \item \label{item_AmpKS}
          $\Amp_{\epsilon, \delta}^f(y)$ can be computed in $O(n + \log (1 / \delta) + \log (1 / \epsilon))$ space,
          given $f, y, \epsilon$ and $\delta$ as an input.
  \end{enumerate}
\end{theorem}
\cref{lemma_HardnessAmplification} slightly differs from derandomized hardness amplification theorems of \cite{ImpagliazzoW97_stoc_conf, HealyVV06_siamcomp_journals} in that
we are also interested in the case when the hardness parameter $\delta$ is $o(1)$,
which is not required in a standard application of derandomized hardness amplification theorems.
We defer a proof of \cref{lemma_HardnessAmplification} to \cref{section_HardnessAmplification}.
%

Using \cref{lemma_HardnessAmplification}, we give a reduction among
different MCLPs.
\begin{theorem}
  \label{thm_Amplification}
  For any constants $\alpha, \delta > 0$,
  for all sufficiently small $\beta > 0$,
  there exists
  a $\XOR_{O(\beta \log N)}$-computable reduction from
  $(\vsProblem{\E}{\widetilde{\SIZE}(2^{\alpha n}; \delta)})$
  to
  $(\vsProblem{\E}{\widetilde{\SIZE}(2^{\beta n}; \frac{1}{2} - 2^{-\beta n}}))$.
\end{theorem}

\begin{proof}
  The reduction is to simply take the hardness amplification procedure $\Amp$ of \cref{lemma_HardnessAmplification}.
  Specifically,
  given the truth table of a function $f \colon \binset^n \to \binset$, the reduction maps $f$ to the truth table of $\Amp^f_{\epsilon, \delta} \colon \binset^{n'} \to \binset$,
  where $\epsilon := 2^{- \beta n}$ and $n' = O(n)$.
  By the second item of \cref{lemma_HardnessAmplification},
  each output of the reduction is computable by a XOR of $k$ bits, where $k = O(\log (1 / \epsilon) / \delta ) = O(\beta n)$.

  We claim the correctness of the reduction.
  Suppose that $\Kt(f) \le O( \log N )$.
  Then, by the third item of \cref{lemma_HardnessAmplification},
  we obtain that $\Kt(\Amp_{\epsilon, \delta}^f) \le O(\log N)$.

  Conversely, suppose that $f \not\in \widetilde{\SIZE}(2^{\alpha n}; \delta)$;
  that is, $\ApproxCC{\delta}{f} > 2^{\alpha n} > 2^{\beta n} \cdot \poly(1 / \epsilon, 1 / \delta)$,
  where the last inequality holds by choosing $\beta > 0$ small enough.
  By the first item of \cref{lemma_HardnessAmplification},
  we obtain that $\ApproxCC{\frac{1}{2} - 2^{- \beta n}}{\Amp_{\epsilon, \delta}^f} > 2^{\beta n}$,
  which implies that $\Amp_{\epsilon, \delta}^f \not\in \widetilde{\SIZE}(2^{\beta' n'}; \frac{1}{2} - 2^{-\beta' n'})$
  for some constant $\beta' > 0$.
\end{proof}


Applying the reduction of \cref{thm_Amplification} to \cref{cor_NWAC0},
we obtain the following.
\begin{corollary}
  Let $d \in \Nat, \delta > 0$ be constants.
  Suppose that
  $(\vsProblem{\E}{\widetilde{\SIZE} (2^{\alpha n} ; \delta)} ) \not\in \io \AC^0_{d+2} (N^{1 + \beta})$
  for some constants $\alpha, \beta > 0$.
  Then,
  there exists a hitting set generator $G = \ESet{ G_n : \binset^{O(\log n)} \to \binset^n }_{n \in \Nat} $ computable in time $n^{O(1)}$ and secure against linear-size $\AC^0_{d}$ circuits.
\end{corollary}
\begin{proof}
  We prove the contrapositive.
  Under the assumption that no hitting set generator is secure against $\AC^0_d$,
  it follows from \cref{cor_NWAC0}
  that
  $(\vsProblem{\E}{\widetilde{\AC^{0}_{d+2}} (2^{\beta n} ; \frac{1}{2} - 2^{- \beta n})}) \in \io \AC^0_{d+1} (N^{1 + \gamma})$
  for any constants $\beta, \gamma > 0$.
  Our goal is to prove that
  $(\vsProblem{\E}{\widetilde{\SIZE} (2^{\alpha n} ; \delta)}) \in \io \AC^0_{d+2} (N^{1 + \eta})$
  for any constants $\alpha, \eta > 0$.

  Fix any constants $\alpha, \eta > 0$.
  By \cref{thm_Amplification},
  there exists a $\XOR_{O(\beta n)}$-computable reduction from
  $(\vsProblem{\E}{\widetilde{\SIZE} (2^{\alpha n} ; \delta)})$
  to
  $(\vsProblem{\E}{\widetilde{\AC^{0}_{d+2}} (2^{\beta n} ; \frac{1}{2} - 2^{- \beta n})})$,
  for all sufficiently small $\beta > 0$.
  The latter problem can be solved by an $\AC^0_{d+1}(N^{1 + \gamma})$.
  Thus we obtain an $\AC^0_{d + 1}(N^{1 + \gamma}) \circ \XOR_{O(\beta n)}$ circuit that computes the former problem.
  Since $\XOR_{O(\beta n)}$ can be computed by a depth-$2$ circuit of size $N^{O(\beta)}$,
  by merging one bottom layer of the $\AC^0_{d+1}$ circuit,
  we obtain a circuit in $\AC^0_{d+2} (N^{1 + \gamma + O(\beta)})$ that computes
  $(\vsProblem{\E}{\widetilde{\SIZE} (2^{\alpha n} ; \delta)})$.
  The result follows by choosing $\beta, \gamma > 0$ small enough depending on $\eta > 0$.
\end{proof}

Note that $\AC^0$ circuits are not capable of computing $\XOR$ gates of large fan-in.
If a computational model can compute $\XOR$ gates,
it is possible to compute a locally-decodable error-correcting code.
Specifically, we provide an efficient reduction
from
$(\vsProblem{\E}{{\SIZE}(2^{o( n )})})$
to
$(\vsProblem{\E}{\widetilde{\SIZE}(2^{o( n )}; \frac{1}{2} - 2^{-o( n )}}))$
that is computable by a single layer of $\XOR$ gates.

\begin{theorem}
  [Reductions by Error-Correcting Codes]
  \label{thm_ECCReduction}
  For any constants $\alpha, \gamma > 0$,
  for all sufficiently small $\beta > 0$,
  there exists
  a reduction from
  $(\vsProblem{\E}{{\SIZE}(2^{\alpha n})})$
  to
  $(\vsProblem{\E}{\widetilde{\SIZE}(2^{\beta n}; \frac{1}{2} - 2^{-\beta n}}))$
  that is computable by one layer of $O(N^{1 + \gamma})$ $\XOR$ gates.
\end{theorem}

\begin{lemma}
  %
  [see   %
    \cite{SudanTV01_jcss_journals, Vadhan12_fttcs_journals, Yekhanin12_fttcs_journals}]
  %
  \label{lemma_ECC}
  For any small constant $\beta > 0$,
  for all large $n \in \Nat$,
  there exists an error-correcting code $\Enc^f$
  that encodes a function $f \colon \binset^n \to \binset$
  as a function $\Enc^f \colon \binset^{(1 + O(\sqrt{\beta})) n} \to \binset$
  satisfying following:
  \begin{enumerate}
    \item
          $\size(f) \le
            \ApproxCC{\frac{1}{2} - 2^{-\beta n}}{\Enc^f}
            \cdot 2^{O(\sqrt{\beta} n)}$.
    \item
          For any fixed $y \in \binset^{(1 + O(\sqrt{\beta})) n}$,
          $\Enc^f(y)$ can be computed by an $\XOR$ of some bits of the truth table of $f$.
    \item
          $\Enc^f$ can be computed in time $2^{O(n)}$.
  \end{enumerate}
\end{lemma}
\begin{proof}
  [Proof Sketch]
  We use a Reed--Muller code concatenated with a Hadamard code.
  The crux is that the length of a codeword of can be made small because the query complexity of local-list-decoding algorithms is allowed to be quite large.

  Let $f \colon \binset^n \to \binset$.
  Let $\mathbb{F}_{q}$ be a finite field, where $q = 2^k$ for some $k$.
  Pick $H \subset \mathbb{F}_q$, and encode any element of $\binset^n$ as an element of $H^t$
  by taking an injection $\eta$ from $\binset^n$ to $H^t$,
  where $t$ is some large constant chosen later.
  Let the size $|H|$ of $H$ be $2^{n/t}$.
  Any $f \colon \binset^n \to \binset$ can be
  encoded as a low-degree extension $\widehat{f} \colon \mathbb{F}_q^t \to \mathbb{F}_q$
  such that $\widehat{f}$ and $f \circ \eta^{-1}$ agree on $H^t$.
  Specifically, $\widehat{f}$ can be defined as
  \[
    \widehat{f}(x) = \sum_{y \in H^t} f(\eta^{-1}(y)) \prod_{i=1}^t \delta_{y_i} (x_i)
  \]
  for $x = (x_1, \ldots, x_t) \in \mathbb{F}_q^t$,
  where $\delta_a$ is a polynomial of degree at most $|H|$ such that $\delta_a(a) = 1$ and $\delta_a(b) = 0$ for every $b \in H \setminus \ESet{a}$.
  The total degree of $\widehat{f}$ is at most $d := t |H| = t 2^{n / t}$.
  We will set $q = t 2^{n / t + O(\beta n)}$.

  Then, each alphabet $\widehat{f}(x)$ in the Reed--Muller code is encoded with a Hadamard code.
  Namely, $\Enc^f(x, y) := \langle \widehat{f}(x), y \rangle$,
  where $x \in \mathbb{F}_{q}^t $, $y \in \mathbb{F}_2^k$ and $\langle \blank, \blank \rangle$ denotes the inner product function over $\mathbb F_2$.
  The length of the truth table of $\Enc^f$
  is at most $q^{t + 1} = O(t^{t + 1} 2^{(n / t + O(\beta n))(t + 1)})$.
  By choosing $t := 1 / \sqrt{\beta}$, this is bounded by $2^{(1 + O(\sqrt{\beta})) \cdot n}$.
  The second item holds because
  for every fixed $x$,
  each coordinate of $\widehat{f}(x) \in \mathbb{F}_q \cong \mathbb{F}_2^k$
  can be written as an $\XOR$ of some bits of the truth table of $f$.%
  \footnote{Note that an addition over $\mathbb{F}_q$ is given by a coordinate-wise addition over $\mathbb{F}_2^k$.}

  To see the first item, we use the local list decoding algorithm of
  Sudan, Trevisan, and Vadhan~\cite{SudanTV01_jcss_journals}.
  They gave a local list-decoding algorithm for
  the code $\Enc^f$
  running in time $\poly(t, q)$
  that can handle a $\roundbrackets{\frac{1}{2} - (d / q)^{\Omega(1)}}$-fraction of errors,
  which is more than $\frac{1}{2} - 2^{- \beta n}$
  by choosing $q := t 2^{n / t + O(\beta n)}$ large enough.
  %
  Given a circuit that approximates $\Enc^f$ on a $\frac{1}{2} - 2^{- \beta n}$ fraction of inputs,
  one can apply the local list-decoding algorithm to obtain a circuit that computes $f$ on every input;
  thus we have $\size(f) \le \ApproxCC{\frac{1}{2} - 2^{- \beta n}}{\Enc^f} \cdot 2^{O(\sqrt{\beta}) n}$.
  %
\end{proof}


\begin{proof}
  [Proof of \cref{thm_ECCReduction}]
  We apply the error-correcting code $\Enc$ of \cref{thm_ECCReduction}.
  Specifically, given $f \colon \binset^n \to \binset$ as input,
  we map $f$ to $\Enc^f \colon \binset^{n'} \to \binset$.
  Since the length of $\truthtable(\Enc^f)$ is $2^{n'} = N^{1 + O(\sqrt{\beta})}$ and each bit is computable by a $\XOR$ of some bits of $\truthtable(f)$,
  the reduction is computable by one layer of $O(N^{1 + \gamma})$ $\XOR$ gates
  for all sufficiently small $\beta > 0$.

  We claim the correctness of the reduction.
  Suppose that
  $\Kt(f) = O(\log N)$; then, by the third item of \cref{lemma_ECC}, we have $\Kt(\Enc^f) = O(\log N)$.

  %
  Now suppose that $f \not\in \SIZE(2^{\alpha n})$.
  By \cref{lemma_ECC}, we have
  %
  \[ 2^{\alpha n} < \size(f) \le
    \ApproxCC{\frac{1}{2} - 2^{- \beta n}}{\Enc^f} \cdot 2^{O(\sqrt{\beta} n)}\,.
  \]
  Therefore, we obtain that $\ApproxCC{\frac{1}{2} - 2^{- \beta n}}{\Enc^f} > 2^{(\alpha - O(\sqrt{\beta})) n} \ge 2^{\beta n}$,
  where the last inequality holds by choosing $\beta > 0$ small enough.
  Therefore,
  $\Enc^f \not\in \widetilde{\SIZE}(2^{\beta' n'}; \frac{1}{2} - 2^{-\beta' n'})$
  holds for some constant $\beta' > 0$.
\end{proof}

Applying the reduction, we obtain the following.
\begin{theorem}
  \label{thm_AC0XOR}
  Let $d$ be a constant.
  Suppose that
  $(\vsProblem{\E}{\SIZE (2^{\alpha n})}) \not\in \io \AC^0_{d+1} \circ \XOR (N^{1 + \beta})$
  for some constants $\alpha, \beta > 0$.
  Then,
  there exists a hitting set generator $G = \ESet{ G_n : \binset^{O(\log n)} \to \binset^n }_{n \in \Nat} $ computable in time $n^{O(1)}$ and secure against linear-size $\AC^0_{d} \circ \XOR$ circuits.
\end{theorem}

\begin{proof}
  We prove the contrapositive.
  Let $\alpha, \beta > 0$ be arbitrary constants.
  Assuming that no hitting set generator is secure against $\AC^0_d \circ \XOR$,
  by \cref{thm_GeneralNW}, we have
  $(\vsProblem{\E}{\widetilde{\SIZE} (2^{\alpha_0 n} ; \frac{1}{2} - 2^{- \alpha_0 n})}) \in \io \AC^0_{d+1} \circ \XOR (N^{1 + \beta_0})$
  for any constants $\alpha_0, \beta_0 > 0$.
  By \cref{thm_ECCReduction},
  $(\vsProblem{\E}{{\SIZE}(2^{\alpha n})})$
  is reducible to
  $(\vsProblem{\E}{\widetilde{\SIZE}(2^{\alpha_0 n}; \frac{1}{2} - 2^{-\alpha_0 n})})$
  by one layer of $O(N^{1 + \gamma})$ XOR gates
  for any constant $\gamma > 0$ and any small enough constant $\alpha_0 > 0$.
  Therefore,
  we obtain
  a circuit in $\io\AC^0_{d+1} \circ \XOR \circ \XOR((N^{1 + \gamma})^{1 + \beta_0})$
  that computes
  $(\vsProblem{\E}{{\SIZE}(2^{\alpha n})})$.
  By merging the bottom two $\XOR$ layers,
  the circuit can be written as an $\io\AC^0_{d+1} \circ \XOR$ circuit of size $N^{1 + \gamma + \beta_0 + \gamma \beta_0}$.%
  \footnote{
    Recall that we count the number of gates except for input gates.
  }
  The result follows by choosing positive constants $\gamma, \beta_0 \ll \beta$ small enough.
\end{proof}

We are now ready to complete a proof of \cref{thm_AC0XOREquiv}, which establishes the equivalence between
the existence of a hitting set generator secure against $\AC^0 \circ \XOR$
and the $\AC^0 \circ \XOR$ circuit lower bound for the $\vsProblem{\E}{{\SIZE}(2^{o(n)})}$ Problem.

\begin{restatetheorem}
  {thm_AC0XOREquiv}
  The following are equivalent.
  \begin{enumerate}
    \item
          \label{item_HSGagainstAC0XORRestated}
          For any constant $d$,
          there exists a hitting set generator $G = \ESet{ G_n : \binset^{O(\log n)} \to \binset^n }_{n \in \Nat} $ computable in time $n^{O(1)}$ and secure against linear-size $\AC_d^0 \circ \XOR$ circuits.
    \item
          \label{item_EvsSIZERestated}
          For any constant $d$,
          for some constant $\beta > 0$,
          $(\vsProblem{\E}{\SIZE (2^{o(n)})}) \not\in \io\AC^0_d(N^{1 + \beta})$.
    \item
          \label{item_EvsKtRestated}
          For any constant $d$, for some constant $\beta > 0$,
          $\MKtP[O(\log N), N^{o(1)}] \not\in \io\AC^0_d \circ \XOR (N^{1 + \beta})$.
    \item
          \label{item_EvsKtLargeRestated}
          For any constants $d, k \in \Nat$,
          $\MKtP[O(\log N), N^{o(1)}] \not\in \io\AC^0_d \circ \XOR (N^k)$.
  \end{enumerate}
\end{restatetheorem}

\begin{proof}
  [Proof of \cref{thm_AC0XOREquiv}]
  The implications from \cref{item_EvsKtLargeRestated} to \cref{item_EvsKtRestated} and from \cref{item_EvsKtRestated} to \cref{item_EvsSIZERestated} are trivial.
  The implication from \cref{item_EvsSIZERestated} to \cref{item_HSGagainstAC0XORRestated} immediately follows from \cref{thm_AC0XOR}.
  The implication from \cref{item_HSGagainstAC0XORRestated} to \cref{item_EvsKtLargeRestated} is a standard approach for showing a lower bound for $\MKtP$,
  and follows from \cref{prop_Converse}.
\end{proof}



\subsection{KS complexity and read-once branching program}
\label{section_KSandROBP}

We now turn our attention to $\KS$ complexity.
This amounts to considering a hitting set generator that is computable in a limited amount of space.
Applying our proof ideas to the case of read-once branching programs,
we provide a potential approach for resolving $\BPL = \L$.


\begin{restatetheorem}
  {thm_ROBP}
  There exists a universal constant $\rho > 0$ satisfying the following.
  Suppose that, for some constants $\alpha, \beta > 0$,
  $(\vsProblem{\DSPACE(n)} { \widetilde{\SIZE}(2^{\alpha n}; 2^{- \rho \alpha n }) })$ cannot
  be computed by a read-once co-nondeterministic branching program of size $N^{1 + \beta}$
  for all large input length $N \in \Nat$.
  Then there exists a hitting set generator $G = \ESet{ G_n : \binset^{O(\log n)} \to \binset^n }_{n \in \Nat} $ computable in $O(\log n)$ space and secure against linear-size read-once branching programs.
\end{restatetheorem}

Since the class of read-once branching programs is not closed under taking several reductions presented so far,
we provide a self-contained proof below.

%
%
%
%
%
%
%
%
%
%
%
%
%
%
%
%
%
%
%
%

\begin{proof}
  %
  We prove the contrapositive of \cref{thm_ROBP} for the universal hitting set generator $\UnivGS$.
  Assume that,
  for every function $s(m) = O(\log m)$,
  there exists a linear-size read-once branching program $D$ that avoids $\UnivGS_m^s$
  for infinitely many $m \in \Nat$.
  Given arbitrary constants $c, \alpha, \beta > 0$,
  we will choose a small constant $\gamma > 0$,
  and define $s(m) := c' \log m / \gamma$ for some large constant $c'$.
  Then using $D$ that avoids $\UnivGS_m^s$,
  we present a $\co\RP$-type randomized read-once branching program that solves
  $(\YES, \NO) := (\vsProblem{\dspace{cn}} { \widetilde{\SIZE}(2^{\alpha n}; 2^{- \rho \alpha n }) })$
  on inputs of length $N = 2^n$,
  for some sufficiently large $N := m^{1 / \gamma}$.
  Here, a \emph{randomized} branching program means a probability distribution on branching programs.
  %
  %
  %
  %

  For simplicity, we first explain a construction of a branching program that may not be read-once.
  A randomized branching program $D_0$ is defined as follows.
  Given an input $f \colon \binset^n \to \binset$,
  set $\epsilon := 1/4m$.
  Define ${\widetilde{f}} := \Amp_{\epsilon, \delta}^{f} \colon \binset^{O(\log N)} \to \binset$.
  Let $\ell$ denote the input length of ${\widetilde{f}}$,
  and let $d = O(\ell) = O(n)$ denote the seed length of the Nisan--Wigderson generator instantiated with $\ell$-sized $m$ subsets (\cref{def_NWgenerator}).
  Pick $z \usample \binset^d$
  and output $\neg D(\NW^{{\widetilde{f}}}(z))$.
  This is a randomized branching program because for each fixed $z$,
  each bit of $\NW^{\widetilde{f}}(z)$ is equal to ${\widetilde{f}}(z_{S})$ for some subset $S$,
  and hence by \cref{item_AmpXOR} of \cref{lemma_HardnessAmplification}
  it is some linear combination of at most $k$ bits of $\truthtable(f)$, which can be computed by a read-once width-2 branching program of size $O(k)$.%
  \footnote{Note that, since $z$ is fixed, the branching program does not have to compute $z_S$ from $z$.}
  Thus $\neg D(\NW^{{\widetilde{f}}}(z))$ can be implemented as a branching program of size $O(m \cdot k)$, by replacing each node of $D$ by the read-once width-2 branching programs that compute some linear combinations of $\truthtable(f)$.


  We claim the correctness of the randomized branching program $D_0$.
  \begin{claim}
    \label{claim_RandomizedROBP}
    \hspace{0em}
    \begin{enumerate}
      \item
            $D_0$ accepts every $f \in \YES$ with probability 1.
      \item
            For every $f \in \NO$,
            the probability that $D_0$ rejects $f$ is at least $\frac{1}{4}$.
      \item
            The size of $D_0$ is at most $N^{\beta}$.
    \end{enumerate}
  \end{claim}
  Take any \Yes instance $f \in \YES$.
  We observe that the output $\NW^{\widetilde{f}}(z)$ of the generator has small $\KS$ complexity.
  Indeed,
  \begin{align*}
    \KS(\NW^{{\widetilde{f}}}(z)) \le |z| + \KS({\widetilde{f}}) + O(\log N) \le \KS(f) + O(\log N),
  \end{align*}
  where we used \cref{lemma_design} and \cref{item_AmpKS} of \cref{lemma_HardnessAmplification}.
  In particular, for a large enough constant $c'$, we have
  \begin{align*}
    \KS(f) \le c' \log N = s(m),
  \end{align*}
  and thus $D(\NW^{\widetilde{f}}(z)) = 0$ by \cref{prop_UniversalHSGforSpace} and the assumption that $D$ avoids $\UnivGS_m^s$.
  This means that the algorithm $D_0$ accepts for every choice of $z$.

  Conversely, suppose that the algorithm $D_0$ accepts some input $f$ with probability at least $\frac{3}{4}$.
  We claim that $f \not\in \NO$. %
  The assumption means that
  $ \Pr_z[ D(\NW^{\widetilde{f}}(z)) = 0 ] \ge \frac{3}{4}$, which is equivalent to saying that $\Pr_z[ D(\NW^{\widetilde{f}}(z)) = 1 ] \le \frac{1}{4}$.
  On the other hand, since $D$ avoids $\UnivGS_m^s$, we have $\Pr_w [ D(w) = 1 ] \ge \frac{1}{2}$.
  In particular, we obtain
  \begin{align*}
    \Pr_w[ D(w) = 1 ] - \Pr_z[ D(\NW^{\widetilde{f}}(z)) = 1 ] \ge \frac{1}{4}.
  \end{align*}
  By the security proof of the Nisan--Wigderson generator (\cref{lemma_SecurityOfNW}),
  there exists a one-query oracle circuit $C$ of size $O(m \cdot 2^{\gamma \ell})$ such that
  \begin{align*}
    \Pr_{y \usample \binset^\ell}[ C^D(y) = {\widetilde{f}}(y) ]  \ge \frac{1}{2} + \frac{1}{4m}.
  \end{align*}
  Now replacing the oracle gate of $C$ with a circuit that simulates $D$,
  we obtain a circuit of size $m^{O(1)} + m \cdot N^{O(\gamma)}$.
  By the property of $\Amp_{\epsilon, \delta}$ (\cref{item_AmpApprox} of \cref{lemma_HardnessAmplification}),
  we obtain another circuit $C'$ of size $m^{O(1)} \cdot N^{O(\gamma)} \cdot (1/\delta)^{O(1)}$ such that
  $$
    \Pr_{x \usample \binset^n}[ C'(x) = f(x) ]  \ge 1 - \delta.
  $$
  Choosing $\delta := N^{-\gamma}$,
  we obtain that $\ApproxCC{N^{-\gamma}}{f} \le N^{O(\gamma)}$,
  where $\gamma > 0$ is an arbitrary small constant.
  This completes the proof of the second item of \cref{claim_RandomizedROBP},
  by choosing $\gamma$ small enough so that $O(\gamma) < \alpha$.

  Recall that the size of $D_0$ is $O(m \cdot k)$.
  Here $k$ is the parameter from \cref{lemma_HardnessAmplification} and $k = O( \log ( 1 / \epsilon ) / \delta ) = N^{O(\gamma)}$.
  Thus the size of $D_0$ is at most $N^{O(\gamma)} \le N^\beta$, by choosing $\gamma$ small enough so that $O(\gamma) \le \beta$.
  This completes the proof of \cref{claim_RandomizedROBP}.

  Now we modify the construction of $D_0$ in order to obtain a \emph{read-once} branching program $D_1$.
  The branching program $D_0(\NW^{\widetilde{f}}(z))$ may read some $x$-th bit of the input $\truthtable(f)$ twice only if
  there exists a pair of distinct indices $(i, j)$ such that the $i$th bit of $\NW^{\widetilde{f}}(z)$ and the $j$th bit of $\NW^{\widetilde{f}}(z)$
  are linear combinations of $\truthtable(f)$ that contain $f(x)$.
  We say that a coin flip $z$ is \emph{bad} if this happens.
  To ensure that a coin flip is bad with small probability,
  we take a pairwise independent generator $G_2 \colon \binset^{O(\ell)} \to (\binset^{\ell})^{m}$,
  and define a modified Nisan--Wigderson generator $\NW'^{\widetilde{f}} := {\widetilde{f}}^m \circ (\ND \oplus G_2)$,
  where $\ND \oplus G_2$ denotes the function such that $\ND \oplus G_2(u, v) := \ND(u) \oplus G_2(v)$.
  Using this modified construction,
  the read-once branching program $D_1$ is defined as $\neg D(\NW'^{\widetilde{f}}(z))$ if $z$ is not bad, and otherwise defined as a trivial branching program that outputs 1 always.
  (Note here that since we deal with a non-uniform computation, one does not need to check the badness of $z$ by using a branching program.)
  By the definition, it is obvious that $D_1$ is read-once; hence it remains to claim that $D_1$ satisfies the promise of $(\YES, \NO)$.

  As in the case of $D_0$, for $f \in \YES$, it can be seen that $D(\NW'^{\widetilde{f}}(z)) = 0$ for every $z$, and hence $D_1$ always accepts.
  (We note that the $\KS$ complexity increases by an additive term of the input length of $G_2$, which is $O(\log N)$.)
  Conversely, we claim that if $D_1$ accepts with probability at least $\frac{7}{8}$, then $f \not\in \NO$.
  Assume that $\Pr_{z} [ \text{$D_1$ accepts} ] \ge \frac{7}{8}$.
  We claim that the probability that $z$ is bad is small:
  Fix any distinct indices $(i, j) \in \NumSet{m}^2$.
  Recall that by \cref{item_AmpXOR} of \cref{lemma_HardnessAmplification},
  there exist functions $h_1, \cdots, h_k$ such that
  $
    \widetilde{f}(y) =
    f(h_1(y)) \oplus \cdots \oplus f(h_k(y))$
  for every $y \in \binset^\ell$,
  where $k = O(\log(1/\epsilon) / \delta)$.
  Then, the $i$th bit of $\NW'^{\widehat{f}}(u, v)$ is equal to
  $f(h_1(u_{S_i} \oplus G_2(v)_i))
    \oplus \cdots \oplus f(h_k(u_{S_i} \oplus G_2(v)_i)).
  $
  Fix any distinct indices $i', j' \in \NumSet{k}$.
  %
  %
  %
  %
  %
  %
  Then, we have
  \begin{align*}
    \Pr_{u, v}[h_{i'}(u_{S_i} \oplus G_2(v)_i)
      = h_{j'}(u_{S_j} \oplus G_2(v)_j)]
    = \Pr_{\substack{u \\ a, b \usample \binset^\ell}}[h_{i'}(u_{S_i} \oplus a) = h_{j'}(u_{S_j} \oplus b)]
    %
    = \frac{1}{N},
  \end{align*}
  where the first inequality holds by the pairwise independence of $G_2$
  and the second inequality holds because $h(a)$ is distributed identically with the uniform distribution for $a \usample \binset^\ell$.
  By the union bound over all $(i, j)$ and $(i', j')$,
  the probability that $z$ is bad is bounded above by $(km)^2 \cdot 1/N \le N^{O(\gamma) - 1} \le \frac{1}{8}$ for a sufficiently small $\gamma > 0$ and a large $N \in \Nat$.
  Thus we obtain
  \begin{align*}
    \Pr_z [ D(\NW'^{\widetilde{f}}(z)) = 0 ]
    \ge
    \Pr [ \text{$D_1$ accepts} ] - \Pr [ \text{$z$ is bad} ] \ge \frac{3}{4}.
  \end{align*}
  The rest of a proof of the correctness is essentially the same with the case of $D_0$,
  observing that the security proof of the Nisan--Wigderson generator also works for the modified version $\NW'$.

  Finally, we convert the randomized read-once branching program $D_1$ of size $N^\beta$
  into a co-nondeterministic read-once branching program of size $N^{1 + \beta}$.
  This can be done by using the standard Adleman's trick \cite{Adleman78_focs_conf}:
  Specifically, the success probability of $D_1$ can be amplified to $1 - 2^{-N}$ by taking AND of $O(N)$ independent copies of $D_1$.
  By the union bound, there exists a good coin flip sequence such that
  AND of $O(N)$ copies of $D_1$ solves $(\YES, \NO)$ on every input of length $N$.
  Hard-wiring such a coin flip sequence and simulating the AND gate by using a co-nondeterministic computation,
  we obtain a co-nondeterministic read-once branching program of size $O(N^{1 + \beta})$.
\end{proof}

\section{Non-trivial derandomization and MKtP}

In this section, we provide a characterization of non-trivial derandomization for uniform algorithms by
a lower bound for $\MKtP$.
We start with a formal definition of non-trivial derandomization for uniform algorithms.
%

\begin{definition}
  [Non-trivial derandomization]
  An algorithm $A$ is said to be
  a \emph{derandomization algorithm for $\DTIME(t(n))$ that runs in time $s(n)$}
  if the following hold.
  %
  %
  %
  %
  %
  %
  %
  %
  %
  %
  %
  The algorithm takes $\unary{n}$ and a description of a machine $M$ and outputs an $n$-bit string $A(\unary{n}, M) \in \binset^n$
  in time $s(n)$.
  For any machine $M$ running in time $t(n)$ on inputs of length $n \in \Nat$,
  there exist infinitely many $n \in \Nat$ such that,
  if $\Pr_{x \usample \binset^n} \squarebrackets{ M(x) = 1 } \ge \frac{1}{2}$,
  then $M(A(\unary{n}, M)) = 1$.
\end{definition}

In the following, for a function $s \colon \Nat \to \Nat$, we denote by $R_\Kt^s$ the set $\ISet{ x \in \binset^* }{ \Kt(x) \ge s(|x|) }$
of $\Kt$-random strings with threshold $s$.
We say that a set $R \subset \binset^*$ is dense if $\Pr_{x \usample \binset^n} [ x \in R ] \ge \frac{1}{2}$ for all large $n \in \Nat$.

\begin{proposition}
  \label{prop_DenseSubsetOfKt}
  The following are equivalent for any time-constructible functions $t(n) \ge n$ and $s(n)$.
  \begin{enumerate}
    \item
          \label{item_NontrivialDerandomization}
          There exists a derandomization algorithm for $\DTIME(t(n))$ that runs in time $2^{s(n) + O(\log t(n))}$.
    \item
          \label{item_DenseSubsetOfMKtP}
          Any dense subset of $R_\Kt^{s(n) + O(\log t(n))}$ cannot be accepted by any $t(n)$-time algorithm.
  \end{enumerate}
\end{proposition}
\begin{proof}
  (\cref{item_NontrivialDerandomization} $\implies$ \cref{item_DenseSubsetOfMKtP})
  Let $A$ be a derandomization algorithm for $\DTIME(t(n))$.
  Since $A$ runs in time $2^{s(n) + O(\log t(n))}$ on input $(\unary{n}, M)$,
  the $\Kt$-complexity of $A(\unary{n}, M)$ is at most $s(n) + O(\log t(n)) + |M|$.

  Let $M$ be any $t(n)$-time algorithm such that $\Pr_{x \usample \binset^n} \squarebrackets{ M(x) = 1 } \ge \frac{1}{2}$
  for all large $n \in \Nat$.
  By the property of $A$, we have $M(A(\unary{n}, M)) = 1$ for infinitely many $n$.
  This means that $M$ accepts the string $A(\unary{n}, M)$ that has $\Kt$-complexity at most $s(n) + O(\log t(n))$;
  therefore, $M$ does not accept a dense subset of $\Kt$-random strings.

  (\cref{item_DenseSubsetOfMKtP} $\implies$ \cref{item_NontrivialDerandomization})
  Let $A$ be the algorithm that
  takes $(\unary{n}, M)$,
  enumerates all the strings $x$ whose $\Kt$-complexity is at most $s(n) + O(\log t(n))$,
  and,
  for each string $x$ with small $\Kt$-complexity,
  simulates $M$ on input $x$; if $M$ accepts $x$ in time $t(n)$, output $x$ and halt.
  The running time of $A$ is clearly at most $2^{s(n) + O(\log t(n))}$.
  To prove the correctness, assume towards a contradiction that some algorithm $M$ runs in time $t(n)$ and for all large $n$,
  $\Pr_{x \usample \binset^n} \squarebrackets{ M(x) = 1 } \ge \frac{1}{2}$ and $M(A(\unary{n}, M)) = 0$.
  The latter condition means that $A$ cannot find a string $x$ that is accepted by $M$; thus,
  $M$ rejects all the strings $x$ whose $\Kt$-complexity is at most $s(n) + O(\log t(n))$.
  Therefore, $M$ accepts a dense subset of $\Kt$-random strings, which is a contradiction.
\end{proof}

\begin{theorem}
  \label{thm_WorstToAverageForMKtP}
  Let $t, s \colon \Nat \to \Nat$ be time-constructible functions such that
  $t$ is non-decreasing and
  $\omega(\log^2 n) \le s(n) \le n$
  and
  $\poly(n) \le t(2 s(n))$
  for all large $n$,
  where $\poly$ is some universal polynomial.

  For any constant $c > 0$, there exists a constant $c'$
  such that,
  if
  there is a $t(n)$-time algorithm that accepts a dense subset of $R_{\Kt}^{n - c \sqrt{n} \log n}$,
  then
  the following promise problem can be solved in $\DTIME(O(t(2 s(n)) \cdot 2^{\sqrt{s(n)} \log n})) \cap \co\RTIME(O(t(2 s(n))))$.
  %
  \begin{enumerate}
    \item
          [\Yes:]
          strings $x$ such that
          $\Kt(x) \le s$, where $s := s(|x|)$.
    \item
          [\No:]
          strings $x$ such that
          $\KtTIME{t'}(x) > s + c' \sqrt{s} \log n$,
          where
          $s := s(|x|)$
          and $t' = t(2 s) \cdot \poly(|x|)$.
  \end{enumerate}
\end{theorem}

\begin{lemma}
  %
  [see \cite{Hirahara18_focs_conf}] %
  \label{lemma_NWwithEnc}
  For any $d, m \le n, \epsilon > 0$,
  there exists a black-box pseudorandom generator construction $G^{x} \colon \binset^d \to \binset^m$
  that takes a string $x \in \binset^{n}$ and satisfies the following.
  \begin{enumerate}
    \item
          For any $\epsilon$-distinguisher $T \colon \binset^m \to \binset$ for $G^x$,
          it holds that
          $$\KtTIME[T]{t}(x) \le \exp(\ell^2 / d) \cdot m + d + O(\log (n / \epsilon)),$$ where $\ell = O(\log n)$ and $t = \poly(n, 1 / \epsilon)$.
    \item
          $G^x(z)$ is computable in time $\poly(n, 1 / \epsilon)$.
  \end{enumerate}
\end{lemma}
\begin{proof}
  [Proof Sketch]
  The pseudorandom generator construction $G^x$ is defined as
  $G^x(z) := \NW^{\Enc(x)}(z)$ for any $z \in \binset^d$,
  where $\NW$ is the Nisan--Wigderson generator \cite{NisanW94_jcss_journals} that is instantiated with the weak design of \cite{RazRV02_jcss_journals}
  and the function whose truth table is $\Enc(x)$ as a candidate hard function.
\end{proof}

\begin{proof}
  [Proof of \cref{thm_WorstToAverageForMKtP}]
  Fix $n \in \Nat$ and an input $x \in \binset^n$.
  For simplicity, let $s := s(n)$.
  Define $d := \sqrt{s} \log n$.  Since $s(n) = \omega(\log^2 n)$, we have $4 c d \le s$ for a sufficiently large $n$.
  Set $m := s + 4 c d \le 2 s$.
  Let $D$ be the $t(n)$-time algorithm that accepts a dense subset of $R_\Kt^{n - c \sqrt{n} \log n}$.

  We claim that there exists a $\co\RP$-type algorithm $A$ that
  uses $d$ random bits and
  solves the promise problem.
  The algorithm $A$ operates as follows.
  Pick a random seed $z \in \binset^d$, and accept if and only if $D(G^x(z)) = 0$,
  where the output length of $G^x$ is set to be $m$.
  This runs in time $t(m) + \poly(n) \le O(t(2 s))$ and can be derandomized in time $O(t(2 s)) \cdot 2^d$ by exhaustively trying all the random bits.

  It remains to prove the correctness of the algorithm $A$.
  Assume that $\Kt(x) \le s$.
  Then, since $G^x(z)$ is computable in polynomial time,
  for any $z \in \binset^d$, we have
  $$\Kt(G^x(z)) \le d + s + O(\log n) \le s + 2 d \le s + 4 c d - 2 c d < m - c \sqrt{m} \log n,$$
  where, in the last inequality, we used the fact that $\sqrt{m} \log n \le \sqrt{2 s} \log n < 2 d$.
  Therefore, $G^x(z) \not\in R_\Kt^{m - c \sqrt{m} \log m}$, which is rejected by $D$; hence, $A$ accepts.

  Conversely, assume that the probability that $A$ accepts is at least $\frac{3}{4}$.
  This means that $$\Pr_{z \usample \binset^d} \squarebrackets{ D(G^x(z)) = 0 } \ge \frac{3}{4}.$$
  Since we have $$\Pr_{w \usample \binset^m} \squarebrackets{ D(w) = 0 } \le \frac{1}{2},$$
  $D$ is a distinguisher for $G^x$.
  By the security of \cref{lemma_NWwithEnc},
  we obtain that, for $t' := \poly(n) \cdot t(m)$,
  \begin{align*}
    \KtTIME{t'}(x) & \le \exp(\ell^2 / d) \cdot m + O(\log n)
    \\ & \le (1 + 2 \ell^2 / d) \cdot (s + 4 c d) + O(\log n)
    \\ & \le s + O(\sqrt{s} \log n),
  \end{align*}
  which can be bounded above by $s + c' \sqrt{s} \log n$
  by choosing a large enough constant $c'$.
  Thus, $x$ is not a \No instance of the promise problem.
  Taking the contrapositive,
  we conclude that any \No instance $x$ is rejected by $A$ with probability at least $\frac{1}{4}$.
\end{proof}

In the special case when $t(n) = n^{O(1)}$, \cref{thm_WorstToAverageForMKtP} implies that the $\vsProblem{\Kt}{\KtTIME{t}}$ Problem can be solved
in $\co\RP$.
\begin{proof}
  [Proof Sketch of \cref{thm_KtvsKt}]
  We claim that the $\vsProblem{\Kt}{\KtTIME{t}}$ Problem can be solved in
  time $2^{\widetilde{O}(\sqrt{n})}$
  under the assumption that $\MKtP \in \P$.
  We use the same proof of \cref{thm_WorstToAverageForMKtP} for $t(n) := n^{\log n}$ when the parameter $s = \omega(\log^2 n)$;
  when $s = O(\log^2 n)$, we set $d := \ell^2$ and $m := O(d)$ as in \cite{Hirahara18_focs_conf}.
\end{proof}

\begin{theorem}
  [A formal version of \cref{thm_NontrivialDerandomizationInformal}]
  \label{thm_NontrivialDerandomizationFormal}
  For any constant $0 < \epsilon < 1$, the following are equivalent.
  \begin{enumerate}
    \item
          \label{item_NontrivialDerandomizationInTheorem}
          For some $c_1$, for any large $c_2$,
          there exists no derandomization algorithm for $\DTIME(2^{c_1 \sqrt{n} \log n})$ that runs in time $2^{n - c_2 \sqrt{n} \log n}$.
    \item
          \label{item_MKtPDenseSubset}
          For some $c_1$, for any large $c_2$,
          there exists an algorithm running in time $2^{c_1 \sqrt{n} \log n}$ that accepts a dense subset of $R_{\Kt}^{n - c_2 \sqrt{n} \log n}$.
    \item
          \label{item_MKtPLowerBoundInTheorem}
          For some $c_1$, for any large $c_2$,
          $\MKtP[n - c_2 \sqrt{n} \log n, n-1] \in \DTIME(2^{c_1 \sqrt{n} \log n})$.
    \item
          \label{item_MKtPLowerBoundSmallInTheorem}
          For some $c_1$, for any large $c_2$,
          $\MKtP[n^\epsilon, n^\epsilon + c_2 \sqrt{n^\epsilon} \log n] \in \DTIME(2^{c_1 \sqrt{n^\epsilon} \log n})$.
  \end{enumerate}
\end{theorem}

\begin{proof}
  (\cref{item_NontrivialDerandomizationInTheorem} $\Longleftrightarrow$ \cref{item_MKtPDenseSubset})
  This equivalence follows from \cref{prop_DenseSubsetOfKt}.

  (\cref{item_MKtPDenseSubset} $\implies$ \cref{item_MKtPLowerBoundSmallInTheorem})
  We apply \cref{thm_WorstToAverageForMKtP} for $t(n) := 2^{c_1 \sqrt{n} \log n}$ and $s(n) := n^\epsilon$.
  Then we obtain a $\DTIME(O(t(2 s)) \cdot 2^{\sqrt{s} \log n})$-algorithm $A$ that distinguishes
  \Yes instances $x$ such that $\Kt(x) \le s$
  from \No instances $x$ such that $\KtTIME{t'}(x) > s + c_2' \sqrt{s} \log n$, where $t' = t(2s) \cdot \poly(n)$ and $c_2'$ is some constant.
  We choose a constant $c_1'$ large enough so that
  $O(t(2 s)) \cdot 2^{\sqrt{s} \log n} \le 2^{c_1' \sqrt{s} \log n}$, which bounds from above the running time of the algorithm $A$.

  We claim that $\MKtP[s, s + c_2'' \sqrt{s} \log n]$ is also solved by $A$ for any sufficiently large $c_2''$.
  Take any string $x$ such that $\Kt(x) \ge s + c_2'' \sqrt{s} \log n$.
  Since $\Kt(x) \le \KtTIME{t'}(x) + O(\log t') \le \KtTIME{t'}(x) + O(c_1 \sqrt{s} \log n)$,
  it follows that $\KtTIME{t'}(x) \ge s + c_2'' \sqrt{s} \log n - O(c_1 \sqrt{s} \log n) > s + c_2' \sqrt{s} \log n$,
  where the last inequality holds for any sufficiently large $c_2''$; thus, $x$ is rejected by $A$.

  (\cref{item_MKtPLowerBoundSmallInTheorem} $\implies$ \cref{item_MKtPLowerBoundInTheorem})
  By the assumption, there exists a constant $c_1$ such that, for all large $c_2$,
  there exists an algorithm $A$ that solves $\MKtP[n^\epsilon, n^\epsilon +  c_2 \sqrt{n^\epsilon} \log n]$ in time $2^{c_1 \sqrt{n^\epsilon} \log n}$.
  Define $c_1' := c_1 / \epsilon$.
  For all large $c_2'$,
  we construct an algorithm $B$ that solves $\MKtP[n - c_2' \sqrt{n} \log n, n - 1]$ in time $2^{c_1' \sqrt{n} \log n}$.
  The algorithm $B$ takes an input $x$ of length $n$ and simulates $A$ on input $x 1 0^{m - n - 1}$ for $m := (n - 2 c_2 \sqrt{n} \log n)^{1/\epsilon}$.
  The running time of $B$ is at most $2^{c_1 \sqrt{m} \log m} \le 2^{c_1 / \epsilon \cdot \sqrt{m} \log n} = 2^{c_1' \sqrt{m} \log n}$.

  It remains to prove the correctness of $B$.
  Take any $x \in \binset^n$ such that $\Kt(x) \le n - c_2' \sqrt{n} \log n$.
  Then we have
  $\Kt(x10^{m-n-1}) \le n - c_2' \sqrt{n} \log n + O(\log n) \le m^\epsilon$ for any large $c_2$; thus $B$ accepts $x$.
  Conversely, take any $x$ such that $\Kt(x) \ge n - 1$.
  Then we have $\Kt(x10^{m-n-1}) \ge n - O(\log n) = m^{\epsilon} + 2 c_2 \sqrt{n} \log n - O(\log n) \ge m^\epsilon + c_2 \sqrt{m^\epsilon} \log n$;
  thus, $B$ rejects.

  (\cref{item_MKtPLowerBoundInTheorem} $\implies$ \cref{item_MKtPDenseSubset})
  This immediately follows from the fact that
  the complement of $\MKtP[n - c_2 \sqrt{n} \log n, n - 1]$ is a dense subset of $R_\Kt^{n - c_2 \sqrt{n} \log n + 1}$.
  %
  %
  %
\end{proof}

%
%
%
%
%
%
%
%
%
%
%
%
%
%
%
%
%
%
%
%
%
%
%
%
%
%
%
%
%
%
%
%
%
%
%
%

\appendix

\section{Derandomized hardness amplification theorem}
\label{section_HardnessAmplification}

We review a simplified proof of derandomized hardness amplification given by Healy, Vadhan and Viola \cite{HealyVV06_siamcomp_journals}, and observe that the parameter shown in \cref{lemma_HardnessAmplification} can be achieved by slightly modifying the construction.

\begin{theoremref}
  [Derandomized hardness amplification]
  {lemma_HardnessAmplification}
  Let $\gamma > 0$ be an arbitrary constant.
  There exists a hardness amplification procedure $\Amp$
  that takes a function $f \colon \binset^n \to \binset$ and parameters $\delta, \epsilon > 0$,
  and returns
  a Boolean function
  $\Amp_{\epsilon, \delta}^f \colon \binset^{O(n + \log (1/\epsilon))} \to \binset$
  satisfying the following:
  \begin{enumerate}
    \item
          If $\ApproxCC{1/2 - \epsilon}{\Amp_{\epsilon, \delta}^f} \le s$,
          then
          $\ApproxCC{\delta}{f} \le s \cdot \poly(1/\epsilon, 1/\delta)$.
    \item
          For any fixed $y \in \binset^{O(n + \log (1 / \epsilon))}$,
          there exist strings $v_1, \cdots, v_k \in \binset^{n}$
          for some $k = O(\log (1 / \epsilon) / \delta)$
          such that
          $\Amp_{\epsilon, \delta}^f(y) =  f(v_1) \oplus \cdots \oplus f(v_k)$.
          Moreover, if $y$ is distributed uniformly at random,
          for each $i \in \NumSet{k}$,
          $v_i$ is distributed uniformly at random.
    \item
          $\Amp_{\epsilon, \delta}^f(y)$ can be computed in $O(n + \log (1 / \delta) + \log (1 / \epsilon))$ space,
          given $f, y, \epsilon$ and $\delta$ as an input.
          %
  \end{enumerate}

\end{theoremref}


First, we explain the construction of our hardness amplification procedure.
The construction is
a XOR of the nearly disjoint generator $\ND$ and a hitter $\Hit$
%
(see \cite{Goldreich11j_goldreich2011_sp_books}). %
In fact, we need a slightly generalized version of a hitter, for which we show that the same construction of \cite{Goldreich11j_goldreich2011_sp_books} suffices.
We note that in the previous works \cite{ImpagliazzoW97_stoc_conf, HealyVV06_siamcomp_journals}, an expander walk was used instead of $\Hit$;
this does not give us a nearly optimal construction of a hitter when $\delta$ is not a constant. %

\begin{lemma}
  %
  \label{lemma_hitter}
  There exists a \emph{``hitter''} $\Hit$ such that, given parameters $n \in \Nat, \epsilon, \delta > 0$,
  a function $\Hit_{n, \epsilon, \delta} \colon \binset^{O(n + \log( 1 / \epsilon) )} \to (\binset^n)^{k_\ParamForAmp}$ takes a seed of length $O(n + \log ( 1 / \epsilon ))$
  and outputs
  a list $L$ of
  $k_\ParamForAmp = O(\log(1 / \epsilon) / \delta)$ strings of length $n$
  such that,
  for any subsets $H_1, \cdots, H_{k_\ParamForAmp} \subset \binset^n$ of size $\ge \delta 2^n$,
  with probability at least $1 - \epsilon$,
  there exists an index $i \in \NumSet{k_\ParamForAmp}$ such that the $i$th string in the list $L$ is in $H_i$.
  %
  Moreover, $\KS(L) \le O(n + \log ( 1 / \epsilon ))$.
\end{lemma}

\begin{proof}
  [Proof Sketch]
  We observe that the construction of \cite[Appendix C]{Goldreich11j_goldreich2011_sp_books} satisfies the requirement of \cref{lemma_hitter}.
  The construction is as follows:
  First, we take a pairwise independent generator $G_2 \colon \binset^{O(n)} \to (\binset^n)^{O(1 / \delta)}$.
  By Chebyshev's inequality, with probability at least $\frac{1}{2}$ over the choice of a seed $r \in \binset^{O(n)}$,
  $G_2$ hits some $H_i, \cdots, H_{i + O(1 / \delta)}$, where $i$ is an arbitrary index.
  Now we take an explicit construction of a constant-degree expander on $2^{O(n)}$ vertices,
  and generate a random walk $v_1, \cdots, v_\ell$ of length $\ell = O(\log (1 / \epsilon))$ over $\binset^{O(n)}$,
  which takes random bits of length $O(n + \log (1 / \epsilon))$.
  The output of $\Hit$ is defined as the concatenation of $G_2(v_1), \cdots, G_2(v_\ell)$.
  The correctness follows by using the Expander Random Walk Theorem \cite[Theorem A.4]{Goldreich11j_goldreich2011_sp_books}.
\end{proof}

%
%

\begin{definition}
  Let $f$ be a function $f \colon \binset^n \to \binset$, and $\epsilon, \delta > 0$ be arbitrary parameters.
  Let $k := k_{n, \epsilon, \delta}$.
  Let $\ND \colon \binset^{O(n)} \to (\binset^n)^{k}$ be the nearly disjoint generator (\cref{def_NDgenerator}) defined with the design of \cref{lemma_design}.
  We define a generator $$\IW_\ParamForAmp \colon \binset^{O(n + \log(1 / \epsilon))} \to (\binset^n)^{k}$$ as $$\IW_\ParamForAmp(x, y) := \ND(x) \oplus \Hit(y).$$
  Then we
  define a hardness amplified version $$\Amp_{\epsilon, \delta}^f \colon \binset^{O(n + \log(1 / \epsilon))} \to \binset$$ of $f$
  as
  the function $\Amp_{\epsilon, \delta}^f := f^{\oplus k} \circ \IW_\ParamForAmp$.
  %
\end{definition}

We proceed to a proof of \cref{lemma_HardnessAmplification}.
%
%
We recall  %
several notions from \cite{HealyVV06_siamcomp_journals}.
For a subset $H \subset \binset^n$ (which is supposed to be a hard-core set) and a function $f \colon \binset^n \to \binset$,
we consider a \emph{probabilistic function} $f_H \colon \binset^n \to \binset$, \ie, a distribution over %
Boolean functions,  %
defined as follows:
For each $x \in H$, the value $f_H(x)$ is defined as a random bit chosen uniformly at random and independently;
For each $x \in \binset^n \setminus H$, the value $f_H(x)$ is defined as $f(x)$.
Assuming that $H$ is indeed a hard-core set for $f$, the distributions $(x, f(x))$ and $(x, f_H(x))$ where $x \usample \binset^n$ are computationally indistinguishable. Indeed:
\begin{proposition}
  %
  [see \cite{Trevisa03_focs_conf}]  %
  \label{prop_HardcoreAndProbabilisticFunction}
  Let $H \subset \binset^n$ be an arbitrary subset, $f \colon \binset^n \to \binset$ be an arbitrary function, and $\epsilon > 0$ be an arbitrary parameter.
  If
  \begin{align*}
    \Pr_{x \usample \binset^n}[ A(x, f(x)) = 1 ] - \Pr_{x \usample \binset^n}[ A(x, f_H(x)) = 1 ] \ge \epsilon,
  \end{align*}
  for some function $A \colon \binset^{n + 1} \to \binset$,
  then
  there exists a one-query oracle circuit of size $O(1)$ such that
  \begin{align*}
    \Pr_{x \usample H} [ C^A(x) = f(x) ] \ge \frac{1}{2} + \frac{\epsilon} {\delta}.
  \end{align*}
\end{proposition}

%
%
For a probabilistic function $h \colon \binset^n \to \binset$,
the expected bias of $h$ is defined as $\ExpBias[h] := \Exp_{x \usample \binset^n} [ \Bias(h(x)) ]$,
%
%
%
%
%
where the bias $\Bias(b)$ of a binary random variable $b \in \binset$ is defined as
\[ \Bias(b) = \left| \Pr_b[b = 0] - \Pr_b[b = 1] \right| \,. \]

It was shown in \cite{HealyVV06_siamcomp_journals} that the hardness of $f^{\oplus k} \circ G(z)$ is essentially characterized
by the expected bias of $\ExpBias[ f_H^{\oplus k} \circ \IW ]$,
%
%
using a property of the nearly disjoint generator $\ND$.
\begin{lemma}
  [{\cite[Lemma 5.2 and Lemma 5.12]{HealyVV06_siamcomp_journals}}]
  \label{lemma_ExpectedBiasCharacterization}
  Let $g$ be an arbitrary probabilistic function, and $\epsilon> 0$ be arbitrary parameters.
  Suppose that there exists a function $A$ such that
  $$
    \Pr_z [ A(z) = f^{\oplus k} \circ \IW (z) ] \ge \frac{1}{2} + \frac{1}{2} \ExpBias[ g^{\oplus k} \circ \IW ] + \frac{\epsilon}{2}
  $$
  Then there exists a one-query oracle circuit $C$ of size $O(k \cdot 2^{\gamma n})$ such that
  $$
    \Exp_x [ C^A(x, f(x)) - C^A(x, g(x)) ] \ge \frac{\epsilon}{ 2k},
  $$
  where $\gamma > 0$ is an arbitrary constant of \cref{lemma_design}.
\end{lemma}

We will apply \cref{lemma_ExpectedBiasCharacterization} for $g := f_H$.
By using a property of the hitter, we show that the expected bias of $f_H$ is small whenever a hard-core set $H$ is large enough.
\begin{lemma}
  \label{lemma_HittingHardCore}
  Let $\epsilon, \delta > 0$ be arbitrary parameters, and $H \subset \binset^n$ be a subset of size at least $\delta 2^n$.
  Let $k := k_\ParamForAmp$.
  Then
  $\ExpBias[f_H^{\oplus k} \circ \IW_\ParamForAmp] \le \epsilon$
\end{lemma}
\begin{proof}
  The idea is that when a hitter hits a hard-core set $H$, the expected bias becomes $0$.
  More specifically,
  recall that $\IW(x, y)$ is defined as $\ND(x) \oplus \Hit_\ParamForAmp(y)$.
  Fix any $x$, and define $H_i$ as a shifted version of $H$: namely, $H_i := \ND(x)_i \oplus H := \ISet{\ND(x)_i \oplus h}{h \in H}$ for every $i \in \NumSet{k}$.
  We apply \cref{lemma_hitter} for $H_1, \cdots, H_k$.
  Since $|H_i| = |H| \ge \delta 2^n$ for every $i \in \NumSet{k}$, by the property of the hitter,
  there exists some index $i \in \NumSet{k}$
  such that
  $\Hit_\ParamForAmp(y)_i \in H_i = \ND(x)_i \oplus H$ with probability at least $1 - \epsilon$ over the choice of $y$.
  In this case,
  $\IW_\ParamForAmp(x, y)_i = \ND(x)_i \oplus \Hit_\ParamForAmp(y)_i \in H$,
  and hence the bias of $f_H(\IW_\ParamForAmp(x, y))$ is exactly $0$.
  Therefore,
  \begin{align*}
    \  & \ExpBias[f_H^{\oplus k} \circ \IW_\ParamForAmp]
    \\ \le\ & \Pr_{x, y} \left[ \Bias( f_H^{\oplus k} \circ \IW_\ParamForAmp (x, y) )  > 0\right]
    \\ \le\ & \Pr_{x, y} \left[\, \IW_\ParamForAmp(x, y)_i \not\in H \text{ for every $i\in\NumSet{k}$}\, \right]\ \le\ \epsilon
  \end{align*}
\end{proof}

Now we combine \cref{prop_HardcoreAndProbabilisticFunction,lemma_ExpectedBiasCharacterization,lemma_HittingHardCore}
and obtain the following:

\begin{corollary}
  \label{cor_BoostingPremise}
  Let $n \in \Nat$ and $\epsilon, \delta > 0$.
  Let $H \subset \binset^n$ be an arbitrary subset of size at least $\delta 2^n$.
  If some function $A$ satisfies
  \begin{align*}
    \Pr_z[ A(z) = f^{\oplus k} \circ \IW_\ParamForAmp(z) ] \ge \frac{1}{2} + \epsilon
  \end{align*}
  then
  there exists a one-query oracle circuit $C$ of size $O(k \cdot 2^{\gamma n})$ such that,
  \begin{align*}
    \Pr_{x \usample H} [ C^A(x) = f(x) ] \ge \frac{1}{2} + \frac{\epsilon } { 2 \delta k }.
  \end{align*}
\end{corollary}

At this point, we make use of Impagliazzo's hard-core lemma \cite{Impagliazzo95_coco_conf}.
Equivalently, one can view it as a boosting algorithm
%
(see \cite{KlivansS03_ml_journals}).   %
\begin{lemma}
  %
  [see   %
    \cite{Impagliazzo95_coco_conf, KlivansS03_ml_journals}]
  \label{lemma_Hardcore}
  Under the condition of \cref{cor_BoostingPremise},
  there exists an oracle circuit $C$ of size $O(k \cdot 2^{\gamma n}) \cdot \poly(k / \epsilon)$
  such that
  $$
    \Pr_{x \usample \binset^n} [ C^A(x) = f(x) ] \ge 1 - \delta.
  $$
\end{lemma}

Now we take any circuit $A$ of size $s$ such that
$$
  \Pr_{z} [  A(z) = \Amp_{\epsilon, \delta}^f(z) ] \ge \frac{1}{2} + \epsilon.
$$
By using \cref{cor_BoostingPremise,lemma_Hardcore}
and replacing each oracle gate by a circuit $A$,
we obtain a circuit $C$ of size $O(k \cdot 2^{\gamma n}) \cdot \poly(k / \epsilon) \cdot s$
such that
$$
  \Pr_{x \usample \binset^n} [ C(x) = f(x) ] \ge 1 - \delta.
$$
This completes the proof of \cref{lemma_HardnessAmplification}.

\section{\texorpdfstring{$\AC^0$}{AC0} Lower Bounds for MKtP}
\label{section_AC0MKtP}

In this section, we prove that there exists no constant-depth fixed-polynomial-size circuit that computes $\MKtP[O(\log N), N^{o(1)}]$.
\begin{proposition}
  \label{prop_AC0versusMKtP}
  For any constants $\alpha < 1, k, d \in \Nat$,
  there exists a constant $c$ such that
  $$\MKtP[c \log N, N^{\alpha}] \not\in \io \AC_d^0(N^k).$$
\end{proposition}

Previously, \cite{ChenHOPRS20_innovations_conf} used the pseudorandom restriction of Trevisan and Xue~\cite{TrevisanX13_coco_conf}
to obtain $\AC^0$ lower bounds for $\MKtP[\polylog N]$.
Here we make use of a polynomial-time-computable pseudorandom restriction that shrinks $\AC^0$ circuits,
which enables us to prove a lower bound for $\MKtP[O(\log N)]$.
The following lemma is proved in the context of quantified derandomization.
\begin{lemma}
  [Goldreich and Wigderson~\cite{GoldreichW14_stoc_conf}]
  \label{lemma_pseudorandomrestriction}
  For any constants $\alpha < 1$ and $d \in \Nat$ and any polynomials $p, q$,
  there exists a constant $k$ and a polynomial-time algorithm of $O(\log n)$ randomness complexity that
  produces restrictions on $n$ variables such that the following conditions hold:
  \begin{enumerate}
    \item
          The number of undetermined variables in each restriction is at least $2 n^\alpha$.
    \item
          For any $n$-input circuit of depth $d$ and size at most $p(n)$,
          with probability at least $1 - 1 / q(n)$,
          the corresponding restricted circuit depends on at most $k$ variables.
  \end{enumerate}
\end{lemma}

\begin{proof}
  [Proof of \cref{prop_AC0versusMKtP}]
  Fix any polynomial $p$ and a depth $d$.
  We claim that,
  for some constant $c > 0$,
  there exists no depth-$d$ circuit of size $p(n)$ that computes $\MKtP[c \log n, n^{\alpha}]$ for all large $n$.
  Assume, towards a contradiction, that there exists a depth-$d$ circuit $C$ of size $p(n)$ that computes $\MKtP[c \log n, n^{\alpha}]$.
  We use \cref{lemma_pseudorandomrestriction} for $q \equiv 2$.
  Then, there exists a polynomial-time algorithm that produces a pseudorandom restriction $\rho$ that shrinks $C$ to a circuit of size $k = O(1)$
  with probability at least $\frac{1}{2}$.
  Fix one pseudorandom restriction $\rho$ such that the restricted circuit $C\restriction_\rho$ is a constant-size circuit.

  We claim that $C$ does not compute $\MKtP[c \log n, n^{\alpha}]$.
  Let $\sigma$ be the restriction that fixes $k$ variables on which $C\restriction_\rho$ depend to $0$.
  Consider $0^n \circ \rho = 0^n \circ \sigma \circ \rho$, \ie, the string obtained by fixing undetermined variables in $\rho$ to $0$.
  Since $\rho$ is generated by a polynomial-time algorithm,
  there exists a constant $c$ such that $\Kt(0^n \circ \rho) \le c \log n$.
  Thus, $0^n \circ \sigma \circ \rho$ is a \Yes instance of $\MKtP[c \log n, n^{\alpha}]$.
  %
  %
  Since $C\restriction_{\sigma \circ \rho}$ is a constant circuit, we have $1 = C\restriction_{\sigma \circ \rho}(0^n) = C\restriction_{\sigma \circ \rho}(x)$
  for any $x \in \binset^n$.
  However, for a random $x \in \binset^n$,
  by a simple counting argument, it holds that $\Kt(x \circ \sigma \circ \rho) \ge 2 n^\alpha - k - 1 > n^\alpha$ with high probability.
  Therefore, $x \circ \sigma \circ \rho$ is a \No instance of $\MKtP[c \log n, n^{\alpha}]$ for some $x$, which is a contradiction.
\end{proof}




%% Theory of Computing
%% File: v019a004.bbl
%% Publication date: October 14, 2023
%% Authors: Shuichi Hirahara
\newcommand{\prelim}{Preliminary version in\ }\newcommand{\prelims}{Preliminary
  versions in\ }
% Define empty bibhead if not already defined
\providecommand{\bibhead}[1]{}
% Define tocrefpdfbookmark if not already defined
\expandafter\ifx\csname pdfbookmark\endcsname\relax%
  \providecommand{\tocrefpdfbookmark}{}
\else\providecommand{\tocrefpdfbookmark}{%
   \hypertarget{tocreferences}{}%
   \pdfbookmark[1]{References}{tocreferences}}%
\fi

\tocrefpdfbookmark
\begin{thebibliography}{10}

\bibitem{Adleman78_focs_conf}\bibhead{Adleman78_focs_conf}
{\sc Leonard~M. Adleman}: Two theorems on random polynomial time.
\newblock In {\em Proc. 19th FOCS}, pp. 75--83. IEEE Comp. Soc., 1978.
\newblock [\epfmtdoi{10.1109/SFCS.1978.37}]

\bibitem{AllenderBKMR06_siamcomp_journals}\bibhead{AllenderBKMR06_siamcomp_journals}
{\sc Eric Allender, Harry Buhrman, Michal Kouck{\'{y}}, Dieter van Melkebeek,
  and Detlef Ronneburger}: Power from random strings.
\newblock {\em SIAM J. Comput.}, 35(6):1467--1493, 2006.
\newblock [\epfmtdoi{10.1137/050628994}]

\bibitem{AllenderHMPS08_siamcomp_journals}\bibhead{AllenderHMPS08_siamcomp_journals}
{\sc Eric Allender, Lisa Hellerstein, Paul McCabe, Toniann Pitassi, and
  Michael~E. Saks}: Minimizing disjunctive normal form formulas and
  $\mathrm{AC}^0$ circuits given a truth table.
\newblock {\em SIAM J. Comput.}, 38(1):63--84, 2008.
\newblock [\epfmtdoi{10.1137/060664537}]

\bibitem{AllenderH19_toct_journals}\bibhead{AllenderH19_toct_journals}
{\sc Eric Allender and Shuichi Hirahara}: New insights on the (non-)hardness of
  circuit minimization and related problems.
\newblock {\em ACM Trans. Comput. Theory}, 11(4):27:1--27, 2019.
\newblock [\epfmtdoi{10.1145/3349616}]

\bibitem{AllenderK10_jacm_journals}\bibhead{AllenderK10_jacm_journals}
{\sc Eric Allender and Michal Kouck{\'{y}}}: Amplifying lower bounds by means
  of self-reducibility.
\newblock {\em J. ACM}, 57(3):14:1--36, 2010.
\newblock [\epfmtdoi{10.1145/1706591.1706594}]

\bibitem{BabaiFNW93_cc_journals}\bibhead{BabaiFNW93_cc_journals}
{\sc L{\'{a}}szl{\'{o}} Babai, Lance Fortnow, Noam Nisan, and Avi Wigderson}:
  \emph{BPP} has subexponential time simulations unless \emph{EXPTIME} has
  publishable proofs.
\newblock {\em Comput. Complexity}, 3(4):307--318, 1993.
\newblock [\epfmtdoi{10.1007/BF01275486}]

\bibitem{BlumM84_siamcomp_journals}\bibhead{BlumM84_siamcomp_journals}
{\sc Manuel Blum and Silvio Micali}: How to generate cryptographically strong
  sequences of pseudo-random bits.
\newblock {\em SIAM J. Comput.}, 13(4):850--864, 1984.
\newblock [\epfmtdoi{10.1137/0213053}]

\bibitem{BogdanovT06_siamcomp_journals}\bibhead{BogdanovT06_siamcomp_journals}
{\sc Andrej Bogdanov and Luca Trevisan}: On worst-case to average-case
  reductions for {NP} problems.
\newblock {\em SIAM J. Comput.}, 36(4):1119--1159, 2006.
\newblock [\epfmtdoi{10.1137/S0097539705446974}]

\bibitem{BorodinRS93_cc_journals}\bibhead{BorodinRS93_cc_journals}
{\sc Allan Borodin, Alexander~A. Razborov, and Roman Smolensky}: On lower
  bounds for read-$k$-times branching programs.
\newblock {\em Comput. Complexity}, 3(1):1--18, 1993.
\newblock [\epfmtdoi{10.1007/BF01200404}]

\bibitem{BuhrmanFP05_mst_journals}\bibhead{BuhrmanFP05_mst_journals}
{\sc Harry Buhrman, Lance Fortnow, and Aduri Pavan}: Some results on
  derandomization.
\newblock {\em Theory Computing Sys.}, 38(2):211--227, 2005.
\newblock [\epfmtdoi{10.1007/s00224-004-1194-y}]

\bibitem{BuhrmanFS09_icalp_conf}\bibhead{BuhrmanFS09_icalp_conf}
{\sc Harry Buhrman, Lance Fortnow, and Rahul Santhanam}: Unconditional lower
  bounds against advice.
\newblock In {\em Proc. 36th Internat. Colloq. on Automata, Languages, and
  Programming (ICALP'09)}, pp. 195--209. Springer, 2009.
\newblock [\epfmtdoi{10.1007/978-3-642-02927-1\_18}]

\bibitem{BuhrmanH92_fsttcs_conf}\bibhead{BuhrmanH92_fsttcs_conf}
{\sc Harry Buhrman and Steven Homer}: Superpolynomial circuits, almost sparse
  oracles and the exponential hierarchy.
\newblock In {\em Proc. 12th Found. Softw. Techn. Theoret. Comp. Sci. Conf.
  (FSTTCS'92)}, pp. 116--127. Springer, 1992.
\newblock [\epfmtdoi{10.1007/3-540-56287-7\_99}]

\bibitem{CarmosinoIKK16_coco_conf}\bibhead{CarmosinoIKK16_coco_conf}
{\sc Marco~L. Carmosino, Russell Impagliazzo, Valentine Kabanets, and Antonina
  Kolokolova}: Learning algorithms from {N}atural {P}roofs.
\newblock In {\em Proc. 31st Comput. Complexity Conf. (CCC'16)}, pp. 10:1--24.
  Schloss Dagstuhl--Leibniz-Zentrum fuer Informatik, 2016.
\newblock [\epfmtdoi{10.4230/LIPIcs.CCC.2016.10}]

\bibitem{ChenHOPRS20_innovations_conf}\bibhead{ChenHOPRS20_innovations_conf}
{\sc Lijie Chen, Shuichi Hirahara, Igor~Carboni Oliveira, J{\'{a}}n Pich, Ninad
  Rajgopal, and Rahul Santhanam}: Beyond {N}atural {P}roofs: {H}ardness
  magnification and locality.
\newblock {\em J. ACM}, 69(4):25:1--49, 2022.
\newblock \prelim \href{https://doi.org/10.4230/LIPIcs.ITCS.2020.70}{ITCS'20}.
\newblock [\epfmtdoi{10.1145/3538391}]

\bibitem{ChenJW19_focs_conf}\bibhead{ChenJW19_focs_conf}
{\sc Lijie Chen, Ce~Jin, and R.~Ryan Williams}: Hardness magnification for all
  sparse {NP} languages.
\newblock In {\em Proc. 60th FOCS}, pp. 1240--1255. IEEE Comp. Soc., 2019.
\newblock [\epfmtdoi{10.1109/FOCS.2019.00077}]

\bibitem{ChengH20_coco_conf}\bibhead{ChengH20_coco_conf}
{\sc Kuan Cheng and William~M. Hoza}: Hitting sets give two-sided
  derandomization of small space.
\newblock {\em Theory of Computing}, 18(21):1--32, 2022.
\newblock \prelim \href{https://doi.org/10.4230/LIPIcs.CCC.2020.10}{CCC'20}.
\newblock [\epfmtdoi{10.4086/toc.2022.v018a021}]

\bibitem{CheraghchiHMY20}\bibhead{CheraghchiHMY20}
{\sc Mahdi Cheraghchi, Shuichi Hirahara, Dimitrios Myrisiotis, and Yuichi
  Yoshida}: One-tape {T}uring machine and branching program lower bounds for
  {MCSP}.
\newblock {\em Theory Computing Sys.}, (10113-9):32 pages, 2022.
\newblock \prelim \href{https://doi.org/10.4230/LIPIcs.STACS.2021.23}{MFCS'21}.
\newblock [\epfmtdoi{10.1007/s00224-022-10113-9}]

\bibitem{CheraghchiKLM19_icalp_conf}\bibhead{CheraghchiKLM19_icalp_conf}
{\sc Mahdi Cheraghchi, Valentine Kabanets, Zhenjian Lu, and Dimitrios
  Myrisiotis}: Circuit lower bounds for {MCSP} from local pseudorandom
  generators.
\newblock {\em ACM Trans. Comput. Theory}, 12(3):21:1--27, 2020.
\newblock \prelim
  \href{https://doi.org/10.4230/LIPIcs.ICALP.2019.39}{ICALP'19}.
\newblock [\epfmtdoi{10.1145/3404860}]

\bibitem{EvenSY84_iandc_journals}\bibhead{EvenSY84_iandc_journals}
{\sc Shimon Even, Alan~L. Selman, and Yacov Yacobi}: The complexity of promise
  problems with applications to public-key cryptography.
\newblock {\em Inform. Control}, 61(2):159--173, 1984.
\newblock [\epfmtdoi{10.1016/S0019-9958(84)80056-X}]

\bibitem{FeigenbaumF93_siamcomp_journals}\bibhead{FeigenbaumF93_siamcomp_journals}
{\sc Joan Feigenbaum and Lance Fortnow}: Random-self-reducibility of complete
  sets.
\newblock {\em SIAM J. Comput.}, 22(5):994--1005, 1993.
\newblock [\epfmtdoi{10.1137/0222061}]

\bibitem{FindGHK16_focs_conf}\bibhead{FindGHK16_focs_conf}
{\sc Magnus~Gausdal Find, Alexander Golovnev, Edward~A. Hirsch, and
  Alexander~S. Kulikov}: A better-than-$3n$ lower bound for the circuit
  complexity of an explicit function.
\newblock In {\em Proc. 57th FOCS}, pp. 89--98. IEEE Comp. Soc., 2016.
\newblock [\epfmtdoi{10.1109/FOCS.2016.19}]

\bibitem{ForbesK18_focs_conf}\bibhead{ForbesK18_focs_conf}
{\sc Michael~A. Forbes and Zander Kelley}: Pseudorandom generators for
  read-once branching programs, in any order.
\newblock In {\em Proc. 59th FOCS}, pp. 946--955. IEEE Comp. Soc., 2018.
\newblock [\epfmtdoi{10.1109/FOCS.2018.00093}]

\bibitem{Fortnow01_coco_conf}\bibhead{Fortnow01_coco_conf}
{\sc Lance Fortnow}: Comparing notions of full derandomization.
\newblock In {\em Proc. 16th IEEE Conf. on Comput. Complexity (CCC'01)}, pp.
  28--34. IEEE Comp. Soc., 2001.
\newblock [\epfmtdoi{10.1109/CCC.2001.933869}]

\bibitem{Goldreich06a_birthday_conf}\bibhead{Goldreich06a_birthday_conf}
{\sc Oded Goldreich}: On promise problems: {A} survey.
\newblock In {\sc Oded Goldreich, Arnold~L. Rosenberg, and Alan~L. Selman},
  editors, {\em Theoretical Computer Science, Essays in Memory of Shimon Even},
  pp. 254--290. Springer, 2006.
\newblock [\epfmtdoi{10.1007/11685654\_12}]

\bibitem{Goldreich11j_goldreich2011_sp_books}\bibhead{Goldreich11j_goldreich2011_sp_books}
{\sc Oded Goldreich}: A sample of samplers: {A} computational perspective on
  sampling.
\newblock In {\sc Oded Goldreich}, editor, {\em Studies in Complexity and
  Cryptography. Miscellanea on the Interplay between Randomness and
  Computation}, pp. 302--332. Springer, 2011.
\newblock [\epfmtdoi{10.1007/978-3-642-22670-0\_24}]

\bibitem{GoldreichGM86_jacm_journals}\bibhead{GoldreichGM86_jacm_journals}
{\sc Oded Goldreich, Shafi Goldwasser, and Silvio Micali}: How to construct
  random functions.
\newblock {\em J. ACM}, 33(4):792--807, 1986.
\newblock [\epfmtdoi{10.1145/6490.6503}]

\bibitem{GoldreichW14_stoc_conf}\bibhead{GoldreichW14_stoc_conf}
{\sc Oded Goldreich and Avi Wigderson}: On derandomizing algorithms that err
  extremely rarely.
\newblock In {\em Proc. 46th STOC}, pp. 109--118. ACM Press, 2014.
\newblock [\epfmtdoi{10.1145/2591796.2591808}]

\bibitem{GoldreichW13_eccc_journals}\bibhead{GoldreichW13_eccc_journals}
{\sc Oded Goldreich and Avi Wigderson}: On the size of depth-three boolean
  circuits for computing multilinear functions.
\newblock In {\sc Oded Goldreich}, editor, {\em Comput. Complexity and Property
  Testing}, LNCS 12050, pp. 41--86. Springer Cham, 2020.
\newblock (2013).
\newblock [\epfmt{eccc}{TR13-043}]

\bibitem{Hastad87_PhD_Thesis}\bibhead{Hastad87_PhD_Thesis}
{\sc Johan H{\aa}stad}: {\em Computational limitations of small-depth
  circuits}.
\newblock Ph.\,D.\ thesis, MIT, 1986.
\newblock \href{https://www.csc.kth.se/~johanh/thesis_with_figs.pdf}{Available
  on author's website}, published as a book by
  \href{https://mitpress.mit.edu/9780262081672/computational-limitations-for-small-depth-circuits/}{MIT
  Press, 1987}.

\bibitem{HastadILL99_siamcomp_journals}\bibhead{HastadILL99_siamcomp_journals}
{\sc Johan H{\aa}stad, Russell Impagliazzo, Leonid~A. Levin, and Michael Luby}:
  A pseudorandom generator from any one-way function.
\newblock {\em SIAM J. Comput.}, 28(4):1364--1396, 1999.
\newblock [\epfmtdoi{10.1137/S0097539793244708}]

\bibitem{HealyVV06_siamcomp_journals}\bibhead{HealyVV06_siamcomp_journals}
{\sc Alexander Healy, Salil~P. Vadhan, and Emanuele Viola}: Using
  nondeterminism to amplify hardness.
\newblock {\em SIAM J. Comput.}, 35(4):903--931, 2006.
\newblock [\epfmtdoi{10.1137/S0097539705447281}]

\bibitem{Hirahara18_focs_conf}\bibhead{Hirahara18_focs_conf}
{\sc Shuichi Hirahara}: Non-black-box worst-case to average-case reductions
  within {NP}.
\newblock In {\em Proc. 59th FOCS}, pp. 247--258. IEEE Comp. Soc., 2018.
\newblock [\epfmtdoi{10.1109/FOCS.2018.00032}]

\bibitem{Hirahara20_focs_conf}\bibhead{Hirahara20_focs_conf}
{\sc Shuichi Hirahara}: Characterizing average-case complexity of {PH} by
  worst-case meta-complexity.
\newblock In {\em Proc. 61st FOCS}, pp. 50--60. IEEE Comp. Soc., 2020.
\newblock [\epfmtdoi{10.1109/FOCS46700.2020.00014}]

\bibitem{Hirahara20_coco_conf}\bibhead{Hirahara20_coco_conf}
{\sc Shuichi Hirahara}: Non-disjoint promise problems from meta-computational
  view of pseudorandom generator constructions.
\newblock In {\em Proc. 35th Comput. Complexity Conf. (CCC'20)}, pp. 20:1--47.
  Schloss Dagstuhl--Leibniz-Zentrum fuer Informatik, 2020.
\newblock [\epfmtdoi{10.4230/LIPIcs.CCC.2020.20}]

\bibitem{Hirahara20_stoc_conf}\bibhead{Hirahara20_stoc_conf}
{\sc Shuichi Hirahara}: Unexpected hardness results for {K}olmogorov complexity
  under uniform reductions.
\newblock In {\em Proc. 52nd STOC}, pp. 1038--1051. ACM Press, 2020.
\newblock [\epfmtdoi{10.1145/3357713.3384251}]

\bibitem{HiraharaOS18_coco_conf}\bibhead{HiraharaOS18_coco_conf}
{\sc Shuichi Hirahara, Igor~Carboni Oliveira, and Rahul Santhanam}:
  {NP}-hardness of {M}inimum {C}ircuit {S}ize {P}roblem for {OR-AND-MOD}
  circuits.
\newblock In {\em Proc. 33rd Comput. Complexity Conf. (CCC'18)}, pp. 5:1--31.
  Schloss Dagstuhl--Leibniz-Zentrum fuer Informatik, 2018.
\newblock [\epfmtdoi{10.4230/LIPIcs.CCC.2018.5}]

\bibitem{HiraharaS17_coco_conf}\bibhead{HiraharaS17_coco_conf}
{\sc Shuichi Hirahara and Rahul Santhanam}: On the average-case complexity of
  {MCSP} and its variants.
\newblock In {\em Proc. 32nd Comput. Complexity Conf. (CCC'17)}, pp. 7:1--20.
  Schloss Dagstuhl--Leibniz-Zentrum fuer Informatik, 2017.
\newblock [\epfmtdoi{10.4230/LIPIcs.CCC.2017.7}]

\bibitem{HiraharaW16_coco_conf}\bibhead{HiraharaW16_coco_conf}
{\sc Shuichi Hirahara and Osamu Watanabe}: Limits of {M}inimum {C}ircuit {S}ize
  {P}roblem as oracle.
\newblock In {\em Proc. 31st Comput. Complexity Conf. (CCC'16)}, pp. 18:1--20.
  Schloss Dagstuhl--Leibniz-Zentrum fuer Informatik, 2016.
\newblock [\epfmtdoi{10.4230/LIPIcs.CCC.2016.18}]

\bibitem{HiraharaW20_approx_conf}\bibhead{HiraharaW20_approx_conf}
{\sc Shuichi Hirahara and Osamu Watanabe}: On nonadaptive security reductions
  of hitting set generators.
\newblock In {\em Proc. 24th Internat. Conf. on Randomization and Computation
  (RANDOM'20)}, pp. 15:1--14. Schloss Dagstuhl--Leibniz-Zentrum fuer
  Informatik, 2020.
\newblock [\epfmtdoi{10.4230/LIPIcs.APPROX/RANDOM.2020.15}]

\bibitem{HitchcockP15_fsttcs_conf}\bibhead{HitchcockP15_fsttcs_conf}
{\sc John~M. Hitchcock and Aduri Pavan}: On the {NP}-completeness of the
  {M}inimum {C}ircuit {S}ize {P}roblem.
\newblock In {\em Proc. 35th Found. Softw. Techn. Theoret. Comp. Sci. Conf.
  (FSTTCS'15)}, pp. 236--245. Schloss Dagstuhl--Leibniz-Zentrum fuer
  Informatik, 2015.
\newblock [\epfmtdoi{10.4230/LIPIcs.FSTTCS.2015.236}]

\bibitem{Impagliazzo95_focs_conf}\bibhead{Impagliazzo95_focs_conf}
{\sc Russell Impagliazzo}: Hard-core distributions for somewhat hard problems.
\newblock In {\em Proc. 36th FOCS}, pp. 538--545. IEEE Comp. Soc., 1995.
\newblock [\epfmtdoi{10.1109/SFCS.1995.492584}]

\bibitem{Impagliazzo95_coco_conf}\bibhead{Impagliazzo95_coco_conf}
{\sc Russell Impagliazzo}: A personal view of average-case complexity.
\newblock In {\em Proc. 10th IEEE Conf. Structure in Complexity Theory
  (SCT'95)}, pp. 134--147. IEEE Comp. Soc., 1995.
\newblock [\epfmtdoi{10.1109/SCT.1995.514853}]

\bibitem{Impagliazzo11_coco_conf}\bibhead{Impagliazzo11_coco_conf}
{\sc Russell Impagliazzo}: Relativized separations of worst-case and
  average-case complexities for {NP}.
\newblock In {\em Proc. 26th IEEE Conf. on Comput. Complexity (CCC'11)}, pp.
  104--114. IEEE Comp. Soc., 2011.
\newblock [\epfmtdoi{10.1109/CCC.2011.34}]

\bibitem{ImpagliazzoW97_stoc_conf}\bibhead{ImpagliazzoW97_stoc_conf}
{\sc Russell Impagliazzo and Avi Wigderson}: \emph{P = BPP} if \emph{E}
  requires exponential circuits: {D}erandomizing the {XOR} lemma.
\newblock In {\em Proc. 29th STOC}, pp. 220--229. ACM Press, 1997.
\newblock [\epfmtdoi{10.1145/258533.258590}]

\bibitem{KabanetsC00_stoc_conf}\bibhead{KabanetsC00_stoc_conf}
{\sc Valentine Kabanets and Jin{-}yi Cai}: Circuit minimization problem.
\newblock In {\em Proc. 32nd STOC}, pp. 73--79. ACM Press, 2000.
\newblock [\epfmtdoi{10.1145/335305.335314}]

\bibitem{KarpL82_advice}\bibhead{KarpL82_advice}
{\sc Richard~M. Karp and Richard~J. Lipton}: Turing machines that take advice.
\newblock {\em Enseign. Math.}, 28:191--209, 1982.
\newblock [\epfmtdoi{10.5169/seals-52237}]

\bibitem{KlivansS03_ml_journals}\bibhead{KlivansS03_ml_journals}
{\sc Adam~R. Klivans and Rocco~A. Servedio}: Boosting and hard-core set
  construction.
\newblock {\em Machine Learning}, 51(3):217--238, 2003.
\newblock [\epfmtdoi{10.1023/A:1022949332276}]

\bibitem{KlivansM02_siamcomp_journals}\bibhead{KlivansM02_siamcomp_journals}
{\sc Adam~R. Klivans and Dieter van Melkebeek}: Graph nonisomorphism has
  subexponential size proofs unless the {P}olynomial-time {H}ierarchy
  collapses.
\newblock {\em SIAM J. Comput.}, 31(5):1501--1526, 2002.
\newblock [\epfmtdoi{10.1137/S0097539700389652}]

\bibitem{Ko91_siamcomp_journals}\bibhead{Ko91_siamcomp_journals}
{\sc Ker{-}I Ko}: On the complexity of learning minimum time-bounded {T}uring
  machines.
\newblock {\em SIAM J. Comput.}, 20(5):962--986, 1991.
\newblock [\epfmtdoi{10.1137/0220059}]

\bibitem{Levin84_iandc_journals}\bibhead{Levin84_iandc_journals}
{\sc Leonid~A. Levin}: Randomness conservation inequalities; {I}nformation and
  independence in mathematical theories.
\newblock {\em Inform. Control}, 61(1):15--37, 1984.
\newblock [\epfmtdoi{10.1016/S0019-9958(84)80060-1}]

\bibitem{Levin86_siamcomp_journals}\bibhead{Levin86_siamcomp_journals}
{\sc Leonid~A. Levin}: Average case complete problems.
\newblock {\em SIAM J. Comput.}, 15(1):285--286, 1986.
\newblock [\epfmtdoi{10.1137/0215020}]

\bibitem{Masek79}\bibhead{Masek79}
{\sc William~J. Masek}: Some {NP}-complete set covering problems.
\newblock Master's thesis, MIT CS Lab, 1978.
\newblock Cited in the
  \href{https://www.sciencedirect.com/science/article/pii/0196677487900216}{NP-completeness
  column} by David Johnson.

\bibitem{McKayMW19_stoc_conf}\bibhead{McKayMW19_stoc_conf}
{\sc Dylan~M. McKay, Cody~D. Murray, and R.~Ryan Williams}: Weak lower bounds
  on resource-bounded compression imply strong separations of complexity
  classes.
\newblock In {\em Proc. 51st STOC}, pp. 1215--1225. ACM Press, 2019.
\newblock [\epfmtdoi{10.1145/3313276.3316396}]

\bibitem{MurrayW17_toc_journals}\bibhead{MurrayW17_toc_journals}
{\sc Cody~D. Murray and R.~Ryan Williams}: On the (non) {NP}-hardness of
  computing circuit complexity.
\newblock {\em Theory of Computing}, 13(4):1--22, 2017.
\newblock [\epfmtdoi{10.4086/toc.2017.v013a004}]

\bibitem{Nisan92_combinatorica_journals}\bibhead{Nisan92_combinatorica_journals}
{\sc Noam Nisan}: Pseudorandom generators for space-bounded computation.
\newblock {\em Combinatorica}, 12(4):449--461, 1992.
\newblock [\epfmtdoi{10.1007/BF01305237}]

\bibitem{NisanW94_jcss_journals}\bibhead{NisanW94_jcss_journals}
{\sc Noam Nisan and Avi Wigderson}: Hardness vs randomness.
\newblock {\em J. Comput. System Sci.}, 49(2):149--167, 1994.
\newblock [\epfmtdoi{10.1016/S0022-0000(05)80043-1}]

\bibitem{OliveiraPS19_coco_conf}\bibhead{OliveiraPS19_coco_conf}
{\sc Igor~Carboni Oliveira, J{\'{a}}n Pich, and Rahul Santhanam}: Hardness
  magnification near state-of-the-art lower bounds.
\newblock {\em Theory of Computing}, 17(11):1--38, 2021.
\newblock \prelim \href{https://doi.org/10.4230/LIPIcs.CCC.2019.27}{CCC'19}.
\newblock [\epfmtdoi{10.4086/toc.2021.v017a011}]

\bibitem{OliveiraS18_focs_conf}\bibhead{OliveiraS18_focs_conf}
{\sc Igor~Carboni Oliveira and Rahul Santhanam}: Hardness magnification for
  natural problems.
\newblock In {\em Proc. 59th FOCS}, pp. 65--76. IEEE Comp. Soc., 2018.
\newblock [\epfmtdoi{10.1109/FOCS.2018.00016}]

\bibitem{Ostrovsky91_coco_conf}\bibhead{Ostrovsky91_coco_conf}
{\sc Rafail Ostrovsky}: One-way functions, hard on average problems, and
  statistical zero-knowledge proofs.
\newblock In {\em Proc. 6th IEEE Conf. Structure in Complexity Theory
  (SCT'91)}, pp. 133--138. IEEE Comp. Soc., 1991.
\newblock [\epfmtdoi{10.1109/SCT.1991.160253}]

\bibitem{PavanSV06_fsttcs_conf}\bibhead{PavanSV06_fsttcs_conf}
{\sc Aduri Pavan, Rahul Santhanam, and N.~V. Vinodchandran}: Some results on
  average-case hardness within the {P}olynomial {H}ierarchy.
\newblock In {\em Proc. 26th Found. Softw. Techn. Theoret. Comp. Sci. Conf.
  (FSTTCS'06)}, pp. 188--199. Springer, 2006.
\newblock [\epfmtdoi{10.1007/11944836\_19}]

\bibitem{RazRV02_jcss_journals}\bibhead{RazRV02_jcss_journals}
{\sc Ran Raz, Omer Reingold, and Salil~P. Vadhan}: Extracting all the
  randomness and reducing the error in {Trevisan's} extractors.
\newblock {\em J. Comput. System Sci.}, 65(1):97--128, 2002.
\newblock [\epfmtdoi{10.1006/jcss.2002.1824}]

\bibitem{Razborov87_lowerbound}\bibhead{Razborov87_lowerbound}
{\sc Alexander~A. Razborov}: Lower bounds on the size of bounded depth circuits
  over a complete basis with logical addition.
\newblock {\em Math. Notes. Acad. Sci. USSR (English translation)},
  41(4):333--338, 1987.
\newblock [\epfmtdoi{10.1007/BF01137685}]

\bibitem{RazborovR97_jcss_journals}\bibhead{RazborovR97_jcss_journals}
{\sc Alexander~A. Razborov and Steven Rudich}: Natural {P}roofs.
\newblock {\em J. Comput. System Sci.}, 55(1):24--35, 1997.
\newblock [\epfmtdoi{10.1006/jcss.1997.1494}]

\bibitem{ReingoldTV04_tcc_conf}\bibhead{ReingoldTV04_tcc_conf}
{\sc Omer Reingold, Luca Trevisan, and Salil~P. Vadhan}: Notions of
  reducibility between cryptographic primitives.
\newblock In {\em Proc. Theory of Cryptography Conf. (TCC'04)}, pp. 1--20.
  Springer, 2004.
\newblock [\epfmtdoi{10.1007/978-3-540-24638-1\_1}]

\bibitem{Santhanam20_innovations_conf}\bibhead{Santhanam20_innovations_conf}
{\sc Rahul Santhanam}: Pseudorandomness and the {M}inimum {C}ircuit {S}ize
  {P}roblem.
\newblock In {\em Proc. 11th Innovations in Theoret. Comp. Sci. Conf.
  (ITCS'20)}, pp. 68:1--26. Schloss Dagstuhl--Leibniz-Zentrum fuer Informatik,
  2020.
\newblock [\epfmtdoi{10.4230/LIPIcs.ITCS.2020.68}]

\bibitem{ShaltielV10_siamcomp_journals}\bibhead{ShaltielV10_siamcomp_journals}
{\sc Ronen Shaltiel and Emanuele Viola}: Hardness amplification proofs require
  majority.
\newblock {\em SIAM J. Comput.}, 39(7):3122--3154, 2010.
\newblock [\epfmtdoi{10.1137/080735096}]

\bibitem{Sudan97_jc_journals}\bibhead{Sudan97_jc_journals}
{\sc Madhu Sudan}: Decoding of {R}eed {S}solomon codes beyond the
  error-correction bound.
\newblock {\em J. Complexity}, 13(1):180--193, 1997.
\newblock [\epfmtdoi{10.1006/jcom.1997.0439}]

\bibitem{SudanTV01_jcss_journals}\bibhead{SudanTV01_jcss_journals}
{\sc Madhu Sudan, Luca Trevisan, and Salil~P. Vadhan}: Pseudorandom generators
  without the {XOR} lemma.
\newblock {\em J. Comput. System Sci.}, 62(2):236--266, 2001.
\newblock [\epfmtdoi{10.1006/jcss.2000.1730}]

\bibitem{Trakhtenbrot84_annals_journals}\bibhead{Trakhtenbrot84_annals_journals}
{\sc Boris~A. Trakhtenbrot}: A survey of {R}ussian approaches to {P}erebor
  (brute-force searches) algorithms.
\newblock {\em IEEE Annals of the History of Computing}, 6(4):384--400, 1984.
\newblock [\epfmtdoi{10.1109/MAHC.1984.10036}]

\bibitem{Trevisa03_focs_conf}\bibhead{Trevisa03_focs_conf}
{\sc Luca Trevisan}: List-decoding using the {XOR} lemma.
\newblock In {\em Proc. 44th FOCS}, pp. 126--135. IEEE Comp. Soc., 2003.
\newblock [\epfmtdoi{10.1109/SFCS.2003.1238187}]

\bibitem{TrevisanV07_cc_journals}\bibhead{TrevisanV07_cc_journals}
{\sc Luca Trevisan and Salil~P. Vadhan}: Pseudorandomness and average-case
  complexity via uniform reductions.
\newblock {\em Comput. Complexity}, 16(4):331--364, 2007.
\newblock [\epfmtdoi{10.1007/s00037-007-0233-x}]

\bibitem{TrevisanX13_coco_conf}\bibhead{TrevisanX13_coco_conf}
{\sc Luca Trevisan and Tongke Xue}: A derandomized {S}witching {L}emma and an
  improved derandomization of {AC0}.
\newblock In {\em Proc. 28th IEEE Conf. on Comput. Complexity (CCC'13)}, pp.
  242--247. IEEE Comp. Soc., 2013.
\newblock [\epfmtdoi{10.1109/CCC.2013.32}]

\bibitem{Vadhan06_siamcomp_journals}\bibhead{Vadhan06_siamcomp_journals}
{\sc Salil~P. Vadhan}: An unconditional study of {C}omputational {Z}ero
  {K}nowledge.
\newblock {\em SIAM J. Comput.}, 36(4):1160--1214, 2006.
\newblock [\epfmtdoi{10.1137/S0097539705447207}]

\bibitem{Vadhan12_fttcs_journals}\bibhead{Vadhan12_fttcs_journals}
{\sc Salil~P. Vadhan}: Pseudorandomness.
\newblock {\em Found. Trends Theor. Comp. Sci.}, 7(1--3):1--336, 2012.
\newblock [\epfmtdoi{10.1561/0400000010}]

\bibitem{ValiantV86_tcs_journals}\bibhead{ValiantV86_tcs_journals}
{\sc Leslie~G. Valiant and Vijay~V. Vazirani}: {NP} is as easy as detecting
  unique solutions.
\newblock {\em Theoret. Comput. Sci.}, 47:85--93, 1986.
\newblock [\epfmtdoi{10.1016/0304-3975(86)90135-0}]

\bibitem{Viola05_cc_journals}\bibhead{Viola05_cc_journals}
{\sc Emanuele Viola}: The complexity of constructing pseudorandom generators
  from hard functions.
\newblock {\em Comput. Complexity}, 13(3--4):147--188, 2005.
\newblock [\epfmtdoi{10.1007/s00037-004-0187-1}]

\bibitem{Yao82a_focs_conf}\bibhead{Yao82a_focs_conf}
{\sc Andrew~Chi{-}Chih Yao}: Theory and applications of trapdoor functions
  (extended abstract).
\newblock In {\em Proc. 23rd FOCS}, pp. 80--91. IEEE Comp. Soc., 1982.
\newblock [\epfmtdoi{10.1109/SFCS.1982.45}]

\bibitem{Yekhanin12_fttcs_journals}\bibhead{Yekhanin12_fttcs_journals}
{\sc Sergey Yekhanin}: Locally decodable codes.
\newblock {\em Found. Trends Theor. Comp. Sci.}, 6(3):139--255, 2012.
\newblock [\epfmtdoi{10.1561/0400000030}]

\end{thebibliography}

\end{document}
